% Work, Energy and Power — Pure Mathematics with Mechanics Upper Sixth — Cours signé OnBuch+
% Official MINESEC syllabus. Compiler avec : tectonic PMM15-work-energy-and-power.tex
\newcommand{\DOCMATIERE}{Pure Mathematics with Mechanics}
\newcommand{\DOCNIVEAU}{Upper Sixth}
\newcommand{\DOCMODULE}{Upper Sixth — Pure mathematics with mechanics}
\newcommand{\DOCLECON}{15}
\newcommand{\DOCTITRE}{Work, Energy and Power}
% OnBuch+ — préambule « cours signés » ENRICHI (compilé avec tectonic / XeLaTeX)
% Système visuel « Pop » d'OnBuch+ : fond crème (jamais blanc), bordures encre
% épaisses, ombres « dures » décalées, titres Archivo Black en capitales.
% Version MATHÉMATIQUES : graphes (pgfplots/tikz), boîtes pédagogiques
% (définition, propriété, méthode, attention, le savais-tu, exercice résolu,
% exercice, TP, bilan), page de garde et signature OnBuch+ ancrée en bas.
%
% Macros attendues dans main.tex AVANT \input{preamble} :
%   \DOCMATIERE  \DOCNIVEAU  \DOCMODULE  \DOCLECON (numéro)  \DOCTITRE
\documentclass[11pt]{article}

\usepackage{fontspec}
\usepackage[english]{babel}
\usepackage[a4paper,margin=2cm,top=3.4cm,bottom=2.8cm,headheight=1.4cm,footskip=1.3cm]{geometry}
\usepackage{amsmath,amssymb,amsfonts}
\usepackage{mathtools} % \xmapsto, \coloneqq... souvent utilisés par les modèles
\usepackage{cancel} % simplifications barrées (bilans de forces)
% \abs{z} (module, valeur absolue) et \norm{u} (norme), utilisés par les modèles
\providecommand{\abs}{}\let\abs\relax\DeclarePairedDelimiter{\abs}{\lvert}{\rvert}
\providecommand{\norm}{}\let\norm\relax\DeclarePairedDelimiter{\norm}{\lVert}{\rVert}
\usepackage[table]{xcolor} % [table] : fournit aussi \rowcolors (alternance de lignes)
\usepackage{pagecolor}
\usepackage{tikz}
\usetikzlibrary{arrows.meta,positioning,calc,shapes.geometric,shapes.misc,shapes.arrows,decorations.pathmorphing,decorations.pathreplacing,decorations.markings,patterns,shadows,mindmap,trees,fit,backgrounds,matrix,intersections,angles,quotes}
\usepackage{pgfplots}
\usepackage[version=4]{mhchem} % équations nucléaires \ce{^{235}_{92}U}
\usepackage[european,siunitx]{circuitikz} % schémas électriques
\pgfplotsset{compat=1.18}
\usepgfplotslibrary{fillbetween,groupplots} % aires entre courbes (calcul intégral)
\usepackage{enumitem}
\usepackage{fancyhdr}
% [most] charge toutes les bibliothèques tcolorbox (listings, minted, raster,
% documentation...) et gonfle inutilement la mémoire TeX sur un document long
% (~300 boîtes) : on ne charge que ce qui est réellement utilisé.
\usepackage[skins,breakable]{tcolorbox}
\usepackage{graphicx}
\usepackage{colortbl}
\usepackage{booktabs}
\usepackage{tabularx}
\usepackage{diagbox} % case d'en-tête diagonale des tableaux à double entrée
% Colonnes extensibles centrée / gauche / droite (C, L, R), souvent utilisées par les modèles
\newcolumntype{C}{>{\centering\arraybackslash}X}
\newcolumntype{L}{>{\raggedright\arraybackslash}X}
\newcolumntype{R}{>{\raggedleft\arraybackslash}X}
\usepackage{array}
\usepackage{multicol}
\usepackage{siunitx}
% parse-numbers=false : accepte aussi \SI{2,5 \times 10^{-3}}{...}
\sisetup{locale=UK,output-decimal-marker={.},group-minimum-digits=4,parse-numbers=false}
\DeclareSIUnit\ohm{\ensuremath{\symup{\Omega}}} % U+2126 absent de la police
\DeclareSIUnit\division{div} % réglages d'oscilloscope (ms/div, V/div)
\DeclareSIUnit\tr{tr} % tours (vitesse de rotation, ex. \SI{1500}{\tr\per\minute})
\usepackage{qrcode}
% \degree : commande fréquemment utilisée par les modèles pour un angle ou une
% température en dehors de \SI{}{} (non fournie par les paquets chargés).
\providecommand{\degree}{^{\circ}}
% \qed (fin de démonstration), utilisé par les modèles sans amsthm
\providecommand{\qed}{\hfill$\square$}
% \barre : double barre des valeurs interdites dans un tableau de variations (tkz-tab)
\providecommand{\barre}{\ensuremath{\|}}
% environnement proof (amsthm non chargé) utilisé par les modèles
\ifdefined\proof\else\newenvironment{proof}[1][Proof]{\par\noindent\textit{#1.}\ }{\qed\par}\fi
\usepackage{lastpage}
\usepackage{hyperref}
\hypersetup{hidelinks}

% ============================================================
% Palette « Pop » OnBuch+ — mêmes hex que la classe P (plus_ui.dart)
% ============================================================
\definecolor{poporange}{HTML}{F59321}
\definecolor{popdark}{HTML}{E07A0C}
\definecolor{popcream}{HTML}{FFF6E8}
\definecolor{popink}{HTML}{1C1714}
\definecolor{poppurple}{HTML}{7A5AE0}
\definecolor{popgreen}{HTML}{2E9E6A}
\definecolor{poppink}{HTML}{E0457B}
\definecolor{popblue}{HTML}{2F7DE1}
\definecolor{popgold}{HTML}{C98A12}
\definecolor{popmuted}{HTML}{8A7F76}
\definecolor{popline}{HTML}{EADFD2}
\definecolor{popgray}{HTML}{8A7F76}
\colorlet{popgrey}{popgray}
% Alias tolérants (le modèle nomme parfois les couleurs différemment)
\colorlet{popwhite}{white}
\colorlet{popblack}{popink}
\colorlet{popred}{poppink}
\colorlet{popyellow}{popgold}
% Teintes claires pour fonds de tableaux / zones
\colorlet{popblueL}{popblue!10!white}
\colorlet{poporangeL}{poporange!14!white}
\colorlet{popgreenL}{popgreen!12!white}
\colorlet{poppurpleL}{poppurple!12!white}
\colorlet{poppinkL}{poppink!12!white}
\colorlet{popgoldL}{popgold!14!white}

\pagecolor{popcream}

% ============================================================
% Typographie
% ============================================================
\defaultfontfeatures{Ligatures=TeX}
\setmainfont{PlusJakartaSans-Medium.ttf}[Path=../../../fonts/,BoldFont=PlusJakartaSans-Bold.ttf]
% Maths en sans-serif (Fira Math) avec chiffres et lettres droites de Plus
% Jakarta, pour que les formules se fondent dans le texte.
\usepackage[math-style=ISO,bold-style=ISO]{unicode-math}
\setmathfont{FiraMath-Regular.otf}[Path=../../../fonts/]
\setmathfont{PlusJakartaSans-Medium.ttf}[Path=../../../fonts/,range={up/num,up/{Latin,latin},\mathup}]
% (icomma retiré : décimales avec point en anglais)
% Caractères mathématiques Unicode absents des polices de texte (Plus Jakarta,
% Archivo Black) : rendus par leur équivalent mathématique, en texte comme en
% formule — sinon ils s'affichent en carré vide (ex. « Arithmétique dans ℤ »).
\usepackage{newunicodechar}
\newunicodechar{ℝ}{\ensuremath{\mathbb{R}}}
\newunicodechar{ℤ}{\ensuremath{\mathbb{Z}}}
\newunicodechar{ℂ}{\ensuremath{\mathbb{C}}}
\newunicodechar{ℕ}{\ensuremath{\mathbb{N}}}
\newunicodechar{ℚ}{\ensuremath{\mathbb{Q}}}
\newunicodechar{𝔻}{\ensuremath{\mathbb{D}}}
\newunicodechar{α}{\ensuremath{\alpha}}
\newunicodechar{β}{\ensuremath{\beta}}
\newunicodechar{γ}{\ensuremath{\gamma}}
\newunicodechar{δ}{\ensuremath{\delta}}
\newunicodechar{ε}{\ensuremath{\varepsilon}}
\newunicodechar{θ}{\ensuremath{\theta}}
\newunicodechar{λ}{\ensuremath{\lambda}}
\newunicodechar{μ}{\ensuremath{\mu}}
\newunicodechar{σ}{\ensuremath{\sigma}}
\newunicodechar{τ}{\ensuremath{\tau}}
\newunicodechar{φ}{\ensuremath{\varphi}}
\newunicodechar{ψ}{\ensuremath{\psi}}
\newunicodechar{ω}{\ensuremath{\omega}}
\newunicodechar{Σ}{\ensuremath{\Sigma}}
% majuscules produites par \MakeUppercase dans les titres (α-aminés -> Α)
\newunicodechar{Α}{\ensuremath{\alpha}}
\newunicodechar{Β}{\ensuremath{\beta}}
\newunicodechar{Γ}{\ensuremath{\Gamma}}
\newunicodechar{Δ}{\ensuremath{\Delta}}
\newunicodechar{Ε}{\ensuremath{\varepsilon}}
\newunicodechar{Θ}{\ensuremath{\Theta}}
\newunicodechar{Λ}{\ensuremath{\Lambda}}
\newunicodechar{Μ}{\ensuremath{\mu}}
\newunicodechar{Τ}{\ensuremath{\tau}}
\newunicodechar{Φ}{\ensuremath{\Phi}}
\newunicodechar{Ψ}{\ensuremath{\Psi}}
\newunicodechar{Ω}{\ensuremath{\Omega}}
\newunicodechar{↦}{\ensuremath{\mapsto}}
\newunicodechar{∈}{\ensuremath{\in}}
\newunicodechar{∉}{\ensuremath{\notin}}
\newunicodechar{∧}{\ensuremath{\wedge}}
\newunicodechar{⇔}{\ensuremath{\Leftrightarrow}}
\newunicodechar{⟺}{\ensuremath{\Longleftrightarrow}}
\newunicodechar{⇒}{\ensuremath{\Rightarrow}}
\newunicodechar{⟹}{\ensuremath{\Longrightarrow}}
\newunicodechar{∘}{\ensuremath{\circ}}
\newunicodechar{∪}{\ensuremath{\cup}}
\newunicodechar{∩}{\ensuremath{\cap}}
\newunicodechar{≡}{\ensuremath{\equiv}}
\newunicodechar{⊂}{\ensuremath{\subset}}
\newunicodechar{⊥}{\ensuremath{\perp}}
\newunicodechar{∃}{\ensuremath{\exists}}
\newunicodechar{∀}{\ensuremath{\forall}}
\newunicodechar{∛}{\ensuremath{\sqrt[3]{\,}}}
\newunicodechar{∜}{\ensuremath{\sqrt[4]{\,}}}
\newunicodechar{⃗}{\ensuremath{{}^{\to}}}
\newunicodechar{ⁿ}{\ifmmode^{n}\else\textsuperscript{\ensuremath{n}}\fi}
\newunicodechar{ˣ}{\ifmmode^{x}\else\textsuperscript{\ensuremath{x}}\fi}
\newunicodechar{ᵇ}{\ifmmode^{b}\else\textsuperscript{\ensuremath{b}}\fi}
\newunicodechar{ʳ}{\ifmmode^{r}\else\textsuperscript{\ensuremath{r}}\fi}
\newunicodechar{ᵘ}{\ifmmode^{u}\else\textsuperscript{\ensuremath{u}}\fi}
\newunicodechar{ʸ}{\ifmmode^{y}\else\textsuperscript{\ensuremath{y}}\fi}
\newunicodechar{ᵃ}{\ifmmode^{a}\else\textsuperscript{\ensuremath{a}}\fi}
\newunicodechar{ᵉ}{\ifmmode^{e}\else\textsuperscript{\ensuremath{e}}\fi}
\newunicodechar{ᵏ}{\ifmmode^{k}\else\textsuperscript{\ensuremath{k}}\fi}
\newunicodechar{ᵖ}{\ifmmode^{p}\else\textsuperscript{\ensuremath{p}}\fi}
\newunicodechar{ᴺ}{\ifmmode^{N}\else\textsuperscript{\ensuremath{N}}\fi}
\newunicodechar{ᶜ}{\ifmmode^{c}\else\textsuperscript{\ensuremath{c}}\fi}
\newunicodechar{ʰ}{\ifmmode^{h}\else\textsuperscript{\ensuremath{h}}\fi}
\newunicodechar{ᵠ}{\ifmmode^{q}\else\textsuperscript{\ensuremath{q}}\fi}
\newunicodechar{ₙ}{\ifmmode_{n}\else\textsubscript{\ensuremath{n}}\fi}
\newunicodechar{ₐ}{\ifmmode_{a}\else\textsubscript{\ensuremath{a}}\fi}
\newunicodechar{ᵢ}{\ifmmode_{i}\else\textsubscript{\ensuremath{i}}\fi}
\newunicodechar{ᵣ}{\ifmmode_{r}\else\textsubscript{\ensuremath{r}}\fi}
\newunicodechar{ₖ}{\ifmmode_{k}\else\textsubscript{\ensuremath{k}}\fi}
\newunicodechar{ᵦ}{\ifmmode_{\beta}\else\textsubscript{\ensuremath{\beta}}\fi}
\newunicodechar{ₓ}{\ifmmode_{x}\else\textsubscript{\ensuremath{x}}\fi}
\newfontfamily\popdisplay{ArchivoBlack-Regular.ttf}[Path=../../../fonts/]
\newfontfamily\poplogo{SpaceGrotesk-Bold.ttf}[Path=../../../fonts/]
\newfontfamily\popsemi{PlusJakartaSans-SemiBold.ttf}[Path=../../../fonts/]
\newfontfamily\popxbold{PlusJakartaSans-ExtraBold.ttf}[Path=../../../fonts/]
\newfontfamily\popxboldls{PlusJakartaSans-ExtraBold.ttf}[Path=../../../fonts/,LetterSpace=3.0]

\setlist[enumerate]{leftmargin=1.7em,itemsep=5pt,topsep=4pt}
\setlist[itemize]{leftmargin=1.7em,itemsep=4pt,topsep=4pt,label={\color{poporange}\textbullet}}
\setlength{\parindent}{0pt}
\setlength{\parskip}{5pt}
\renewcommand{\arraystretch}{1.25}

% pgfplots : style Pop par défaut
\pgfplotsset{
  every axis/.append style={axis line style={popink,line width=1pt},tick style={popink},
    grid style={popline,line width=0.5pt},label style={font=\small},tick label style={font=\footnotesize}},
  popcurve/.style={poporange,line width=1.8pt,no markers,smooth},
}
% Même style disponible dans un \draw TikZ simple (hors pgfplots)
\tikzset{popcurve/.style={poporange,line width=1.8pt}}
\tikzset{popgrid/.style={popline,very thin}}
\colorlet{popcurve}{poporange} % le nom du style est parfois utilisé comme couleur

% ============================================================
% Pill / squiggle
% ============================================================
\newcommand{\poppill}[3]{%
  \tikz[baseline=-0.55ex]\node[draw=popink,line width=1.2pt,rounded corners=2.6pt,%
    fill=#2,inner xsep=6.5pt,inner ysep=3pt]{%
    \popxboldls\fontsize{7}{7}\selectfont\color{#3}\MakeUppercase{#1}};%
}
\newcommand{\popsquiggle}[1]{%
  \tikz[baseline=-0.4ex]{
    \draw[#1,line width=2.6pt,line cap=round]
      plot [smooth] coordinates {(0,0.09) (0.34,0.01) (0.68,0.21) (1.02,0.01) (1.36,0.10)};
  }%
}
% Mot-clé surligné (marqueur orange sous le texte)
\newcommand{\cle}[1]{{\popxbold\color{popdark}#1}}

% ============================================================
% En-tête / pied
% ============================================================
\pagestyle{fancy}
\fancyhf{}
\renewcommand{\headrulewidth}{0pt}
\renewcommand{\footrulewidth}{0pt}
\lhead{\raisebox{-0.15ex}{\poplogo\fontsize{13}{13}\selectfont\color{popink}OnBuch\color{poporange}+}}
\rhead{\poppill{\DOCMATIERE\ \textbullet\ \DOCNIVEAU\ \textbullet\ Chapter \DOCLECON}{white}{popink}}
\lfoot{\popxbold\fontsize{7.6}{7.6}\selectfont\color{popmuted}\MakeUppercase{Signed course OnBuch+}\ \textbullet\ {\popsemi\DOCTITRE}}
\rfoot{\tikz[baseline=-0.6ex]\node[rounded corners=8.5pt,fill=popink,minimum height=17pt,minimum width=17pt,inner xsep=5pt,inner sep=0]{\popxbold\fontsize{8}{8}\selectfont\color{popcream}\thepage\,/\,\pageref*{LastPage}};}

% ============================================================
% Titres
% ============================================================
\newcommand{\chaptitre}[2]{%
  \begin{tcolorbox}[enhanced,colback=poporange,colframe=popink,boxrule=2.6pt,arc=11pt,
      drop shadow=popink,left=15pt,right=15pt,top=12pt,bottom=11pt,before skip=3pt,after skip=18pt]
    \poppill{#2}{popink}{popcream}\par\vspace{9pt}
    {\raggedright\hyphenpenalty=10000\exhyphenpenalty=10000\popdisplay\fontsize{22}{24}\selectfont\color{popcream}\MakeUppercase{#1}\par}
  \end{tcolorbox}
}
% Section numérotée : pastille orange avec le numéro + titre Archivo
\newcounter{popsec}
\newcommand{\coursec}[1]{%
  \stepcounter{popsec}\setcounter{popsub}{0}%
  \par\vspace{16pt}\noindent
  \begin{minipage}{\linewidth}
  \tikz[baseline=-0.7ex]\node[draw=popink,line width=1.6pt,fill=poporange,rounded corners=4pt,minimum size=22pt,inner sep=0]{\popdisplay\fontsize{12}{12}\selectfont\color{popcream}\thepopsec};%
  \hspace{8pt}{\popdisplay\fontsize{15}{17}\selectfont\color{popink}\MakeUppercase{#1}}\par
  \vspace{4pt}\noindent{\color{poporange}\rule{36pt}{3.6pt}}
  \end{minipage}\par\nopagebreak\vspace{8pt}
}
\newcounter{popsub}
\newcommand{\courssub}[1]{%
  \stepcounter{popsub}%
  \par\vspace{9pt}\noindent{\popxbold\fontsize{11.5}{13}\selectfont\color{popink}{\color{poporange}\thepopsec.\thepopsub}\hspace{6pt}#1}\par\nopagebreak\vspace{4pt}
}

% ============================================================
% Boîtes pédagogiques — même recette : fond blanc, bordure épaisse colorée,
% ombre dure encre, badge plein. Titre optionnel [#1].
% ============================================================
\tcbset{popbase/.style={breakable,enhanced,colback=white,boxrule=2.3pt,arc=9pt,
  shadow={3pt}{-3pt}{0pt}{popink},left=14pt,right=14pt,top=11pt,bottom=11pt,before skip=11pt,after skip=15pt,
  fontupper=\popsemi\fontsize{10.6}{14.5}\selectfont\color{popink},
  fontlower=\popsemi\fontsize{10.6}{14.5}\selectfont\color{popink}}}
% Coche dessinée (le glyphe ✓ n'existe pas dans les polices du document)
\newcommand{\popcheck}[1]{\tikz[baseline=-0.5ex]\draw[#1,line width=1.6pt,line cap=round,line join=round](-0.13,0.0)--(-0.04,-0.09)--(0.13,0.1);}
\providecommand{\coche}{\popcheck{popgreen}}
\providecommand{\resultat}[1]{\par\smallskip\noindent\fcolorbox{popgreen}{popgreenL}{\parbox{\dimexpr\linewidth-2\fboxsep-2\fboxrule\relax}{\textbf{Result.} #1}}\par\smallskip}
\newcommand{\popboxhead}[3]{\poppill{#1}{#2}{white}\ifx\relax#3\relax\else\hspace{6pt}{\popxbold\fontsize{10.6}{12}\selectfont\color{popink}#3}\fi\par\vspace{7pt}}

\newtcolorbox{aretenir}[1][]{popbase,colframe=popgreen,before upper={\popboxhead{Key points}{popgreen}{#1}}}
\newtcolorbox{exemplebox}[1][]{popbase,colframe=poporange,before upper={\popboxhead{Exemple}{poporange}{#1}}}
\newtcolorbox{definition}[1][]{popbase,colframe=popblue,before upper={\popboxhead{Definition}{popblue}{#1}}}
\newtcolorbox{propriete}[1][]{popbase,colframe=poppurple,before upper={\popboxhead{Property}{poppurple}{#1}}}
\newtcolorbox{methode}[1][]{popbase,colframe=popgold,before upper={\popboxhead{Method}{popgold}{#1}}}
\newtcolorbox{attention}[1][]{popbase,colframe=poppink,before upper={\popboxhead{Warning}{poppink}{#1}}}
\newtcolorbox{experience}[1][]{popbase,colframe=popblue,colback=popblueL,before upper={\popboxhead{Experiment}{popblue}{#1}}}
% « Le savais-tu ? » : carte violette pleine (contexte, culture, Cameroun)
\newtcolorbox{savaistu}[1][]{popbase,colback=poppurple,colframe=popink,
  fontupper=\popsemi\fontsize{10.4}{14.2}\selectfont\color{white},
  before upper={\poppill{Did you know?}{white}{poppurple}\ifx\relax#1\relax\else\hspace{6pt}{\popxbold\color{white}#1}\fi\par\vspace{7pt}}}
% Exercice résolu : énoncé en haut, \tcblower puis la solution sur fond crème
\newcounter{popexo}\newcounter{popexr}
\newtcolorbox{exoresolu}[1][]{popbase,colframe=poporange,colbacklower=popcream,
  segmentation style={popink,line width=1.2pt,dashed},
  before upper={\stepcounter{popexr}\popboxhead{Worked example \thepopexr}{poporange}{#1}},
  before lower={\poppill{Solution}{popink}{popcream}\par\vspace{6pt}}}
% Exercice (non résolu en place) — la correction est dans la partie Corrigés
\newtcolorbox{exercice}[1][]{popbase,colframe=popink,
  before upper={\stepcounter{popexo}\popboxhead{Exercise \thepopexo}{popink}{#1}}}
% Corrigé
\newtcolorbox{corrige}[1][]{popbase,colframe=popgreen,colback=popgreenL,
  before upper={\popboxhead{Solution}{popgreen}{#1}}}
% Objectifs / prérequis / bilan
\newtcolorbox{objectifs}{popbase,colframe=popink,colback=poporangeL,
  before upper={\popboxhead{Chapter objectives}{popink}{}}}
\newtcolorbox{prerequis}{popbase,colframe=popink,colback=popblueL,
  before upper={\popboxhead{Prerequisites}{popblue}{}}}
\newtcolorbox{bilan}{popbase,colframe=popink,colback=white,boxrule=2.8pt,
  before upper={\popboxhead{Fiche bilan}{poporange}{L'essentiel à connaître pour l'examen}}}
% Figure encadrée avec légende
\newtcolorbox{popfigure}[1][]{enhanced,breakable=false,colback=white,colframe=popline,boxrule=1.4pt,arc=8pt,
  left=10pt,right=10pt,top=10pt,bottom=8pt,before skip=10pt,after skip=12pt,halign=center,
  lower separated=false,halign lower=center,
  fontlower=\popsemi\fontsize{9}{11.5}\selectfont\color{popmuted},
  #1}
\newcommand{\legende}[1]{\par\vspace{6pt}{\popsemi\fontsize{9}{11.5}\selectfont\color{popmuted}#1}}

% ============================================================
% Signature OnBuch+
% ============================================================
\newcommand{\poplogoinline}[1]{{\poplogo\fontsize{#1}{#1}\selectfont\color{popink}OnBuch\color{poporange}+}}
\newcommand{\signatureonbuch}{%
  \begin{tcolorbox}[enhanced,colback=popink,colframe=popink,boxrule=2.2pt,arc=10pt,
      drop shadow=poporange,left=14pt,right=14pt,top=10pt,bottom=10pt,before skip=0pt,after skip=0pt]
    \tikz[baseline=-0.7ex]\node[circle,fill=popgreen,draw=popcream,line width=1.3pt,minimum size=22pt,inner sep=0]{\popcheck{white}};%
    \hspace{9pt}\begin{minipage}[c]{0.62\linewidth}
      {\popxbold\fontsize{12}{13}\selectfont\color{popcream}Signed\ \textbullet\ {\poplogo OnBuch\color{poporange}+}}\par\vspace{2pt}
      {\raggedright\popsemi\fontsize{8}{10.5}\selectfont\color{popline}\MakeUppercase{OnBuch+ teaching team}\\\MakeUppercase{Follows the official MINESEC syllabus}\par}
    \end{minipage}\hfill
    {\poplogo\fontsize{22}{22}\selectfont\color{popcream}OnBuch\color{poporange}+}
  \end{tcolorbox}%
}

% ============================================================
% Page de garde
% #1 titre · #2 matière · #3 niveau · #4 module · #5 n° leçon · #6 durée ·
% #7 description / objectifs (texte ou itemize)
% ============================================================
\newcommand{\pagegarde}[7]{%
  \thispagestyle{empty}
  \newgeometry{a4paper,margin=1.7cm,top=1.6cm,bottom=1.4cm}%
  \begin{tikzpicture}[remember picture,overlay]
    \fill[poporange!18] ([xshift=-3.2cm,yshift=-3.6cm]current page.north east) circle (4.6cm);
    \fill[poppurple!14] ([xshift=2.4cm,yshift=7.6cm]current page.south west) circle (3.8cm);
  \end{tikzpicture}%
  {\poplogo\fontsize{30}{30}\selectfont\color{popink}OnBuch\color{poporange}+}\hfill
  \raisebox{0.6ex}{\poppill{Signed course \textbullet\ Premium}{popink}{popcream}}\par
  \vspace{1pt}\popsquiggle{poppurple}\par
  \vspace{8pt}
  \begin{tcolorbox}[enhanced,colback=poporange,colframe=popink,boxrule=2.8pt,arc=13pt,
      drop shadow=popink,left=18pt,right=18pt,top=12pt,bottom=13pt,after skip=10pt]
    \poppill{#2 \textbullet\ #3}{popink}{popcream}\hspace{5pt}\poppill{#4}{white}{popink}\par\vspace{8pt}
    {\popdisplay\fontsize{14}{14}\selectfont\color{popink}CHAPTER #5}\par\vspace{4pt}
    {\raggedright\hyphenpenalty=10000\exhyphenpenalty=10000\popdisplay\fontsize{25}{27}\selectfont\color{popcream}\MakeUppercase{#1}\par}
    \vspace{8pt}
    {\popxbold\fontsize{9.5}{9.5}\selectfont\color{popink}\MakeUppercase{Suggested time: #6}}
  \end{tcolorbox}
  \begin{tcolorbox}[enhanced,colback=white,colframe=popink,boxrule=2.2pt,arc=10pt,
      drop shadow=popink,left=16pt,right=16pt,top=10pt,bottom=10pt,after skip=8pt]
    \poppill{In this chapter}{poppurple}{white}\par\vspace{5pt}
    \popsemi\fontsize{9.3}{12.6}\selectfont\color{popink}#7
  \end{tcolorbox}
  \noindent\begin{minipage}[t]{0.487\linewidth}
    \begin{tcolorbox}[enhanced,colback=popgreen,colframe=popink,boxrule=2.2pt,arc=10pt,
        drop shadow=popink,left=11pt,right=11pt,top=10pt,bottom=10pt,height=3.1cm,valign=center]
      \tcbox[colback=white,colframe=popink,boxrule=1.3pt,arc=5pt,boxsep=0pt,nobeforeafter,
        left=3pt,right=3pt,top=3pt,bottom=3pt,valign=center]{\qrcode[height=1.9cm]{https://play.google.com/store/apps/details?id=com.luvvix.onbuch&hl=en}}%
      \hspace{8pt}\begin{minipage}[c]{0.5\linewidth}
        {\popdisplay\fontsize{11}{12.5}\selectfont\color{white}\MakeUppercase{Download OnBuch}}\par\vspace{4pt}
        {\raggedright\popsemi\fontsize{8.4}{11}\selectfont\color{white}Free \textbullet\ Google Play\\AI tutor, past papers, results\par}
      \end{minipage}
    \end{tcolorbox}
  \end{minipage}\hfill
  \begin{minipage}[t]{0.487\linewidth}
    \begin{tcolorbox}[enhanced,colback=poppurple,colframe=popink,boxrule=2.2pt,arc=10pt,
        drop shadow=popink,left=11pt,right=11pt,top=10pt,bottom=10pt,height=3.1cm,valign=center]
      \tikz[baseline=-0.7ex]\node[circle,fill=white,draw=popink,line width=1.3pt,minimum size=1.6cm,inner sep=0]{\popdisplay\fontsize{24}{24}\selectfont\color{poppurple}+};%
      \hspace{8pt}\begin{minipage}[c]{0.6\linewidth}
        {\popdisplay\fontsize{11}{12.5}\selectfont\color{white}\MakeUppercase{Go OnBuch+}}\par\vspace{4pt}
        {\raggedright\popsemi\fontsize{8.4}{11}\selectfont\color{white}All signed courses, unlimited AI tutor, PDF sheets — OnBuch+ tab in the app\par}
      \end{minipage}
    \end{tcolorbox}
  \end{minipage}
  \vspace{8pt}
  \popcontenu
  \vfill
  \signatureonbuch
  \restoregeometry
  \newpage
}

% Rangée « Ce cours contient » (page de garde)
\newcommand{\poptile}[3]{% #1 couleur #2 gros texte #3 légende
  \begin{tcolorbox}[enhanced,colback=#1,colframe=popink,boxrule=2pt,arc=9pt,shadow={2.5pt}{-2.5pt}{0pt}{popink},
      left=6pt,right=6pt,top=9pt,bottom=9pt,halign=center,valign=center,height=1.75cm,width=\linewidth,nobeforeafter]
    {\popdisplay\fontsize{12.5}{13}\selectfont\color{white}\MakeUppercase{#2}}\par\vspace{3pt}
    {\centering\popsemi\fontsize{7.8}{9.5}\selectfont\color{white}#3\par}
  \end{tcolorbox}}
\newcommand{\popcontenu}{%
  \par\noindent{\popxboldls\fontsize{7.5}{7.5}\selectfont\color{popmuted}THIS SIGNED COURSE CONTAINS}\par\vspace{7pt}
  \noindent\begin{minipage}[t]{0.235\linewidth}\poptile{popblue}{Course}{complete, illustrated, step by step}\end{minipage}\hfill
  \begin{minipage}[t]{0.235\linewidth}\poptile{poporange}{Methods}{and worked examples}\end{minipage}\hfill
  \begin{minipage}[t]{0.235\linewidth}\poptile{poppink}{Exercises}{exam style + solutions}\end{minipage}\hfill
  \begin{minipage}[t]{0.235\linewidth}\poptile{popgreen}{Summary sheet}{the essentials to remember}\end{minipage}\par}

% Sommaire
\newtcolorbox{sommaire}{enhanced,colback=white,colframe=popink,boxrule=2.2pt,arc=10pt,
  drop shadow=popink,left=18pt,right=18pt,top=13pt,bottom=13pt,before skip=6pt,after skip=20pt,
  before upper={\poppill{Contents}{poppurple}{white}\par\vspace{9pt}\popsemi\fontsize{10.6}{14.5}\selectfont\color{popink}}}

% Signature de fin de cours — poussée tout en bas de la dernière page
\newcommand{\findecours}{%
  \par\vfill\nopagebreak
  \begin{center}\popsquiggle{poporange}\end{center}\vspace{4pt}
  \signatureonbuch
}
\AtBeginDocument{\def\llbracket{\mathopen{[\mkern-3.5mu[}}\def\rrbracket{\mathclose{]\mkern-3.5mu]}}} % intervalles d'entiers (glyphe ⟦ absent de la police)
% Glyphes absents de Fira Math : redéfinis à partir de glyphes disponibles
\AtBeginDocument{%
  \def\setminus{\mathbin{\backslash}}\let\smallsetminus\setminus
  \def\vdots{\mathord{\vbox{\baselineskip3.2pt\lineskiplimit0pt\kern4pt\hbox{.}\hbox{.}\hbox{.}}}}%
  \def\ddots{\mathinner{\mkern1mu\raise6.5pt\hbox{.}\mkern2mu\raise3.6pt\hbox{.}\mkern2mu\raise0.7pt\hbox{.}\mkern1mu}}%
  \def\triangle{\mathord{\scalebox{0.9}{$\symup{\Delta}$}}}%
  \def\nabla{\mathord{\raisebox{\depth}{\scalebox{1}[-1]{$\symup{\Delta}$}}}}%
  \def\checkmark{\popcheck{popdark}}%
  \def\star{\mathbin{\ast}}\def\bigstar{\mathord{\ast}}%
  \def\diamond{\mathbin{\scalebox{0.6}{$\Diamond$}}}%
  \def\bigcup{\mathop{\mathchoice{\raisebox{-0.3em}{\scalebox{1.7}{$\cup$}}}{\scalebox{1.25}{$\cup$}}{\cup}{\cup}}\displaylimits}%
  \def\bigcap{\mathop{\mathchoice{\raisebox{-0.3em}{\scalebox{1.7}{$\cap$}}}{\scalebox{1.25}{$\cap$}}{\cap}{\cap}}\displaylimits}%
  \def\lesseqqgtr{\mathrel{\lessgtr}}\def\bigoplus{\mathop{\oplus}\displaylimits}%
}
\providecommand{\ssi}{if and only if} % abréviation fréquente
\AtBeginDocument{\providecommand{\wideparen}{\overparen}} % arc de cercle

%% FIGFIX-BEGIN v3 — correctifs globaux des figures (docs/onbuch-generation/tools/figfix)
\usepackage{adjustbox}\usepackage{etoolbox}
% toute figure TikZ/pgfplots plus large que la ligne est réduite à la largeur disponible
\BeforeBeginEnvironment{tikzpicture}{\begin{adjustbox}{max width=\linewidth}}
\AfterEndEnvironment{tikzpicture}{\end{adjustbox}}
% légendes pgfplots lisibles par-dessus les courbes
\pgfplotsset{every axis/.append style={legend style={fill=white,fill opacity=0.92,text opacity=1,draw=black!15,inner sep=3pt}}}
% étiquettes d'arêtes (syntaxe edge["..."]) avec fond blanc
\tikzset{every edge quotes/.append style={fill=white,inner sep=1.2pt}}
% bulles de cartes mentales : texte à la ligne dans la bulle, bulle un peu plus grande
\tikzset{every concept/.append style={text width=2.0cm,align=center,minimum size=2.3cm,inner sep=2pt}}
%% FIGFIX-END
\begin{document}

\pagegarde{Work, Energy and Power}{Pure Mathematics with Mechanics}{Upper Sixth}{Upper Sixth — Pure mathematics with mechanics}{15}{10 periods}{This chapter develops the concepts of work, energy and power for constant and variable forces, including motion under gravity and against resistance. It equips Upper Sixth students with the work–energy theorem, conservation of mechanical energy and power–velocity relations needed to solve GCE Advanced Level mechanics problems.
\vspace{4pt}
\begin{multicols}{2}\small
\begin{itemize}
\item By the end of this chapter I can define the work done by a constant force and compute it using W = Fs cos θ.
\item By the end of this chapter I can calculate the work done against gravity when a body is raised through a vertical height.
\item By the end of this chapter I can evaluate the work done by a variable force using integration, W = ∫F dx.
\item By the end of this chapter I can define kinetic and potential energy and compute their values in standard situations.
\item By the end of this chapter I can state and apply the work–energy theorem to moving bodies.
\item By the end of this chapter I can apply the principle of conservation of mechanical energy to systems with and without friction.
\end{itemize}\end{multicols}}

\begin{sommaire}
\begin{multicols}{2}
\begin{enumerate}
\item Work Done by a Constant Force and Work Done Against Gravity
\item Work Done by Variable Forces
\item Energy: Kinetic and Potential
\item The Work–Energy Theorem and Conservation of Energy
\item Power: Definition, Instantaneous Power and Power–Velocity
\item Pumping Water, Water Jets and Applications of Work, Energy and Power
\item Integration activity
\item Methods and worked examples
\item Exercises
\item Solutions to the exercises
\item Summary sheet
\end{enumerate}
\end{multicols}
\end{sommaire}

\begin{objectifs}
\begin{itemize}
\item By the end of this chapter I can define the work done by a constant force and compute it using W = Fs cos θ.
\item By the end of this chapter I can calculate the work done against gravity when a body is raised through a vertical height.
\item By the end of this chapter I can evaluate the work done by a variable force using integration, W = ∫F dx.
\item By the end of this chapter I can define kinetic and potential energy and compute their values in standard situations.
\item By the end of this chapter I can state and apply the work–energy theorem to moving bodies.
\item By the end of this chapter I can apply the principle of conservation of mechanical energy to systems with and without friction.
\item By the end of this chapter I can define power, distinguish average from instantaneous power, and use P = Fv.
\item By the end of this chapter I can solve problems on vehicles moving with constant power, including maximum speed on level roads and inclines.
\item By the end of this chapter I can model pumping water and water jets and compute the power required.
\item By the end of this chapter I can solve mixed GCE-style problems combining work, energy and power.
\end{itemize}
\end{objectifs}

\begin{prerequis}
\begin{itemize}
\item Kinematics of motion in a straight line (suvat equations) from Lower Sixth
\item Newton's laws of motion and resolution of forces on inclined planes
\item Definite integration and integration of polynomial functions (Pure Mathematics)
\item Trigonometric ratios and resolving vectors into components
\item Friction and the coefficient of friction (Lower Sixth mechanics)
\end{itemize}
\end{prerequis}

\newpage

\coursec{Work Done by a Constant Force and Work Done Against Gravity}

Every morning in Bamenda, water vendors push heavy carts up steep streets, and at the Memve'ele hydroelectric dam, falling water spins turbines that light up entire towns. How much ``effort'' does the vendor really expend climbing the hill, and how much energy does the falling water deliver each second? A generator mechanic in Douala chooses between two engines rated in kilowatts --- what does that number actually measure? In this chapter we turn these everyday questions into precise mathematics: \cle{work}, \cle{energy} and \cle{power}.

We begin with the simplest situation: a single, constant force acting on a body that moves in a straight line. Our goal in this section is to define work precisely, to learn when work is positive, negative or zero, and to master the most important special case at GCE level --- the work done against gravity.

\courssub{From Everyday Effort to a Mathematical Definition}

Imagine a porter at the Yaound\'e central market dragging a sack of rice along the ground. Two things clearly determine how ``tiring'' the task is: how hard he pulls (the \cle{force}) and how far he drags the sack (the \cle{displacement}). But there is a third, subtler ingredient. If he pulls \emph{horizontally}, the sack slides easily; if he pulls almost \emph{vertically}, the sack barely moves forward, however hard he strains. Only the part of the force \emph{in the direction of motion} is useful in moving the sack. This observation is the key to the definition of work.

\begin{definition}[Work done by a constant force]
The \cle{work done} by a constant force $\vec{F}$ whose point of application moves through a displacement $\vec{s}$ is the scalar product
\[
W = \vec{F} \cdot \vec{s} = Fs\cos\theta,
\]
where $F = |\vec{F}|$, $s = |\vec{s}|$ and $\theta$ is the angle between the force and the displacement ($0 \le \theta \le \pi$). The SI unit of work is the \cle{joule} (symbol $\mathrm{J}$): $1\ \mathrm{J} = 1\ \mathrm{N\,m}$.
\end{definition}

The factor $\cos\theta$ selects exactly the component of the force along the direction of motion, $F\cos\theta$. So we may read the formula in two equivalent ways:
\[
W = \underbrace{(F\cos\theta)}_{\text{component of force along motion}} \times s
\qquad\text{or}\qquad
W = F \times \underbrace{(s\cos\theta)}_{\text{displacement along the force}}.
\]

\begin{popfigure}
\begin{tikzpicture}[scale=1.0, >=stealth]
% ground
\draw[thick, popmuted] (-0.5,0) -- (7.5,0);
\foreach \x in {0,0.5,...,7} \draw[popmuted] (\x,0) -- (\x-0.3,-0.25);
% block
\fill[popblueL] (1,0) rectangle (2.6,1.1);
\draw[thick, popink] (1,0) rectangle (2.6,1.1);
% displacement arrow
\draw[->, very thick, popgreen] (1.8,-0.7) -- (6.3,-0.7) node[midway, below, popgreen]{$\vec{s}$};
% force arrow at angle
\draw[->, very thick, poporange] (1.8,0.55) -- (4.6,2.35) node[above right, poporange]{$\vec{F}$};
% components
\draw[->, thick, poppurple, dashed] (1.8,0.55) -- (4.6,0.55) node[midway, below, poppurple]{$F\cos\theta$};
\draw[->, thick, poppurple, dashed] (4.6,0.55) -- (4.6,2.35) node[midway, right, poppurple]{$F\sin\theta$};
% angle
\draw[popdark] (3.3,0.55) arc (0:38:1.5) node[midway, right, popdark]{$\theta$};
\end{tikzpicture}
\legende{A block pulled by a constant force $\vec{F}$ at angle $\theta$ to the horizontal displacement $\vec{s}$. Only the component $F\cos\theta$ does work.}
\end{popfigure}

\begin{attention}[Work is a scalar]
Although force and displacement are vectors, work is a \cle{scalar}: it has magnitude and sign, but no direction. We add works algebraically, never by the parallelogram rule. Note also that $W$ can be \emph{negative} --- a negative work is not ``less than nothing''; it means the force removes energy from the body rather than giving it energy.
\end{attention}

\courssub{Positive, Negative and Zero Work}

The sign of $W = Fs\cos\theta$ is decided entirely by the angle $\theta$, since $F$ and $s$ are positive magnitudes.

\begin{propriete}[Sign of the work]
Let a constant force $\vec{F}$ act on a body moving through displacement $\vec{s}$, with $\theta$ the angle between them.
\begin{itemize}
\item If $0 \leq \theta < 90^\circ$, then $\cos\theta > 0$ and $W > 0$: the force \emph{helps} the motion (positive work).
\item If $\theta = 90^\circ$, then $\cos\theta = 0$ and $W = 0$: a force perpendicular to the displacement does \emph{no work}.
\item If $90^\circ < \theta \leq 180^\circ$, then $\cos\theta < 0$ and $W < 0$: the force \emph{opposes} the motion (negative work). Friction is the standard example.
\end{itemize}
\end{propriete}

\begin{exemplebox}[Zero work: carrying a load horizontally]
A student carries a $10\ \mathrm{kg}$ bag of books horizontally across the school yard. The upward force she exerts on the bag is perpendicular to the horizontal displacement, so $\theta = 90^\circ$ and the work done by her supporting force is $W = 0$. She gets tired --- but in the mechanical sense, her upward force does no work on the bag.
\end{exemplebox}

\begin{exoresolu}[Pushing a cart in Bamenda]
A water vendor pushes his cart with a constant horizontal force of $120\ \mathrm{N}$ along a straight street for $50\ \mathrm{m}$. A resistive force (friction and air) of $40\ \mathrm{N}$ opposes the motion. Find the work done by each force and the total work done on the cart.
\tcblower
\textbf{Solution.} The pushing force is parallel to the displacement, so $\theta = 0^\circ$:
\[
W_{\text{push}} = Fs\cos 0^\circ = 120 \times 50 \times 1 = 6000\ \mathrm{J}.
\]
The resistive force points opposite to the motion, so $\theta = 180^\circ$:
\[
W_{\text{res}} = 40 \times 50 \times \cos 180^\circ = 40 \times 50 \times (-1) = -2000\ \mathrm{J}.
\]
The total (net) work is the algebraic sum:
\[
W_{\text{total}} = 6000 + (-2000) = 4000\ \mathrm{J}.
\]
The cart gains $4000\ \mathrm{J}$ of energy from the vendor's effort; the remaining $2000\ \mathrm{J}$ is dissipated by the resistive forces.
\end{exoresolu}

\begin{propriete}[Total work is an algebraic sum]
When several forces $\vec{F}_1, \vec{F}_2, \dots, \vec{F}_n$ act simultaneously on a body, the total work done on the body over a displacement is the algebraic sum of the individual works:
\[
W_{\text{total}} = W_1 + W_2 + \cdots + W_n.
\]
(Equivalently, the total work equals the work of the resultant force $\vec{F}_{\text{res}} = \sum \vec{F}_i$.) This property is used constantly at GCE level; its proof is not examinable.
\end{propriete}

\begin{exoresolu}[Dragging a sack of rice with a rope at an angle]
A porter drags a sack of rice across a horizontal floor using a rope inclined at $30^\circ$ above the horizontal. He pulls with a constant tension of $100\ \mathrm{N}$ over a distance of $12\ \mathrm{m}$. Find the work done by the tension.
\tcblower
\textbf{Solution.} The angle between the tension and the horizontal displacement is $\theta = 30^\circ$:
\[
W = Fs\cos\theta = 100 \times 12 \times \cos 30^\circ = 1200 \times \frac{\sqrt{3}}{2} = 600\sqrt{3} \approx 1039\ \mathrm{J}.
\]
Notice that the vertical component $F\sin 30^\circ = 50\ \mathrm{N}$ does no work at all: it merely lightens the sack on the floor.
\end{exoresolu}

\courssub{Work Done Against Gravity}

Lift a bag of cement from the ground onto a lorry and you feel immediately that you are working \emph{against} the weight of the bag. Let us make this precise.

Consider a body of mass $m$ raised slowly (at constant speed, so that the lifting force just balances the weight) through a vertical height $h$. The lifting force has magnitude $F = mg$ and acts vertically upward, in the same direction as the vertical displacement. Hence
\[
W = Fs\cos 0^\circ = mg \times h \times 1 = mgh.
\]

\begin{aretenir}[Work done against gravity]
The work done \emph{against gravity} in raising a body of mass $m$ through a vertical height $h$ is
\[
W_{\text{against gravity}} = mgh,
\]
where $g$ is the acceleration due to gravity (take $g = 9.81\ \mathrm{m\,s^{-2}}$, or $10\ \mathrm{m\,s^{-2}}$ when the question says so). 

Conversely, the work done \emph{by the weight} (gravitational force) when the body moves vertically is:
\begin{itemize}
\item $W_{\text{gravity}} = -mgh$ if the body \emph{rises} through height $h$ (force down, displacement up, $\theta=180^\circ$).
\item $W_{\text{gravity}} = +mgh$ if the body \emph{falls} through height $h$ (force and displacement both down, $\theta=0^\circ$).
\end{itemize}
\end{aretenir}

\begin{methode}[Computing work against gravity]
\begin{enumerate}
\item Identify the mass $m$ of the body (convert to kilograms if necessary).
\item Identify the \textbf{vertical} height $h$ gained (or lost) --- never the distance along a slope.
\item Compute $W = mgh$ with $g = 9.81\ \mathrm{m\,s^{-2}}$ (or the value given).
\item State the sign clearly: work done \emph{against} gravity is positive ($+mgh$); work done \emph{by} gravity is negative for a rise ($-mgh$) and positive for a fall ($+mgh$).
\end{enumerate}
\end{methode}

A remarkable and examinable fact is that the work done by gravity depends \emph{only} on the vertical displacement, not on the path taken. Whether the vendor carries his jerrycans straight up a staircase or along a long gentle ramp to the same floor, the work done against gravity is the same: $mgh$. To see why, split any path into tiny steps; the horizontal steps contribute nothing (the weight is vertical, perpendicular to horizontal motion), and the vertical steps add up to the total height $h$. Forces with this path-independence property are called \cle{conservative} forces, and gravity is the first example you meet.

\begin{attention}[The classic trap: height versus distance along the slope]
A very common mistake at GCE is to compute the work against gravity as $mg \times s$, where $s$ is the distance travelled \emph{along} an inclined plane. This is wrong. The work against gravity is $mgh$ where $h = s\sin\alpha$ is the \emph{vertical} height gained. Always extract the vertical height first.
\end{attention}

\courssub{Pulling a Body Up a Rough Inclined Plane}

We now combine everything: a crate is pulled at constant speed up a rough plane inclined at angle $\alpha$ to the horizontal, through a distance $s$ along the slope. Four forces act on the crate:
\begin{itemize}
\item the applied force $\vec{F}$, parallel to the slope (directed upward);
\item the weight $m\vec{g}$, vertically downward;
\item the normal reaction $\vec{R}$, perpendicular to the slope;
\item the friction $\vec{f}$, acting down the slope (opposing the motion).
\end{itemize}

\begin{popfigure}
\begin{tikzpicture}[scale=1.05, >=stealth]
% Inclined plane: alpha = 26.565 deg (tan = 0.5), length 7, height 3.5
\coordinate (O) at (0,0);
\coordinate (B) at (7,0);
\coordinate (T) at (7,3.5);
% Plane surface
\draw[thick, popmuted] (O) -- (T);
% Horizontal base
\draw[thick, popmuted] (O) -- (B);
% Hatch marks on plane (manual, no calc lib)
\foreach \d in {0.8, 1.8, 2.8, 3.8, 4.8, 5.8, 6.8} {
  \draw[popmuted] (\d, \d*0.5) -- ++(-0.2, -0.4);
}
% Angle alpha
\draw[popdark] (1.5,0) arc (0:26.565:1.5) node[midway, right, popdark] at (1.8,0.3) {$\alpha$};
% Crate on slope (center at 3.5, 1.75)
% Crate dimensions: width 1.2 along slope, height 1.0 perpendicular
% Unit vector along slope: (cos a, sin a) = (0.894, 0.447)
% Unit vector normal: (-sin a, cos a) = (-0.447, 0.894)
% Center C = (3.5, 1.75)
% Corners relative to center: +/-0.6 along slope, +/-0.5 normal
\coordinate (C) at (3.5, 1.75);
\pgfmathsetmacro{\cx}{3.5}
\pgfmathsetmacro{\cy}{1.75}
\pgfmathsetmacro{\ca}{0.8944}
\pgfmathsetmacro{\sa}{0.4472}
\pgfmathsetmacro{\wx}{0.6*\ca}
\pgfmathsetmacro{\wy}{0.6*\sa}
\pgfmathsetmacro{\nx}{-0.5*\sa}
\pgfmathsetmacro{\ny}{0.5*\ca}
% Four corners
\coordinate (CR1) at ({\cx+\wx+\nx}, {\cy+\wy+\ny});
\coordinate (CR2) at ({\cx-\wx+\nx}, {\cy-\wy+\ny});
\coordinate (CR3) at ({\cx-\wx-\nx}, {\cy-\wy-\ny});
\coordinate (CR4) at ({\cx+\wx-\nx}, {\cy+\wy-\ny});
\fill[popblueL] (CR1) -- (CR2) -- (CR3) -- (CR4) -- cycle;
\draw[thick, popink] (CR1) -- (CR2) -- (CR3) -- (CR4) -- cycle;
% Applied force F (up slope)
\draw[->, very thick, poporange] (C) -- ++({\ca*2.2}, {\sa*2.2}) node[above right, poporange]{$\vec{F}$};
% Friction f (down slope)
\draw[->, very thick, poppink] (C) -- ++({-\ca*1.5}, {-\sa*1.5}) node[below left, poppink]{$\vec{f}$};
% Normal reaction R
\draw[->, very thick, popblue] (C) -- ++({\nx*3.2}, {\ny*3.2}) node[above left, popblue]{$\vec{R}$};
% Weight mg
\draw[->, very thick, poppurple] (C) -- ++(0, -2.0) node[below, poppurple]{$m\vec{g}$};
% Displacement s
\draw[->, thick, popgreen, dashed] (0.5, 0.25) -- ++({\ca*4.5}, {\sa*4.5}) node[midway, below right, popgreen]{$s$};
% Height h
\draw[<->, thick, popdark] (7.2, 0) -- (7.2, 3.5) node[midway, right, popdark]{$h = s\sin\alpha$};
\end{tikzpicture}
\legende{A crate pulled up a rough plane inclined at angle $\alpha$. The vertical height gained is $h = s\sin\alpha$; friction acts down the slope, opposing the motion.}
\end{popfigure}

Let us compute the work of each force over the displacement $s$ along the plane.

\textbf{Work of the applied force.} $\vec{F}$ is parallel to the motion (assumed up the slope): $\theta = 0^\circ$, so $W_F = Fs$.

\textbf{Work of the normal reaction.} $\vec{R}$ is perpendicular to the slope, hence to the displacement: $\theta = 90^\circ$, so $W_R = 0$.

\textbf{Work of friction.} Friction opposes the motion, $\theta = 180^\circ$: $W_f = -fs$. The work done \emph{against} friction is $+fs$.

\textbf{Work of the weight.} 
Method 1 (Angle): The angle between the vertically downward weight and the displacement up the slope is $90^\circ + \alpha$, so
\[
W_g = mgs\cos(90^\circ + \alpha) = -mgs\sin\alpha = -mgh.
\]
Method 2 (Components): Resolve weight into components parallel ($mg\sin\alpha$ down the slope) and perpendicular ($mg\cos\alpha$) to the plane. The perpendicular component does no work. The parallel component acts opposite to displacement, so $W_g = (-mg\sin\alpha) \times s = -mgs\sin\alpha = -mgh$.
Both methods confirm that only the vertical height $h = s\sin\alpha$ matters. The work done \emph{against} gravity is $+mgh$.

\begin{propriete}[Work needed to pull a body up a rough plane at constant speed]
If a body of mass $m$ is pulled at \emph{constant speed} through a distance $s$ up a rough plane inclined at $\alpha$, with frictional force $f$, then the applied force does work
\[
W_F = \underbrace{mgh}_{\text{against gravity}} + \underbrace{fs}_{\text{against friction}}, \qquad h = s\sin\alpha.
\]
\textbf{Derivation (examinable):} At constant speed, acceleration is zero, so the resultant force parallel to the plane is zero. 
$F - mg\sin\alpha - f = 0 \implies F = mg\sin\alpha + f$.
Multiplying by the displacement $s$:
$W_F = Fs = mg s\sin\alpha + fs = mgh + fs$.
\end{propriete}

\begin{exoresolu}[Crate up a rough ramp]
A crate of mass $50\ \mathrm{kg}$ is pulled at constant speed $8\ \mathrm{m}$ up a ramp inclined at $25^\circ$ to the horizontal. The frictional force between crate and ramp is $90\ \mathrm{N}$. Taking $g = 9.81\ \mathrm{m\,s^{-2}}$, find (a) the work done against gravity, (b) the work done against friction, (c) the total work done by the pulling force.
\tcblower
\textbf{Solution.} The vertical height gained is
\[
h = s\sin\alpha = 8\sin 25^\circ = 8 \times 0.4226 \approx 3.381\ \mathrm{m}.
\]
(a) Work done against gravity:
\[
W_g = mgh = 50 \times 9.81 \times 3.381 \approx 1658\ \mathrm{J}.
\]
(b) Work done against friction:
\[
W_f = fs = 90 \times 8 = 720\ \mathrm{J}.
\]
(c) Since the speed is constant, the pulling force supplies both:
\[
W_F = W_g + W_f = 1658 + 720 = 2378\ \mathrm{J} \approx 2.38\ \mathrm{kJ}.
\]
Note the signs: the work done \emph{by} gravity is $-1658\ \mathrm{J}$ and the work done \emph{by} friction is $-720\ \mathrm{J}$, while the work done by the pulling force is $+2378\ \mathrm{J}$; they sum to zero, consistent with constant speed (Work-Energy Theorem).
\end{exoresolu}

\begin{attention}[GCE exam tip: state every sign]
In GCE Advanced Level questions, examiners award marks for the \emph{sign} of each work term. Always write explicitly: ``work done by friction $= -fs$ (negative, because friction opposes the motion)'' and ``work done by gravity $= -mgh$ (the body rises)''. A correct magnitude with a missing or wrong sign loses marks.
\end{attention}

\begin{savaistu}[The joule in daily life]
The joule is named after the English physicist James Prescott Joule, who showed that mechanical work and heat are two faces of the same thing: energy. One joule is roughly the work needed to lift a $100\ \mathrm{g}$ tomato by one metre --- so the $2378\ \mathrm{J}$ spent hauling the crate up the ramp is like lifting that tomato nearly two and a half thousand times!
\end{savaistu}

\begin{exercice}[Lifting water at the well]
A girl draws a bucket of water of mass $12\ \mathrm{kg}$ from a well $15\ \mathrm{m}$ deep, at constant speed, using a rope. Taking $g = 9.81\ \mathrm{m\,s^{-2}}$, find (a) the work she does against gravity; (b) the work done by gravity on the bucket during the lift; (c) the work done by gravity if the full bucket is then lowered back down to the water surface.
\end{exercice}

\begin{corrige}[Exercise 1]
(a) Work against gravity: $W = mgh = 12 \times 9.81 \times 15 = 1765.8\ \mathrm{J} \approx 1766\ \mathrm{J}$.
(b) The bucket rises while gravity acts downward, so the work done by gravity is $-1765.8\ \mathrm{J}$.
(c) On the way down, gravity acts in the direction of motion, so the work done by gravity is $+1765.8\ \mathrm{J}$. Over the whole round trip, the total work done by gravity is zero --- another illustration that gravity is conservative.
\end{corrige}

\begin{exercice}[Sack on a building site]
A worker pulls a sack of cement of mass $40\ \mathrm{kg}$ at constant speed $10\ \mathrm{m}$ up a plank inclined at $30^\circ$ to the horizontal. Friction between sack and plank is $60\ \mathrm{N}$. Take $g = 10\ \mathrm{m\,s^{-2}}$. Find (a) the vertical height gained; (b) the work done against gravity; (c) the work done against friction; (d) the work done by the worker.
\end{exercice}

\begin{corrige}[Exercise 2]
(a) $h = s\sin\alpha = 10\sin 30^\circ = 5\ \mathrm{m}$.
(b) $W_g = mgh = 40 \times 10 \times 5 = 2000\ \mathrm{J}$.
(c) $W_f = fs = 60 \times 10 = 600\ \mathrm{J}$.
(d) At constant speed, $W = W_g + W_f = 2000 + 600 = 2600\ \mathrm{J}$.
\end{corrige}

\begin{aretenir}[Key points of this section]
\begin{itemize}
\item Work of a constant force: $W = \vec{F}\cdot\vec{s} = Fs\cos\theta$, in joules ($1\ \mathrm{J} = 1\ \mathrm{N\,m}$); work is a scalar.
\item $\theta < 90^\circ$: positive work; $\theta = 90^\circ$: zero work; $\theta > 90^\circ$: negative work (e.g., friction).
\item Total work $=$ algebraic sum of individual works $=$ work of resultant force.
\item Work against gravity $= mgh$, where $h$ is the \emph{vertical} height; gravity's work is path-independent (conservative force).
\item Up a rough plane at constant speed: work of the applied force $= mgh + fs$, with $h = s\sin\alpha$.
\end{itemize}
\end{aretenir}

\coursec{Work Done by Variable Forces}

\courssub{From Constant Forces to Variable Forces}

In the previous section we established that when a \cle{constant} force $F$ moves a body through a displacement $s$ along its line of action, the work done is $W = Fs$. This formula served us well for a man pushing a wheelbarrow with a steady force, or for the work done against gravity when lifting a bag of cement at constant speed, where the lifting force equals the constant weight $mg$.

But look around you in everyday Cameroonian life and you will quickly notice that truly constant forces are the exception rather than the rule. When a trader at Mokolo market stretches the rubber band holding a bundle of vegetables, the band pulls back harder and harder the more it is stretched. When a mechanic compresses the suspension spring of a \emph{bendskin} (motorbike), the resistance grows with compression. When a crane slowly lifts a long chain from the ground, the weight being lifted increases as more links leave the floor. In all these situations the force \emph{depends on position}, and the simple product $W = Fs$ no longer applies.

\begin{attention}[Why $W = Fs$ fails for a variable force]
The formula $W = Fs$ was derived under the hypothesis that $F$ keeps the \emph{same value} throughout the displacement. If $F$ changes with position $x$, which value of $F$ should we multiply by $s$? The initial value underestimates the work, the final value overestimates it, and no single value is correct. We need a new tool — and that tool is \cle{integration}.
\end{attention}

\courssub{Work as the Area Under the Force--Displacement Graph}

Let us build the idea step by step. Suppose a force $F(x)$ acts on a body moving along the $x$-axis from $x = a$ to $x = b$. Divide the interval $[a, b]$ into many tiny sub-intervals of width $\delta x$. On each tiny sub-interval, the force barely changes, so it is \emph{approximately constant} there, and the work done across that small piece is approximately
\[
\delta W \approx F(x)\,\delta x .
\]
But $F(x)\,\delta x$ is precisely the area of a thin rectangle of height $F(x)$ and width $\delta x$ under the graph of $F$ against $x$. Adding all these small contributions and letting $\delta x \to 0$, the total work becomes the \emph{exact area} under the force--displacement curve between $a$ and $b$ — which is exactly what the definite integral computes.

\begin{definition}[Work done by a variable force]
The \cle{work done} by a variable force $F(x)$ acting along the direction of motion in moving a body from position $x = a$ to position $x = b$ is
\[
W = \int_{a}^{b} F(x)\, dx .
\]
Geometrically, $W$ is the \cle{area under the force--displacement graph} between $x = a$ and $x = b$. The unit of work remains the joule ($1\ \mathrm{J} = 1\ \mathrm{N\,m}$).
\end{definition}

\begin{aretenir}[Key points]
\begin{itemize}
\item Constant force: $W = Fs$ (a special case of the integral, since $\int_0^s F\,dx = Fs$).
\item Variable force: $W = \displaystyle\int_a^b F(x)\,dx$ = area under the $F$--$x$ graph.
\item If the force opposes the motion, the work done \emph{by} the force is negative; the work done \emph{against} it is positive.
\end{itemize}
\end{aretenir}

\begin{exoresolu}[Work done by a force given as a function of $x$]
A force of magnitude $F(x) = 3x^2 + 2x$ newtons acts on a particle moving along the $x$-axis, in the direction of motion. Calculate the work done by the force as the particle moves from $x = 1\ \mathrm{m}$ to $x = 3\ \mathrm{m}$.
\tcblower
\textbf{Solution.} The force varies with position, so we integrate:
\begin{align*}
W &= \int_{1}^{3} \left(3x^2 + 2x\right) dx \\
&= \left[ x^3 + x^2 \right]_{1}^{3} \\
&= \left(3^3 + 3^2\right) - \left(1^3 + 1^2\right) \\
&= (27 + 9) - (1 + 1) = 36 - 2 = 34 .
\end{align*}
The work done by the force is $\boxed{W = 34\ \mathrm{J}}$.
\end{exoresolu}

\courssub{Hooke's Law: Springs and Elastic Strings}

The most important example of a variable force at GCE Advanced Level is the tension in a \cle{spring} or \cle{elastic string}. Take a light spring fixed at one end and pull the free end horizontally. The more you stretch it, the harder it pulls back. This experimental observation is summarised by a famous law.

\begin{experience}[Stretching a spring]
Hang a light spring vertically, fix a pointer to its lower end, and add slotted masses one at a time, recording the extension $x$ each time. Plotting the tension $T = mg$ against $x$ gives, for moderate extensions, a \emph{straight line through the origin}: the tension is proportional to the extension. Beyond a certain load (the \cle{elastic limit}), the line curves and the spring is permanently deformed.
\end{experience}

\begin{propriete}[Hooke's Law]
When a light elastic spring (or string) of natural length $l$ is stretched by an \cle{extension} $x$ (within its elastic limit), the tension pulling back towards the natural length is
\[
T = kx , \qquad\text{equivalently}\qquad T = \frac{\lambda}{l}\,x ,
\]
where $k$ is the \cle{stiffness} (or spring constant, in $\mathrm{N\,m^{-1}}$) and $\lambda$ is the \cle{modulus of elasticity} (in newtons). The same law holds for compression of a spring, with $x$ the compression. (The statement of Hooke's law is examinable; its experimental verification above is a useful guided activity.)
\end{propriete}

\begin{attention}[Extension, not total length]
In $T = kx$, the quantity $x$ is the \emph{extension}, i.e.\ (stretched length) $-$ (natural length). A very common examination error is to substitute the total stretched length for $x$. Always compute $x = \text{new length} - l$ first.
\end{attention}

\begin{popfigure}
\begin{center}
\begin{tikzpicture}[scale=1.0, every node/.style={font=\small}]
% wall
\fill[popmuted] (0,-0.9) rectangle (0.25,0.9);
\draw[popdark, thick] (0.25,-0.9) -- (0.25,0.9);
% natural length spring
\draw[popblue, thick, decorate, decoration={coil, aspect=0.4, segment length=4pt, amplitude=4pt}] (0.25,0.45) -- (3.2,0.45);
\draw[popdark, thick] (3.2,0.45) -- (3.7,0.45);
\fill[popblue] (3.7,0.45) circle (2pt);
\draw[popdark, <->] (0.25,1.15) -- (3.7,1.15) node[midway, above]{natural length $l$};
% stretched spring
\draw[poporange, thick, decorate, decoration={coil, aspect=0.4, segment length=7pt, amplitude=4pt}] (0.25,-0.55) -- (4.9,-0.55);
\draw[popdark, thick] (4.9,-0.55) -- (5.5,-0.55);
\fill[poporange] (5.5,-0.55) circle (2pt);
\draw[poporange, ->, thick] (5.5,-0.55) -- (6.6,-0.55) node[right]{$T = kx$};
\draw[popdark, dashed] (3.7,0.1) -- (3.7,-1.25);
\draw[popdark, <->] (3.7,-1.15) -- (5.5,-1.15) node[midway, below]{extension $x$};
\end{tikzpicture}
\end{center}
\legende{A light spring fixed at one end: at natural length $l$ (top) and stretched by $x$ (bottom), where it exerts a tension $T = kx$ directed back towards the natural length.}
\end{popfigure}

\courssub{Work Done in Stretching a Spring}

To stretch a spring slowly by an extension $x$, you must apply a force equal and opposite to the tension, i.e.\ an applied force $F(x) = kx$. This force is \emph{not constant}: it starts at $0$ and grows linearly to $ke$ when the extension reaches its final value $e$. The work you do is therefore given by integration.

\begin{propriete}[Work done in stretching a spring from its natural length]
The work done in stretching a light spring of stiffness $k$ from its natural length to an extension $e$ is
\[
W = \tfrac{1}{2} k e^2 .
\]
This work is stored in the spring as \cle{elastic potential energy} and is fully recoverable when the spring returns to its natural length (within the elastic limit). (The derivation below is examinable.)
\end{propriete}

\textbf{Derivation.} The applied force at extension $x$ is $F(x) = kx$. Hence
\begin{align*}
W &= \int_{0}^{e} kx\, dx = k\left[\frac{x^2}{2}\right]_{0}^{e} = \frac{1}{2}k e^2 . \qquad\blacksquare
\end{align*}

\begin{popfigure}
\begin{center}
\begin{tikzpicture}
\begin{axis}[
  width=9.5cm, height=6.5cm,
  axis lines=left,
  xlabel={extension $x$ (m)}, ylabel={force $F$ (N)},
  xmin=0, xmax=1.25, ymin=0, ymax=1.25,
  xtick={1}, xticklabels={$e$},
  ytick={1}, yticklabels={$ke$},
  tick label style={font=\small},
  label style={font=\small},
  clip=false
]
\addplot[draw=none, fill=popblueL, domain=0:1] {x} \closedcycle;
\addplot[popblue, very thick, domain=0:1.15] {x} node[pos=0.85, above left, popdark]{$F = kx$};
\draw[popdark, dashed] (axis cs:1,0) -- (axis cs:1,1);
\node[popdark] at (axis cs:0.55,0.28) {$W = \tfrac{1}{2}ke^2$};
\end{axis}
\end{tikzpicture}
\end{center}
\legende{The force--extension graph for a spring is the straight line $F = kx$. The work done stretching from $0$ to $e$ is the shaded triangular area: $W = \tfrac{1}{2}\times e \times ke = \tfrac{1}{2}ke^2$.}
\end{popfigure}

\begin{exemplebox}[Graphical check: area of a triangle]
Since $F = kx$ is a straight line through the origin, the area under it from $0$ to $e$ is a triangle of base $e$ and height $ke$:
\[
W = \frac{1}{2} \times \text{base} \times \text{height} = \frac{1}{2} \times e \times ke = \frac{1}{2}ke^2 ,
\]
confirming the integral. Note that $\tfrac{1}{2}ke$ is the \emph{average} force over the stretch, so $W = (\text{average force})\times(\text{extension})$ — a useful way to remember the factor $\tfrac{1}{2}$.
\end{exemplebox}

\courssub{Work Done Between Two Extensions}

Often the spring is \emph{already} stretched by $e_1$ and is then stretched further to $e_2$. The force still varies, so we integrate between the two extensions.

\begin{methode}[Work done in stretching from extension $e_1$ to extension $e_2$]
\begin{enumerate}
\item Write the applied force as a function of extension: $F(x) = kx$.
\item Integrate between the two extensions:
\[
W = \int_{e_1}^{e_2} kx\, dx = k\left[\frac{x^2}{2}\right]_{e_1}^{e_2}.
\]
\item Conclude:
\[
W = \frac{1}{2}k\left(e_2^2 - e_1^2\right).
\]
Equivalently, this is the \emph{area of a trapezium} under the line $F = kx$ between $e_1$ and $e_2$: parallel sides $ke_1$ and $ke_2$, width $(e_2 - e_1)$, giving $W = \tfrac{1}{2}(ke_1 + ke_2)(e_2 - e_1) = \tfrac{1}{2}k(e_2^2 - e_1^2)$.
\end{enumerate}
\end{methode}

\begin{exoresolu}[Stretching a spring further]
A light spring has stiffness $k = 200\ \mathrm{N\,m^{-1}}$. It is already stretched by $e_1 = 0.10\ \mathrm{m}$. Calculate the extra work needed to stretch it to an extension $e_2 = 0.30\ \mathrm{m}$.
\tcblower
\textbf{Solution.} Integrate the force $F(x) = 200x$ between the two extensions:
\begin{align*}
W &= \int_{0.10}^{0.30} 200x\, dx = 200\left[\frac{x^2}{2}\right]_{0.10}^{0.30} \\
&= 100\left(0.30^2 - 0.10^2\right) = 100(0.09 - 0.01) = 100 \times 0.08 .
\end{align*}
Hence $\boxed{W = 8\ \mathrm{J}}$. Notice that this is \emph{more} than the work needed for the first $0.10\ \mathrm{m}$ (which was only $\tfrac12(200)(0.10)^2 = 1\ \mathrm{J}$): the further the spring is stretched, the harder each additional centimetre becomes.
\end{exoresolu}

\begin{exoresolu}[A spring with a modulus of elasticity]
A light elastic string has natural length $l = 1.2\ \mathrm{m}$ and modulus of elasticity $\lambda = 60\ \mathrm{N}$. Find the work done in stretching it from its natural length to a total length of $1.5\ \mathrm{m}$.
\tcblower
\textbf{Solution.} First find the extension: $e = 1.5 - 1.2 = 0.3\ \mathrm{m}$. The stiffness is
\[
k = \frac{\lambda}{l} = \frac{60}{1.2} = 50\ \mathrm{N\,m^{-1}} .
\]
Therefore
\[
W = \frac{1}{2}ke^2 = \frac{1}{2}\times 50 \times (0.3)^2 = 25 \times 0.09 = 2.25\ \mathrm{J}.
\]
The work done is $\boxed{W = 2.25\ \mathrm{J}}$.
\end{exoresolu}

\begin{attention}[The most common trap: (final force) $\times$ (extension)]
Do \textbf{not} compute the work in stretching a spring as $(ke)\times e = ke^2$. The force is not constant — it grows from $0$ to $ke$ — so using the final force throughout \emph{overestimates the work by a factor of 2}. The correct result is $W = \tfrac{1}{2}ke^2$, i.e.\ the \emph{average} force $\tfrac{1}{2}ke$ times the extension. Similarly, between two extensions use $W = \tfrac{1}{2}k(e_2^2 - e_1^2)$, never $ke_2(e_2 - e_1)$.
\end{attention}

\courssub{Link with Elastic Potential Energy (Looking Ahead)}

\begin{savaistu}[Energy stored in a spring]
The work $\tfrac{1}{2}ke^2$ done in stretching a spring is not lost: it is stored inside the spring as \cle{elastic potential energy} (EPE),
\[
\text{EPE} = \frac{1}{2}ke^2 = \frac{\lambda x^2}{2l},
\]
and is released when the spring contracts. This is why a stretched catapult can launch a stone, and why the suspension springs of a \emph{bendskin} absorb the shocks of Douala's potholes and give the energy back smoothly. We shall exploit this fully in the sections on energy conservation: the work you do on a spring becomes stored energy, which can later become kinetic energy.
\end{savaistu}

\begin{exercice}[Practice]
\begin{enumerate}
\item A force $F(x) = 5x + 4$ newtons moves a body from $x = 0$ to $x = 2\ \mathrm{m}$ along its line of action. Find the work done.
\item A spring of stiffness $k = 150\ \mathrm{N\,m^{-1}}$ is stretched from its natural length by $0.20\ \mathrm{m}$. Find the work done.
\item For the spring of question 2, find the \emph{additional} work needed to increase the extension from $0.20\ \mathrm{m}$ to $0.40\ \mathrm{m}$. Why is this not simply twice the answer to question 2?
\item A light elastic string has natural length $0.8\ \mathrm{m}$ and modulus $\lambda = 40\ \mathrm{N}$. Calculate the work done in stretching it to a length of $1.0\ \mathrm{m}$.
\end{enumerate}
\end{exercice}

\begin{corrige}[Exercise 1]
\begin{enumerate}
\item $W = \displaystyle\int_0^2 (5x+4)\,dx = \left[\tfrac{5x^2}{2} + 4x\right]_0^2 = 10 + 8 = 18\ \mathrm{J}$.
\item $W = \tfrac{1}{2}ke^2 = \tfrac{1}{2}(150)(0.20)^2 = 75 \times 0.04 = 3\ \mathrm{J}$.
\item $W = \tfrac{1}{2}(150)\left(0.40^2 - 0.20^2\right) = 75(0.16 - 0.04) = 75 \times 0.12 = 9\ \mathrm{J}$. This is three times the answer to question 2 (not twice), because the force is larger on average over the second interval: work depends on $e^2$, not on $e$.
\item Extension $e = 1.0 - 0.8 = 0.2\ \mathrm{m}$; stiffness $k = \lambda/l = 40/0.8 = 50\ \mathrm{N\,m^{-1}}$. Hence $W = \tfrac{1}{2}(50)(0.2)^2 = 25 \times 0.04 = 1\ \mathrm{J}$.
\end{enumerate}
\end{corrige}

\begin{aretenir}[Summary of the section]
\begin{itemize}
\item Variable force along a line: $W = \displaystyle\int_a^b F(x)\,dx$ = area under the $F$--$x$ graph.
\item Hooke's law: $T = kx = \dfrac{\lambda}{l}x$, with $x$ the \emph{extension}.
\item Work to stretch a spring from natural length: $W = \tfrac{1}{2}ke^2$ (triangular area).
\item Work between extensions $e_1$ and $e_2$: $W = \tfrac{1}{2}k(e_2^2 - e_1^2)$ (trapezium area).
\item Never use (final force) $\times$ (extension): the factor $\tfrac{1}{2}$ comes from the force growing linearly from zero.
\item The work stored in a stretched spring is elastic potential energy, $\tfrac{1}{2}ke^2$, ready to be reused in energy problems.
\end{itemize}
\end{aretenir}

\coursec{Energy: Kinetic and Potential}

In the previous sections we learnt how to calculate the work done by a force, whether constant or variable. We established that work is the mechanism by which energy is transferred to or from a body. But what exactly \emph{is} energy? In mechanics, energy is the capacity to do work. A moving car can crash into a barrier and deform it — it has the capacity to do work because of its motion. A water tank on a rooftop in Bamenda can drive a turbine when the water falls — it has the capacity to do work because of its position. In this section we formalise these two fundamental forms of mechanical energy: \cle{kinetic energy} (energy of motion) and \cle{potential energy} (energy of position).

\courssub{The Concept of Energy}

\begin{definition}[Energy]
\textbf{Energy} is a scalar quantity that measures the capacity of a body or a system to do work. The SI unit of energy is the \textbf{joule (J)}, the same unit as work:
\[
1\ \mathrm{J} = 1\ \mathrm{N\,m} = 1\ \mathrm{kg\,m^2\,s^{-2}}.
\]
Energy, like work, is a scalar — it has magnitude but no direction.
\end{definition}

Energy exists in many forms: mechanical, chemical, electrical, nuclear, thermal, and so on. In mechanics we focus on \emph{mechanical energy}, which is the sum of kinetic energy and potential energy. The principle of conservation of energy, which we shall study in the next section, tells us that in an isolated system the total energy remains constant, merely transforming from one form to another.

\courssub{Kinetic Energy}

Consider a body of mass $m$ initially at rest on a smooth horizontal surface. A constant net force $F$ acts on it over a displacement $s$, accelerating it from rest to a speed $v$. The work done by the force is $W = Fs$. From Newton's second law, $F = ma$, and from the kinematic relation $v^2 = u^2 + 2as$ with $u = 0$, we obtain $v^2 = 2as$, that is $as = \frac{1}{2}v^2$. Hence
\[
W = Fs = m(as) = m\left(\frac{1}{2}v^2\right) = \frac{1}{2}mv^2.
\]
This work done on the body is now stored as its \emph{kinetic energy} — the energy it possesses by virtue of its motion. This derivation is examinable and should be reproduced fluently.

\begin{definition}[Kinetic Energy]
The \textbf{kinetic energy} (KE) of a body of mass $m$ moving with speed $v$ is defined as
\[
\boxed{\mathrm{KE} = \frac{1}{2}mv^2}.
\]
Kinetic energy is a scalar quantity measured in joules (J). It depends on the \emph{mass} and on the \emph{square of the speed}.
\end{definition}

\begin{propriete}[Properties of Kinetic Energy]
\begin{enumerate}
\item \textbf{Non-negativity:} $\mathrm{KE} \ge 0$ for all $v$; $\mathrm{KE} = 0$ only when $v = 0$ (the body is at rest).
\item \textbf{Depends on speed, not velocity:} Since $v^2 = \vec{v}\cdot\vec{v} = |\vec{v}|^2$, kinetic energy depends only on the magnitude of the velocity, never on its direction. A body moving north at $10\ \mathrm{m\,s^{-1}}$ has the same KE as a body of the same mass moving south at $10\ \mathrm{m\,s^{-1}}$.
\item \textbf{Change in kinetic energy:} If the speed changes from $u$ to $v$, the change in kinetic energy is
\[
\Delta\mathrm{KE} = \frac{1}{2}mv^2 - \frac{1}{2}mu^2 = \frac{1}{2}m(v^2 - u^2).
\]
This quantity can be positive (gain in KE), negative (loss in KE), or zero.
\item \textbf{Quadratic dependence on speed:} Doubling the speed \emph{quadruples} the kinetic energy; tripling the speed increases KE ninefold.
\end{enumerate}
\end{propriete}

\begin{attention}[Common Errors with Kinetic Energy]
\begin{itemize}
\item \textbf{Forgetting the factor $\frac{1}{2}$:} Writing $\mathrm{KE} = mv^2$ is a very frequent mistake. Always remember the half!
\item \textbf{Using velocity instead of speed:} Kinetic energy uses $v^2$ (speed squared), not $\vec{v}^{\,2}$. Never write $\frac{1}{2}m\vec{v}^{\,2}$.
\item \textbf{Confusing $\Delta\mathrm{KE}$ with $\mathrm{KE}$:} The work–energy theorem relates \emph{net work} to the \emph{change} in kinetic energy, not to the kinetic energy itself.
\item \textbf{Unit errors:} Mass must be in kg and speed in $\mathrm{m\,s^{-1}}$ to obtain energy in joules. If mass is given in tonnes or grams, or speed in $\mathrm{km\,h^{-1}}$, convert first (divide $\mathrm{km\,h^{-1}}$ by $3.6$).
\end{itemize}
\end{attention}

\begin{exemplebox}[Kinetic Energy Calculations]
Calculate the kinetic energy of:
\begin{enumerate}
\item A car of mass $1200\ \mathrm{kg}$ travelling at $25\ \mathrm{m\,s^{-1}}$ ($90\ \mathrm{km\,h^{-1}}$).
\item A bullet of mass $8\ \mathrm{g}$ fired at $400\ \mathrm{m\,s^{-1}}$.
\item A lorry of mass $15$ tonnes whose speed increases from $10\ \mathrm{m\,s^{-1}}$ to $20\ \mathrm{m\,s^{-1}}$. Find the gain in kinetic energy.
\end{enumerate}

\tcblower
\textbf{Solution.}
\begin{enumerate}
\item $\mathrm{KE} = \frac{1}{2}(1200)(25^2) = 600 \times 625 = 375\,000\ \mathrm{J} = 375\ \mathrm{kJ}$.
\item Mass $m = 8\ \mathrm{g} = 0.008\ \mathrm{kg}$. $\mathrm{KE} = \frac{1}{2}(0.008)(400^2) = 0.004 \times 160\,000 = 640\ \mathrm{J}$.
\item $m = 15$ tonnes $= 15\,000\ \mathrm{kg}$.
\[
\Delta\mathrm{KE} = \frac{1}{2}(15\,000)(20^2 - 10^2) = 7\,500 \times (400 - 100) = 7\,500 \times 300 = 2\,250\,000\ \mathrm{J} = 2.25\ \mathrm{MJ}.
\]
Notice that the gain in KE when the speed doubles from $10$ to $20\ \mathrm{m\,s^{-1}}$ is three times the initial KE at $10\ \mathrm{m\,s^{-1}}$ (which is $\frac{1}{2}(15\,000)(10^2) = 750\ \mathrm{kJ}$).
\end{enumerate}
\end{exemplebox}

\courssub{Gravitational Potential Energy}

When you lift a book from the floor to a shelf, you do work against the gravitational force. That work is stored as \emph{gravitational potential energy} — energy due to the body's position in a gravitational field. Unlike kinetic energy, potential energy is defined only up to an additive constant: we must choose a \emph{reference level} (zero level) where $\mathrm{PE} = 0$.

\begin{definition}[Gravitational Potential Energy]
The \textbf{gravitational potential energy} (PE) of a body of mass $m$ at a vertical height $h$ above a chosen reference level is
\[
\boxed{\mathrm{PE} = mgh},
\]
where $g = 9.81\ \mathrm{m\,s^{-2}}$ (or $9.8\ \mathrm{m\,s^{-2}}$ as specified in the examination). PE is a scalar measured in joules (J). The height $h$ is the \emph{vertical} distance above the reference level.
\end{definition}

\begin{popfigure}
\begin{tikzpicture}[scale=1.1]
% Ground line
\draw[thick, popdark] (-3,0) -- (5,0) node[right] {Ground};
% Reference level (dashed)
\draw[dashed, thick, popblue] (-3,1) -- (5,1) node[right, popblue] {Reference level ($h=0$)};
% Height markers
\draw[<->, poporange, thick] (3.5,1) -- (3.5,3) node[midway, right, poporange] {$h_1 = 2\,\mathrm{m}$};
\draw[<->, poporange, thick] (4.2,1) -- (4.2,4.5) node[midway, right, poporange] {$h_2 = 3.5\,\mathrm{m}$};
% Ball at height h1
\draw[fill=popblue, draw=popdark] (2,3) circle (0.3);
\node[above] at (2,3.3) {$m$};
\node[below] at (2,2.6) {$\mathrm{PE}_1 = mgh_1$};
% Ball at height h2
\draw[fill=poporange, draw=popdark] (0.5,4.5) circle (0.3);
\node[above] at (0.5,4.8) {$m$};
\node[below] at (0.5,4.1) {$\mathrm{PE}_2 = mgh_2$};
% Vertical dashed lines from reference
\draw[dashed, thin, gray] (2,1) -- (2,3);
\draw[dashed, thin, gray] (0.5,1) -- (0.5,4.5);
% g vector
\draw[->, thick, poppurple] (4.7,2.5) -- (4.7,1.5) node[midway, right, poppurple] {$\vec{g}$};
% Labels
\node[align=center] at (-2,4) {Gravitational potential energy\\depends only on vertical height\\above the reference level.};
\end{tikzpicture}
\legende{Gravitational potential energy at two different heights above a chosen reference level.}
\end{popfigure}

\begin{propriete}[Key Points on Gravitational Potential Energy]
\begin{enumerate}
\item \textbf{Reference level is arbitrary:} You may choose the ground, the floor, the lowest point of the motion, or any convenient horizontal plane as $h = 0$. The \emph{value} of PE changes with the choice, but the \emph{difference} in PE between two points does not.
\item \textbf{Only changes in PE are physically significant:} The work done by gravity when a body moves from height $h_1$ to height $h_2$ is $W = mg(h_1 - h_2) = -\Delta\mathrm{PE}$. This is independent of the reference level.
\item \textbf{Gain and loss:} When a body is raised, $\Delta\mathrm{PE} > 0$ (gain). When it is lowered, $\Delta\mathrm{PE} < 0$ (loss).
\item \textbf{Path independence:} The change in gravitational PE depends only on the initial and final vertical heights, not on the path taken (straight up, inclined plane, or curved track). This is why gravity is called a \emph{conservative} force.
\end{enumerate}
\end{propriete}

\begin{methode}[Choosing and Using a Reference Level for PE]
\begin{enumerate}
\item \textbf{Identify} the initial and final positions of the body in the problem.
\item \textbf{Choose} a convenient horizontal reference level (often the lowest point reached, or the ground). \textbf{State it clearly:} ``Take the reference level as the ground / the bottom of the incline / the lowest point.''
\item \textbf{Measure} the vertical heights $h_i$ and $h_f$ of the initial and final positions above this reference.
\item \textbf{Compute} $\mathrm{PE}_i = mgh_i$ and $\mathrm{PE}_f = mgh_f$.
\item \textbf{Find the change:} $\Delta\mathrm{PE} = \mathrm{PE}_f - \mathrm{PE}_i = mg(h_f - h_i)$.
\end{enumerate}
\end{methode}

\begin{exemplebox}[Potential Energy on an Inclined Plane]
A block of mass $5\ \mathrm{kg}$ slides down a smooth plane inclined at $30^\circ$ to the horizontal. The block travels $4\ \mathrm{m}$ along the plane from rest. Take the reference level as the bottom of the incline (where the block would be after sliding $4\ \mathrm{m}$). Find:
\begin{enumerate}
\item The initial height of the block above the reference level.
\item The initial and final gravitational potential energies.
\item The loss in potential energy.
\end{enumerate}

\tcblower
\textbf{Solution.}
\begin{enumerate}
\item The vertical height dropped is $h = 4\sin 30^\circ = 4 \times 0.5 = 2\ \mathrm{m}$. Since the reference is the bottom, the initial height is $h_i = 2\ \mathrm{m}$ and the final height is $h_f = 0$.
\item $\mathrm{PE}_i = mgh_i = 5 \times 9.81 \times 2 = 98.1\ \mathrm{J}$. $\mathrm{PE}_f = 0$.
\item Loss in PE $= \mathrm{PE}_i - \mathrm{PE}_f = 98.1\ \mathrm{J}$. (Equivalently $\Delta\mathrm{PE} = -98.1\ \mathrm{J}$.)
\end{enumerate}
\textbf{Note:} If we had chosen the starting point as reference ($h_i = 0$, $h_f = -2\ \mathrm{m}$), then $\mathrm{PE}_i = 0$, $\mathrm{PE}_f = -98.1\ \mathrm{J}$, and $\Delta\mathrm{PE} = -98.1\ \mathrm{J}$ — the same change, confirming that only differences matter.
\end{exemplebox}

\courssub{Elastic Potential Energy (Brief Overview)}

While the syllabus focuses on gravitational potential energy, it is useful to know that energy can also be stored in a deformed elastic body, such as a stretched or compressed spring.

\begin{propriete}[Elastic Potential Energy]
For a light spring of stiffness $k$ (spring constant, in $\mathrm{N\,m^{-1}}$) stretched or compressed by an extension $x$ from its natural length, the \textbf{elastic potential energy} (EPE) stored is
\[
\boxed{\mathrm{EPE} = \frac{1}{2}kx^2}.
\]
This follows from the work done against the variable tension $T = kx$ (Hooke's law): $W = \displaystyle\int_0^x kx\,dx = \frac{1}{2}kx^2$. Like gravitational PE, EPE is always non-negative, and it depends on the square of the deformation.
\end{propriete}

\begin{attention}[Elastic PE — Extension Beyond the Core Syllabus]
Elastic potential energy is not explicitly required in the GCE Advanced Level Pure Mathematics with Mechanics syllabus for this chapter, but it appears in Further Mechanics and in many physics contexts. Notice the pattern $\frac{1}{2}kx^2$, exactly analogous to $\frac{1}{2}mv^2$: energy stored in a quadratic degree of freedom always carries the factor $\frac{1}{2}$.
\end{attention}

\courssub{Computational Drills and Applications}

Let us now practise quick calculations that illustrate the quadratic growth of kinetic energy with speed and the linear growth of potential energy with height.

\begin{popfigure}
\begin{tikzpicture}
\begin{axis}[
    ybar,
    bar width=1.2cm,
    width=12cm,
    height=7cm,
    enlarge x limits=0.3,
    ylabel={Kinetic Energy (kJ)},
    xlabel={Speed (m/s)},
    symbolic x coords={10, 20, 30},
    xtick=data,
    nodes near coords,
    nodes near coords align={vertical},
    ymin=0,
    ymax=600,
    ytick={0,100,200,300,400,500},
    title={Kinetic Energy of a 1000 kg Car at Different Speeds},
    title style={font=\bfseries},
]
\addplot[ybar, fill=popblue, draw=popdark] coordinates {(10,50) (20,200) (30,450)};
\end{axis}
\end{tikzpicture}
\legende{Quadratic growth of kinetic energy with speed for a car of mass $1000\ \mathrm{kg}$. KE at $10\ \mathrm{m\,s^{-1}}$ is $50\ \mathrm{kJ}$; at $20\ \mathrm{m\,s^{-1}}$ it is $200\ \mathrm{kJ}$ ($4\times$); at $30\ \mathrm{m\,s^{-1}}$ it is $450\ \mathrm{kJ}$ ($9\times$).}
\end{popfigure}

\begin{exemplebox}[Quick Drills: KE and PE]
\begin{enumerate}
\item A motorbike of mass $200\ \mathrm{kg}$ has a kinetic energy of $50\ \mathrm{kJ}$. What is its speed?
\item Water of mass $5000\ \mathrm{kg}$ is stored in a tank $20\ \mathrm{m}$ above ground level. Taking the ground as reference, find the gravitational potential energy of the water.
\item A stone of mass $0.5\ \mathrm{kg}$ is thrown vertically upwards with an initial speed of $20\ \mathrm{m\,s^{-1}}$. Find its kinetic energy at the start and its maximum gain in potential energy (ignore air resistance).
\end{enumerate}

\tcblower
\textbf{Solution.}
\begin{enumerate}
\item $\mathrm{KE} = \frac{1}{2}mv^2 \Rightarrow 50\,000 = \frac{1}{2}(200)v^2 = 100v^2 \Rightarrow v^2 = 500 \Rightarrow v = \sqrt{500} \approx 22.4\ \mathrm{m\,s^{-1}}$.
\item $\mathrm{PE} = mgh = 5000 \times 9.81 \times 20 = 981\,000\ \mathrm{J} = 981\ \mathrm{kJ}$.
\item Initial KE $= \frac{1}{2}(0.5)(20^2) = 0.25 \times 400 = 100\ \mathrm{J}$. At maximum height the speed is $0$, so the final KE is $0$. By conservation of energy (no air resistance), loss in KE $=$ gain in PE $= 100\ \mathrm{J}$. Maximum height: $h = \dfrac{\mathrm{PE}}{mg} = \dfrac{100}{0.5 \times 9.81} \approx 20.4\ \mathrm{m}$.
\end{enumerate}
\end{exemplebox}

\begin{exoresolu}[Energy Changes on a Roller Coaster]
A roller coaster car of mass $800\ \mathrm{kg}$ (including passengers) starts from rest at point A, which is $40\ \mathrm{m}$ above the ground. It descends to point B at ground level, then rises to point C at $25\ \mathrm{m}$ above ground. Assume the track is smooth (no friction). Take the ground as the reference level for PE ($g = 9.81\ \mathrm{m\,s^{-2}}$). Find:
\begin{enumerate}
\item The total mechanical energy at A.
\item The kinetic energy and speed at B.
\item The kinetic energy and speed at C.
\end{enumerate}

\tcblower
\textbf{Solution.}
Since the track is smooth, mechanical energy is conserved: $\mathrm{KE} + \mathrm{PE} = \text{constant}$.

\textbf{At A:} $v_A = 0$, $h_A = 40\ \mathrm{m}$.
$\mathrm{KE}_A = 0$.
$\mathrm{PE}_A = mgh_A = 800 \times 9.81 \times 40 = 313\,920\ \mathrm{J}$.
Total mechanical energy: $E = \mathrm{KE}_A + \mathrm{PE}_A = 313\,920\ \mathrm{J}$.

\textbf{At B:} $h_B = 0$, so $\mathrm{PE}_B = 0$.
$\mathrm{KE}_B = E - \mathrm{PE}_B = 313\,920\ \mathrm{J}$.
$\frac{1}{2}mv_B^2 = 313\,920 \Rightarrow v_B^2 = \dfrac{2 \times 313\,920}{800} = 784.8 \Rightarrow v_B \approx 28.0\ \mathrm{m\,s^{-1}}$.

\textbf{At C:} $h_C = 25\ \mathrm{m}$.
$\mathrm{PE}_C = 800 \times 9.81 \times 25 = 196\,200\ \mathrm{J}$.
$\mathrm{KE}_C = E - \mathrm{PE}_C = 313\,920 - 196\,200 = 117\,720\ \mathrm{J}$.
$\frac{1}{2}mv_C^2 = 117\,720 \Rightarrow v_C^2 = \dfrac{2 \times 117\,720}{800} = 294.3 \Rightarrow v_C \approx 17.2\ \mathrm{m\,s^{-1}}$.

\textbf{Check:} the mass cancels in $\frac{1}{2}mv^2 + mgh = \text{constant}$, so the speeds depend only on the heights, not on the mass of the car — a useful quick check in the examination.
\end{exoresolu}

\courssub{Summary of Key Points}

\begin{aretenir}[Energy: Kinetic and Potential — Key Points]
\begin{itemize}
\item \textbf{Energy} is the capacity to do work; scalar, unit joule (J).
\item \textbf{Kinetic Energy:} $\mathrm{KE} = \frac{1}{2}mv^2$. Always $\ge 0$. Depends on speed squared, not direction. $\Delta\mathrm{KE} = \frac{1}{2}m(v^2 - u^2)$.
\item \textbf{Gravitational Potential Energy:} $\mathrm{PE} = mgh$ relative to a chosen reference level. Only \emph{changes} in PE are physically meaningful: $\Delta\mathrm{PE} = mg\,\Delta h$.
\item \textbf{Reference level:} Must be stated clearly. Convenient choices: ground, lowest point, initial position.
\item \textbf{Elastic Potential Energy (extension):} $\mathrm{EPE} = \frac{1}{2}kx^2$ for a spring of stiffness $k$ extended by $x$.
\item \textbf{Quadratic vs Linear:} KE grows with $v^2$ (doubling speed $\Rightarrow$ KE $\times 4$). PE grows linearly with $h$.
\item \textbf{Common pitfalls:} Forgetting $\frac{1}{2}$ in KE; using $v$ instead of $v^2$; not converting mass to kg or speed to $\mathrm{m\,s^{-1}}$; not stating the reference level for PE.
\end{itemize}
\end{aretenir}

\begin{exercice}[Practice Exercise]
\begin{enumerate}
\item A truck of mass $12$ tonnes travels at $72\ \mathrm{km\,h^{-1}}$. Calculate its kinetic energy in MJ.
\item A girl of mass $45\ \mathrm{kg}$ climbs a vertical rope through $6\ \mathrm{m}$. Taking the starting point as reference, find her gain in gravitational potential energy.
\item A spring of stiffness $200\ \mathrm{N\,m^{-1}}$ is compressed by $0.15\ \mathrm{m}$. Find the elastic potential energy stored.
\item A ball of mass $0.2\ \mathrm{kg}$ is dropped from a height of $10\ \mathrm{m}$. Find its kinetic energy just before it hits the ground (ignore air resistance, $g = 9.81\ \mathrm{m\,s^{-2}}$).
\item A car of mass $1000\ \mathrm{kg}$ accelerates from $15\ \mathrm{m\,s^{-1}}$ to $30\ \mathrm{m\,s^{-1}}$. Find the increase in its kinetic energy. By what factor does the KE increase?
\end{enumerate}
\end{exercice}

\begin{corrige}[Exercise]
\begin{enumerate}
\item $m = 12\,000\ \mathrm{kg}$, $v = 72\ \mathrm{km\,h^{-1}} = 20\ \mathrm{m\,s^{-1}}$. $\mathrm{KE} = \frac{1}{2}(12\,000)(20^2) = 2\,400\,000\ \mathrm{J} = 2.4\ \mathrm{MJ}$.
\item Gain in PE $= mg\,\Delta h = 45 \times 9.81 \times 6 = 2648.7\ \mathrm{J} \approx 2.65\ \mathrm{kJ}$.
\item $\mathrm{EPE} = \frac{1}{2}kx^2 = \frac{1}{2}(200)(0.15^2) = 100 \times 0.0225 = 2.25\ \mathrm{J}$.
\item Loss in PE $=$ gain in KE $= mgh = 0.2 \times 9.81 \times 10 = 19.62\ \mathrm{J}$.
\item $\Delta\mathrm{KE} = \frac{1}{2}(1000)(30^2 - 15^2) = 500 \times (900 - 225) = 500 \times 675 = 337\,500\ \mathrm{J} = 337.5\ \mathrm{kJ}$. Initial KE $= \frac{1}{2}(1000)(15^2) = 112.5\ \mathrm{kJ}$; final KE $= 450\ \mathrm{kJ}$. Factor $= \dfrac{450}{112.5} = 4$ (speed doubled $\Rightarrow$ KE $\times 4$).
\end{enumerate}
\end{corrige}

\coursec{The Work--Energy Theorem and Conservation of Energy}

In the previous section we met the two great forms of mechanical energy: \cle{kinetic energy} $\tfrac12 mv^2$, the energy a body possesses because it moves, and \cle{gravitational potential energy} $mgh$, the energy it possesses because of its position in the Earth's gravitational field. We also saw that \cle{work} is the mechanism by which energy is transferred: when a force does work on a body, energy flows into (or out of) that body. In this section we make that idea precise. The result is the \cle{work--energy theorem}, one of the most powerful labour-saving devices in all of mechanics, and its celebrated consequence, the \cle{principle of conservation of mechanical energy}. Together they allow us to find speeds directly --- without resolving forces, without finding accelerations, and without touching the equations of motion.

\courssub{The Work--Energy Theorem}

\textbf{Intuition.}\par
Push a trolley along a level floor. The work you do does not disappear --- it reappears as motion: the trolley speeds up. If a rough patch of floor does negative work on the trolley, the trolley slows down: kinetic energy is drained away. This suggests a precise accounting rule:

\[
\text{(total work done on the body)} \;=\; \text{(change in its kinetic energy)}.
\]

\begin{propriete}[Work--Energy Theorem]
The total (net) work done by \emph{all} the forces acting on a body during a displacement equals the change in the body's kinetic energy:
\[
W_{\text{net}} \;=\; \Delta \mathrm{KE} \;=\; \tfrac12 mv^2 - \tfrac12 mu^2,
\]
where $u$ and $v$ are the initial and final speeds. This holds whether the forces are constant or variable, and whether the path is straight or curved.
\end{propriete}

\textbf{Derivation (examinable for constant force in a straight line; general case by integration).}\par
For a body of mass $m$ moving in a straight line under a \emph{constant} resultant force $F$, Newton's second law gives $F = ma$, where $a$ is the constant acceleration. Over a displacement $s$, the constant-acceleration formula
\[
v^2 = u^2 + 2as \qquad\Longrightarrow\qquad as = \tfrac12\!\left(v^2 - u^2\right).
\]
The work done by the resultant force is $W = Fs$. Substituting $F = ma$:
\[
W = mas = m\cdot \tfrac12\!\left(v^2 - u^2\right) = \tfrac12 mv^2 - \tfrac12 mu^2. \qquad \blacksquare
\]
For a variable force $F(x)$ acting along a straight line, the work is $W = \int F(x)\,dx$. Using $F = m\,\frac{dv}{dt} = m v \frac{dv}{dx}$, we obtain
\[
W = \int m v \frac{dv}{dx}\,dx = \int_{u}^{v} m v\,dv = \tfrac12 mv^2 - \tfrac12 mu^2.
\]
For a curved path, the same argument applies to the tangential component of the resultant force. Hence the theorem is completely general.

\begin{attention}[Signs matter]
Work done \emph{by} a force that helps the motion is positive and \emph{increases} the kinetic energy; work done \emph{against} the motion (friction, resistance, braking) is negative and \emph{decreases} it. Always compute the \emph{net} work: add the work of every force with its correct sign.
\end{attention}

\courssub{Horizontal Motion with Driving and Resisting Forces}

A very common GCE situation: a vehicle or body moves along a \emph{horizontal} road under a driving force $F$ and a resistance $R$. Gravity and the normal reaction do no work (they are perpendicular to the motion), so the work--energy theorem over a distance $s$ gives
\[
(F - R)\,s \;=\; \tfrac12 mv^2 - \tfrac12 mu^2.
\]

\begin{exoresolu}[Accelerating car on a level road]
A car of mass $900\ \mathrm{kg}$ accelerates from $5\ \mathrm{m\,s^{-1}}$ to $15\ \mathrm{m\,s^{-1}}$ over a distance of $200\ \mathrm{m}$ on a horizontal road, against a constant resistance of $400\ \mathrm{N}$. Find the driving force exerted by the engine.
\tcblower
\textbf{Solution.} Apply the work--energy theorem with $u = 5$, $v = 15$, $s = 200$, $m = 900$:
\[
(F - 400)(200) = \tfrac12(900)\left(15^2 - 5^2\right) = 450 \times 200 = 90\,000\ \mathrm{J}.
\]
Hence
\[
F - 400 = \frac{90\,000}{200} = 450 \qquad\Longrightarrow\qquad \boxed{F = 850\ \mathrm{N}}.
\]
Notice that we never needed the acceleration or the time taken.
\end{exoresolu}

\courssub{Conservation of Mechanical Energy}

Now suppose a body moves under gravity \emph{alone} (no friction, no resistance, no driving force). Over a vertical drop $h$, gravity does work $mgh$, so by the work--energy theorem
\[
mgh = \tfrac12 mv^2 - \tfrac12 mu^2
\qquad\Longrightarrow\qquad
\tfrac12 mu^2 + mgh_{\text{initial}} = \tfrac12 mv^2 + mgh_{\text{final}}.
\]
In other words, the sum $\mathrm{KE} + \mathrm{PE}$ is unchanged.

\begin{propriete}[Conservation of Mechanical Energy]
If gravity is the only force that does work on a body (i.e. no non-conservative forces such as friction or air resistance do work), then its total mechanical energy is constant:
\[
\mathrm{KE} + \mathrm{PE} = \tfrac12 mv^2 + mgh = \text{constant}.
\]
Equivalently, any loss of potential energy appears as an equal gain of kinetic energy, and conversely.
\end{propriete}

\begin{attention}[When conservation fails]
Conservation of mechanical energy requires that \emph{no non-conservative force does work}. On a \textbf{rough} surface, friction does negative work and mechanical energy is \emph{not} conserved --- it is converted into heat. Do not write $\mathrm{KE}+\mathrm{PE}=\text{constant}$ on a rough slope; use the energy balance with a friction term instead (below).
\end{attention}

\begin{popfigure}
\begin{center}
\begin{tikzpicture}[scale=1.0]
  % pendulum pivot and string positions
  \coordinate (O) at (0,0);
  \fill[popdark] (O) circle (1.5pt);
  % arc of swing
  \draw[popline, dashed] (210:3.2) arc (210:330:3.2);
  % left (highest) position
  \draw[popblue, thick] (O) -- (210:3.2);
  \fill[popblue] (210:3.2) circle (0.22);
  \node[popblue, left] at (210:3.35) {\small $A$};
  % right (highest) position
  \draw[popblue!50, thick] (O) -- (330:3.2);
  \fill[popblue!50] (330:3.2) circle (0.22);
  % lowest position
  \draw[poporange, thick] (O) -- (270:3.2);
  \fill[poporange] (270:3.2) circle (0.22);
  \node[poporange, below] at (270:3.35) {\small $B$};
  % height annotation
  \draw[<->, popgreen, thick] ($(210:3.2)+(0.55,0)$) -- ($(270:3.2)+(0.55,0)$)
     node[midway, right, popgreen] {\small $h$};
  % labels
  \node[popblue, align=center, text width=2.5cm] at (-3.6,-1.1) {\small $\mathrm{PE}=mgh$\\ \small $\mathrm{KE}=0$};
  \node[poporange, align=center, text width=2.5cm] at (0.9,-3.6) {\small $\mathrm{KE}=\tfrac12 mv^2$\\ \small $\mathrm{PE}=0$};
  % velocity arrow at bottom
  \draw[->, poporange, very thick] ($(270:3.2)+(0.25,0)$) -- ++(1.1,0) node[right] {\small $v$};
\end{tikzpicture}
\end{center}
\legende{A pendulum bob released from rest at $A$. At the top all the energy is potential; at the bottom $B$ it is entirely kinetic: $mgh=\tfrac12 mv^2$.}
\end{popfigure}

\begin{exemplebox}[The swinging pendulum]
A pendulum bob of mass $0.2\ \mathrm{kg}$ is released from rest with the string taut at a point $0.45\ \mathrm{m}$ vertically above its lowest position. Taking $g = 9.8\ \mathrm{m\,s^{-2}}$, conservation of energy between the top and the bottom gives
\[
mgh = \tfrac12 mv^2 \;\Longrightarrow\; v = \sqrt{2gh} = \sqrt{2 \times 9.8 \times 0.45} = \sqrt{8.82} \approx 3.0\ \mathrm{m\,s^{-1}}.
\]
The mass cancels --- the speed at the bottom is independent of the bob's mass.
\end{exemplebox}

\courssub{Energy Balance with Friction: the General Energy Equation}

When friction (or another resistance) does work, mechanical energy is no longer conserved, but \emph{energy} itself is: the ``missing'' mechanical energy is exactly the work done against the resistance. The work--energy theorem, applied to all forces, gives the master equation of this chapter:

\begin{aretenir}[General energy balance]
For a body moving under the action of external forces (other than gravity) and a resistance $R$ over a path of length $s$:
\[
W_{\text{ext}} \;=\; \Delta \mathrm{KE} \;+\; \Delta \mathrm{PE} \;+\; W_{\text{resistance}},
\]
where $W_{\text{ext}}$ is the work done by all external forces except gravity, $\Delta \mathrm{KE} = \tfrac12 mv^2-\tfrac12 mu^2$, $\Delta \mathrm{PE} = mg\,\Delta h$, and $W_{\text{resistance}} = Rs$ for a constant resistance $R$ opposing the motion.
In words: \emph{initial energy + work input = final energy + energy lost to friction.}
\end{aretenir}

\begin{methode}[Solving ``find the speed'' problems by energy]
\begin{enumerate}
  \item Identify the initial and final positions and speeds; choose a reference level for PE.
  \item List every force that \emph{does work}: gravity (via $\Delta \mathrm{PE}$), driving forces, resistances.
  \item Write the energy balance: initial $\mathrm{KE}+\mathrm{PE}$ + work input $=$ final $\mathrm{KE}+\mathrm{PE}$ + work against friction.
  \item Solve for the unknown speed, force or distance. No acceleration, no time, no resolving needed.
\end{enumerate}
\end{methode}

\begin{popfigure}
\begin{center}
\begin{tikzpicture}[scale=1.05]
  % inclined plane
  \coordinate (A) at (0,3);   % top of slope
  \coordinate (B) at (5,0);   % bottom
  \coordinate (C) at (0,0);
  \draw[popdark, thick] (A) -- (B) -- (C) -- cycle;
  \draw[popline] (C) -- (B);
  % block at top
  \draw[fill=popblueL] (0.35,3.17) rectangle (1.05,3.75);
  \node[popblue, align=center, text width=2cm] at (0.7,4.0) {\small rest: $\mathrm{PE}=mgh$};
  % block at bottom
  \draw[fill=popgreenL] (4.15,0.18) rectangle (4.85,0.76);
  \node[popgreen, align=center, text width=2.5cm] at (6.6,0.5) {\small $\mathrm{KE}=\tfrac12 mv^2$};
  % height
  \draw[<->, poppurple, thick] (-0.35,0) -- (-0.35,3) node[midway, left] {\small $h$};
  % friction arrow on slope
  \draw[->, poporange, very thick] (2.6,1.42) -- (1.7,2.0) node[above left, poporange] {\small friction $R$};
  % energy flow annotation
  \node[align=center, popdark] at (2.5,-0.75) {\small $mgh \;=\; \tfrac12 mv^2 \;+\; R\,s$};
\end{tikzpicture}
\end{center}
\legende{Energy flow for a block sliding from rest down a rough slope: the initial potential energy splits into final kinetic energy plus the work done against friction.}
\end{popfigure}

\begin{exoresolu}[Sliding down a rough slope]
A block of mass $4\ \mathrm{kg}$ is released from rest at the top of a rough slope of length $10\ \mathrm{m}$ and vertical height $6\ \mathrm{m}$. The frictional resistance is constant at $8\ \mathrm{N}$. Taking $g = 9.8\ \mathrm{m\,s^{-2}}$, find the speed of the block at the bottom.
\tcblower
\textbf{Solution.} Use the energy balance (no driving force):
\[
\text{initial PE} = \text{final KE} + \text{work against friction}.
\]
\[
mgh = \tfrac12 mv^2 + Rs
\;\Longrightarrow\;
4 \times 9.8 \times 6 = \tfrac12 (4) v^2 + 8 \times 10.
\]
\[
235.2 = 2v^2 + 80 \;\Longrightarrow\; 2v^2 = 155.2 \;\Longrightarrow\; v^2 = 77.6.
\]
\[
\boxed{v \approx 8.8\ \mathrm{m\,s^{-1}}}
\]
Check: on a \emph{smooth} slope the speed would be $\sqrt{2gh} = \sqrt{117.6} \approx 10.8\ \mathrm{m\,s^{-1}}$; friction has reduced it, as expected.
\end{exoresolu}

\begin{exemplebox}[Body projected vertically upwards]
A stone is thrown vertically upwards with speed $14\ \mathrm{m\,s^{-1}}$. Ignoring air resistance, how high does it rise? Conservation of energy between launch and the highest point (where $v=0$):
\[
\tfrac12 mu^2 = mgh \;\Longrightarrow\; h = \frac{u^2}{2g} = \frac{196}{19.6} = 10\ \mathrm{m}.
\]
Compare the labour with the kinematic route ($v^2=u^2-2gs$): the energy method gives the same answer in one line.
\end{exemplebox}

\begin{savaistu}[Why engineers love energy]
Energy methods are the everyday tool of engineers designing hydroelectric schemes such as those on the Sanaga river: the potential energy of water held at height is converted into kinetic energy and then into electrical energy, and the whole design calculation is an energy balance --- exactly the equation $mgh = \tfrac12 mv^2 + \text{losses}$ that you have just used.
\end{savaistu}

\begin{exercice}[]
A cyclist of total mass $80\ \mathrm{kg}$ freewheels (no pedalling) from rest down a hill of vertical drop $20\ \mathrm{m}$ along a road of length $250\ \mathrm{m}$, against a constant resistance of $30\ \mathrm{N}$. Taking $g = 9.8\ \mathrm{m\,s^{-2}}$, find the speed at the bottom of the hill.
\end{exercice}

\begin{corrige}[Exercise 1]
Energy balance: initial PE $=$ final KE $+$ work against resistance:
\[
80 \times 9.8 \times 20 = \tfrac12(80)v^2 + 30 \times 250
\;\Longrightarrow\;
15\,680 = 40v^2 + 7\,500.
\]
\[
40v^2 = 8\,180 \;\Longrightarrow\; v^2 = 204.5 \;\Longrightarrow\; v \approx 14.3\ \mathrm{m\,s^{-1}}.
\]
\end{corrige}

\begin{exercice}[]
A ball of mass $0.5\ \mathrm{kg}$ is projected vertically upwards at $12\ \mathrm{m\,s^{-1}}$. It reaches a greatest height of $6.5\ \mathrm{m}$. Find the work done against air resistance during the ascent, and deduce the mean resisting force. Take $g = 9.8\ \mathrm{m\,s^{-2}}$.
\end{exercice}

\begin{corrige}[Exercise 2]
Initial KE $= \tfrac12(0.5)(144) = 36\ \mathrm{J}$. Final PE $= 0.5 \times 9.8 \times 6.5 = 31.85\ \mathrm{J}$.
Work against resistance $= 36 - 31.85 = 4.15\ \mathrm{J}$.
Mean force $= \dfrac{4.15}{6.5} \approx 0.64\ \mathrm{N}$.
\end{corrige}

\begin{aretenir}[Exam tip (GCE)]
Whenever a question asks you to \emph{``find the speed''} of a body that has moved through a height or along a slope --- especially with friction mentioned --- reach for the energy balance first. It avoids resolving forces, avoids finding the acceleration, and avoids the equations of motion entirely. Reserve Newton's second law for questions that ask for accelerations, times, or tensions.
\end{aretenir}

\coursec{Power: Definition, Instantaneous Power and Power--Velocity}

In the previous section we established the \cle{work--energy theorem}: the work done by the resultant force on a body equals its change in kinetic energy. That theorem tells us \emph{how much} work is done, but it says nothing about \emph{how fast} the work is done. Yet in everyday life this distinction is crucial. Two cranes on a construction site in Douala may lift the same load of cement through the same height --- they do the same work --- but the one that finishes in ten seconds is clearly more ``powerful'' than the one that needs a full minute. Similarly, two cars may both be capable of reaching \SI{100}{km/h}, but the one that gets there in eight seconds has a more powerful engine. This section makes that idea precise.

\courssub{Definition of Power}

\begin{experience}[Comparing two workers]
Manga and Etienne each carry a sack of rice of mass \SI{50}{kg} up a ladder of height \SI{4}{m}. Each does work
\[
W = mgh = 50 \times 9.81 \times 4 = 1962\ \mathrm{J}
\]
against gravity. Manga takes \SI{8}{s}; Etienne takes \SI{16}{s}. Same work, different \emph{rates}. Manga works at $1962/8 \approx 245$ joules per second, Etienne at only $1962/16 \approx 123$ joules per second. The quantity ``joules per second'' is what physicists call \cle{power}.
\end{experience}

\begin{definition}[Power]
\cle{Power} is the rate at which work is done (equivalently, the rate at which energy is transferred or converted).
\begin{itemize}
\item If a constant amount of work $W$ is done in time $t$, the \textbf{average power} is
\[
P_{\text{av}} = \frac{W}{t}.
\]
\item The \textbf{instantaneous power} at time $t$ is
\[
P = \frac{dW}{dt}.
\]
\end{itemize}
The SI unit of power is the \cle{watt} (W): $1\ \mathrm{W} = 1\ \mathrm{J\,s^{-1}}$. Larger units are common: $1\ \mathrm{kW} = 1000\ \mathrm{W}$ and $1\ \mathrm{MW} = 10^{6}\ \mathrm{W}$.
\end{definition}

\begin{attention}[Average power is not instantaneous power]
Saying ``the engine develops \SI{60}{kW}'' refers to the power \emph{at each instant}. If a machine does \SI{6000}{J} of work in \SI{2}{s}, its average power is \SI{3000}{W}, but its instantaneous power may have varied throughout those two seconds. In GCE questions about vehicles, $P$ almost always means the (constant) instantaneous power output of the engine.
\end{attention}

\courssub{The Power--Velocity Relation $P = Fv$}

Consider a constant force $F$ acting on a body and moving its point of application with velocity $v$ in the same direction as the force. In a small time interval $\delta t$, the body moves a distance $\delta s = v\,\delta t$, so the work done is
\[
\delta W = F\,\delta s = Fv\,\delta t.
\]
Dividing by $\delta t$ and letting $\delta t \to 0$:
\[
P = \frac{dW}{dt} = Fv.
\]

\begin{propriete}[Power of a force]
For a constant force $F$ whose point of application moves with velocity $v$ \emph{in the direction of the force}, the instantaneous power developed is
\[
P = Fv.
\]
More generally, if the force makes an angle $\theta$ with the velocity, $P = Fv\cos\theta$. This derivation is short and is frequently asked as a ``show that'' in the GCE.
\end{propriete}

\begin{exemplebox}[Power of a pulling force]
A horizontal force of \SI{400}{N} pulls a crate along a level floor at a steady speed of \SI{2.5}{m/s}. The power developed is
\[
P = Fv = 400 \times 2.5 = 1000\ \mathrm{W} = 1\ \mathrm{kW}.
\]
\end{exemplebox}

\courssub{Vehicles Moving with Constant Power}

A car engine working at constant power $P$ does \emph{not} exert a constant driving force. Since $P = Fv$, the driving force at speed $v$ is
\[
F = \frac{P}{v}.
\]
As the speed increases, the driving force \emph{decreases} hyperbolically. This single fact explains the whole behaviour of vehicles and is the heart of almost every GCE question on this topic.

\begin{attention}[The most common GCE error]
A vehicle with constant power does \textbf{NOT} have constant driving force. Never compute $F = P/v$ at one speed and then use that same $F$ at a different speed. Always ask yourself: ``at \emph{which} speed am I evaluating the driving force?''
\end{attention}

\begin{popfigure}
\begin{center}
\begin{tikzpicture}
\begin{axis}[
    width=11cm, height=7.5cm,
    axis lines=left,
    xlabel={$v\ (\mathrm{m\,s^{-1}})$}, ylabel={Force (N)},
    xmin=0, xmax=32, ymin=0, ymax=5200,
    xtick={0,5,10,15,20,25,30},
    ytick={0,1000,2000,3000,4000,5000},
    domain=2.5:32, samples=100,
    legend style={at={(0.97,0.97)}, anchor=north east, font=\small}
]
\addplot[very thick, popblue] {60000/x};
\addlegendentry{$F = P/v$ (driving force)}
\addplot[very thick, poporange, domain=0:32] {2000};
\addlegendentry{$R$ (resistance)}
\addplot[mark=*, popdark, only marks] coordinates {(30,2000)};
\draw[dashed, popmuted] (axis cs:30,0) -- (axis cs:30,2000);
\node[popdark, font=\small, align=center] at (axis cs:24,700) {$v_{\max}=P/R$};
\end{axis}
\end{tikzpicture}
\end{center}
\legende{Driving force $F=P/v$ (here $P = \SI{60}{kW}$) decreases hyperbolically with speed. It meets the constant resistance $R = \SI{2000}{N}$ at $v_{\max} = P/R = \SI{30}{m/s}$: beyond this speed the resistance exceeds the driving force, so the car cannot accelerate further.}
\end{popfigure}

\courssub{Equation of Motion and Maximum Speed on a Level Road}

Let a vehicle of mass $m$ move on a level road with engine power $P$, against a total resistance $R$ (friction, air resistance, etc.). The forces along the road are the driving force $P/v$ forwards and $R$ backwards. Newton's second law gives the \cle{equation of motion}:
\[
\frac{P}{v} - R = ma.
\]
At low speed $P/v$ is large, so $a > 0$ and the vehicle accelerates. As $v$ grows, $P/v$ falls towards $R$, so the acceleration falls towards zero. The speed approaches a limiting value, the \cle{maximum (terminal) speed}, reached when $a = 0$:
\[
\frac{P}{v_{\max}} = R \qquad\Longrightarrow\qquad v_{\max} = \frac{P}{R}.
\]

\begin{methode}[Maximum speed problems]
\begin{enumerate}
\item Convert all units: power to watts, speeds to $\mathrm{m\,s^{-1}}$, masses to kg.
\item Write the equation of motion $\dfrac{P}{v} - (\text{total opposing force}) = ma$.
\item For maximum speed, set $a = 0$ and solve $\dfrac{P}{v_{\max}} = \text{total resisting force}$.
\item If acceleration at a given speed is required, substitute that speed into $P/v$ and use Newton's second law.
\end{enumerate}
\end{methode}

\begin{exoresolu}[Car on a level road]
A car of mass \SI{900}{kg} moves along a level road against a constant resistance of \SI{600}{N}. Its engine works at a constant rate of \SI{24}{kW}. Find (a) the maximum speed of the car; (b) the acceleration of the car when its speed is \SI{20}{m/s}.
\tcblower
\textbf{Solution.} Convert: $P = 24\,000\ \mathrm{W}$.

(a) At maximum speed, $a = 0$, so
\[
\frac{P}{v_{\max}} = R \;\Longrightarrow\; v_{\max} = \frac{24\,000}{600} = 40\ \mathrm{m\,s^{-1}}.
\]

(b) At $v = 20\ \mathrm{m\,s^{-1}}$, the driving force is
\[
F = \frac{P}{v} = \frac{24\,000}{20} = 1200\ \mathrm{N}.
\]
Newton's second law:
\[
1200 - 600 = 900a \;\Longrightarrow\; a = \frac{600}{900} \approx 0.67\ \mathrm{m\,s^{-2}}.
\]
Note how the driving force at \SI{20}{m/s} (\SI{1200}{N}) differs from that at maximum speed (\SI{600}{N}): the force is \emph{not} constant.
\end{exoresolu}

\courssub{Motion on an Incline with Constant Power}

When the vehicle climbs a hill inclined at angle $\alpha$ to the horizontal, the component of its weight along the slope, $mg\sin\alpha$, acts \emph{down} the slope and must be added to the resistance. The equation of motion uphill is
\[
\frac{P}{v} - R - mg\sin\alpha = ma,
\]
so the maximum speed uphill satisfies
\[
\frac{P}{v_{\max}} = R + mg\sin\alpha.
\]
Going \emph{downhill}, the weight component assists the motion, so (with the engine still working) the equation becomes
\[
\frac{P}{v} + mg\sin\alpha - R = ma,
\qquad\text{and at maximum speed}\qquad
\frac{P}{v_{\max}} = R - mg\sin\alpha.
\]
(If the engine is switched off downhill, simply delete the $P/v$ term.)

\begin{popfigure}
\begin{center}
\begin{tikzpicture}[scale=1.05, >=stealth]
\coordinate (A) at (0,0);
\coordinate (B) at (7,0);
\coordinate (C) at (7,2.45);
\draw[thick, popdark] (A) -- (B) -- (C) -- cycle;
\draw[popmuted] (5.6,0) arc[start angle=0, end angle=19.3, radius=1.4];
\node[popmuted] at (5.1,0.28) {$\alpha$};
% car body on the slope
\draw[fill=popblueL, draw=popblue, thick, rotate around={19.3:(4.6,1.61)}]
    (3.6,1.61) rectangle (5.6,2.31);
\draw[fill=white, draw=popblue, rotate around={19.3:(4.6,1.61)}] (3.95,1.41) circle (0.2);
\draw[fill=white, draw=popblue, rotate around={19.3:(4.6,1.61)}] (5.25,1.41) circle (0.2);
% forces from centre of mass
\coordinate (G) at (4.6,2.13);
\draw[->, very thick, popgreen] (G) -- ++(19.3:2.4) node[above, sloped, pos=0.95, font=\small]{$F=\dfrac{P}{v}$};
\draw[->, very thick, poporange] (G) -- ++(199.3:1.7) node[below, sloped, pos=0.9, font=\small]{$R$};
\draw[->, very thick, poppurple] (G) -- ++(-90:2.0) node[below, font=\small]{$mg$};
\draw[->, thick, poppurple, dashed] (G) -- ++(250.7:1.55) node[pos=1.05, below, font=\small]{$mg\sin\alpha$};
\draw[->, thick, popmuted, dashed] (G) -- ++(160.7:0.85) node[pos=1.15, left, font=\small]{$mg\cos\alpha$};
\end{tikzpicture}
\end{center}
\legende{Forces on a vehicle of mass $m$ climbing a slope of angle $\alpha$ at speed $v$: driving force $F = P/v$ up the slope, resistance $R$ and weight component $mg\sin\alpha$ down the slope. At maximum speed these balance: $P/v_{\max} = R + mg\sin\alpha$.}
\end{popfigure}

\begin{exoresolu}[Lorry climbing a hill against a speed-dependent resistance]
A lorry of mass \SI{4000}{kg} has an engine working at a constant rate of \SI{80}{kW}. The resistance to motion is $R = 40v$ newtons, where $v$ is the speed in $\mathrm{m\,s^{-1}}$. Take $g = 10\ \mathrm{m\,s^{-2}}$.
\begin{enumerate}
\item Find the maximum speed of the lorry on a level road.
\item The lorry climbs a hill inclined at an angle $\alpha$ to the horizontal, where $\sin\alpha = \tfrac{1}{20}$. Find its maximum speed up the hill.
\item Find the acceleration of the lorry when it is climbing this hill at \SI{10}{m/s}.
\end{enumerate}
\tcblower
\textbf{Solution.} $P = 80\,000\ \mathrm{W}$, $m = 4000\ \mathrm{kg}$.

\textbf{1.} Level road, $a = 0$:
\[
\frac{80\,000}{v} = 40v \;\Longrightarrow\; v^2 = 2000 \;\Longrightarrow\; v_{\max} = \sqrt{2000} \approx 44.7\ \mathrm{m\,s^{-1}}.
\]

\textbf{2.} Uphill, the weight component is
\[
mg\sin\alpha = 4000 \times 10 \times \tfrac{1}{20} = 2000\ \mathrm{N}.
\]
At maximum speed:
\[
\frac{80\,000}{v} = 40v + 2000
\;\Longrightarrow\;
40v^2 + 2000v - 80\,000 = 0
\;\Longrightarrow\;
v^2 + 50v - 2000 = 0.
\]
By the quadratic formula:
\[
v = \frac{-50 + \sqrt{2500 + 8000}}{2} = \frac{-50 + \sqrt{10\,500}}{2} \approx \frac{-50 + 102.5}{2} \approx 26.2\ \mathrm{m\,s^{-1}}.
\]
(The negative root is rejected since $v > 0$.)

\textbf{3.} At $v = 10\ \mathrm{m\,s^{-1}}$ uphill, the driving force is $P/v = 80\,000/10 = 8000\ \mathrm{N}$ and the resistance is $R = 40 \times 10 = 400\ \mathrm{N}$. Newton's second law:
\[
8000 - 400 - 2000 = 4000a
\;\Longrightarrow\;
a = \frac{5600}{4000} = 1.4\ \mathrm{m\,s^{-2}}.
\]
\end{exoresolu}

\begin{attention}[Exam tip --- units first]
Before substituting anything, convert: power to \textbf{watts} (multiply kW by $1000$), speeds to $\mathbf{m\,s^{-1}}$ (divide km/h by $3.6$), masses to \textbf{kg}. Mixing kW with N and m/s is the fastest way to lose marks in this topic.
\end{attention}

\begin{aretenir}[Key points]
\begin{itemize}
\item Power is the rate of doing work: $P_{\text{av}} = W/t$ (average), $P = dW/dt$ (instantaneous); unit the watt, $1\ \mathrm{kW} = 1000\ \mathrm{W}$.
\item For a force $F$ moving its point of application at velocity $v$ in the same direction: $P = Fv$.
\item A vehicle at constant power has driving force $F = P/v$, which \emph{decreases} as speed increases.
\item Equation of motion: $\dfrac{P}{v} - R = ma$ (level road); add $mg\sin\alpha$ to the opposing forces uphill, subtract it downhill.
\item Maximum speed occurs when $a = 0$: on level ground $v_{\max} = P/R$; uphill $\dfrac{P}{v_{\max}} = R + mg\sin\alpha$; downhill $\dfrac{P}{v_{\max}} = R - mg\sin\alpha$.
\end{itemize}
\end{aretenir}

\begin{exercice}[Maximum speed with speed-dependent resistance]
A car of mass \SI{1200}{kg} works at a constant rate of \SI{48}{kW}. The total resistance to motion is $R = 30v$ newtons on a level road, where $v$ is in $\mathrm{m\,s^{-1}}$. Take $g = 10\ \mathrm{m\,s^{-2}}$.
\begin{enumerate}
\item Find the maximum speed of the car on the level road.
\item Find the acceleration when the car is travelling at \SI{20}{m/s} on the level road.
\item The car now descends a hill with $\sin\alpha = \tfrac{1}{40}$, with the engine still working at \SI{48}{kW}. Show that the maximum speed downhill $v$ satisfies $v^2 - 10v - 1600 = 0$, and find it.
\end{enumerate}
\end{exercice}

\begin{corrige}[Exercise 1]
$P = 48\,000\ \mathrm{W}$.
\textbf{1.} At maximum speed on the level road, $\dfrac{48\,000}{v} = 30v \Rightarrow v^2 = 1600 \Rightarrow v_{\max} = 40\ \mathrm{m\,s^{-1}}$.
\textbf{2.} At $v = 20\ \mathrm{m\,s^{-1}}$: driving force $F = 48\,000/20 = 2400\ \mathrm{N}$, resistance $R = 30 \times 20 = 600\ \mathrm{N}$, so
\[
a = \frac{2400 - 600}{1200} = 1.5\ \mathrm{m\,s^{-2}}.
\]
\textbf{3.} Downhill, $mg\sin\alpha = 1200 \times 10 \times \tfrac{1}{40} = 300\ \mathrm{N}$ assists the motion. At maximum speed, driving force plus weight component balance resistance:
\[
\frac{48\,000}{v} + 300 = 30v
\;\Longrightarrow\;
48\,000 + 300v = 30v^2
\;\Longrightarrow\;
v^2 - 10v - 1600 = 0,
\]
as required. Solving:
\[
v = \frac{10 + \sqrt{100 + 6400}}{2} = \frac{10 + \sqrt{6500}}{2} \approx \frac{10 + 80.6}{2} \approx 45.3\ \mathrm{m\,s^{-1}}.
\]
(The negative root is rejected.) Note that the maximum speed downhill exceeds that on the level road, as expected.
\end{corrige}

\begin{savaistu}[Why lorries slow down on hills]
The relation $F = P/v$ explains a familiar sight on the road from Yaound\'e to Bafoussam: heavy lorries crawling uphill at a fraction of their speed on the flat. The engine's power is fixed, so to produce the large driving force needed against gravity and resistance, the speed $v = P/F$ must be small --- which is exactly why drivers change down to a lower gear before a steep climb: lower gears let the wheels turn slowly while the engine keeps delivering its full power.
\end{savaistu}

\coursec{Pumping Water, Water Jets and Applications of Work, Energy and Power}

In the previous section we discovered that power is the \emph{rate of doing work}, and that a force $F$ moving its point of application at velocity $v$ delivers the instantaneous power $P = Fv$. We now apply these ideas to one of the most frequent families of GCE Advanced Level problems: \cle{pumps} raising water, \cle{water jets} and fire hoses, and engines working against gravity and resistance. These problems look different from block-and-slope questions, but they are governed by exactly the same principle: \emph{power is the energy given to the system each second}.

\courssub{Modelling a Pump Raising Water}

\textbf{From a static load to a continuous flow.}\par
Until now, the objects we lifted were rigid bodies of fixed mass $m$: raising them through a height $h$ required the work $W = mgh$, done once and for all. A pump is different: it does not lift one mass, it lifts a \emph{continuous stream} of water. Every second, a certain mass of water arrives at the pump and is delivered, higher up or faster, at the outlet. The natural quantity is therefore not the mass but the \cle{mass flow rate}.

\begin{definition}[Mass flow rate]
The \cle{mass flow rate} $\dot{m}$ of a fluid in motion is the mass of fluid passing a given cross-section per unit time:
\[
\dot{m} = \frac{\Delta m}{\Delta t}, \qquad \text{unit: } \mathrm{kg\,s^{-1}}.
\]
For a jet or pipe of cross-sectional area $A$ in which water of density $\rho$ moves at speed $v$, the volume crossing a section in one second is $Av$ (a cylinder of length $v$), so
\[
\boxed{\dot{m} = \rho A v}.
\]
\end{definition}

\begin{popfigure}
\begin{center}
\begin{tikzpicture}[scale=0.95, every node/.style={font=\small}]
% ground
\draw[thick, popdark] (-0.5,0) -- (10.5,0);
\foreach \x in {-0.3,0.3,...,10.3}{\draw[popmuted] (\x,0) -- (\x-0.25,-0.2);}
% reservoir
\draw[fill=popblueL, draw=popblue] (0.3,0) rectangle (2.7,0.9);
\node[popdark] at (1.5,0.45) {reservoir};
% pump
\draw[fill=poporange!25, draw=poporange, thick] (3.4,0) rectangle (4.8,0.9);
\node[popdark] at (4.1,0.45) {pump};
\node[popdark, below=2pt] at (4.1,0) {power input $P$};
% pipe up
\draw[fill=popblueL, draw=popblue] (4.8,0.25) -- (5.6,0.25) -- (5.6,4.2) -- (6.4,4.2) -- (6.4,0.25) -- (7.6,0.25);
\draw[popblue] (4.8,0.65) -- (5.6,0.65);
\draw[popblue] (6.4,0.65) -- (7.6,0.65);
% tank
\draw[fill=popblueL, draw=popblue, thick] (7.6,2.6) rectangle (10.2,4.4);
\draw[popblue] (7.6,3.6) -- (10.2,3.6);
\node[popdark] at (8.9,4.05) {overhead tank};
% water jet into tank
\draw[->, very thick, popblue] (6.4,3.9) -- (7.5,3.9) node[midway, above, popdark]{$v$};
% height arrow
\draw[<->, thick, poppurple] (5.0,0.9) -- (5.0,3.9) node[midway, left, popdark]{$h$};
% labels
\node[popdark, align=center, text width=4cm] at (6.0,5.0) {pipe of cross-section $A$\\ density $\rho$, \; $\dot{m}=\rho A v$};
\draw[->, popmuted] (6.0,4.75) -- (6.0,4.3);
\end{tikzpicture}
\end{center}
\legende{A pump raises water from a ground reservoir to an overhead tank at height $h$; the water leaves the pipe of cross-section $A$ at speed $v$.}
\end{popfigure}

\textbf{Energy given to each kilogram of water.}\par
Consider what the pump must do for each kilogram of water passing through it:
\begin{itemize}
\item raise it through the height $h$: this costs $gh$ joules of gravitational potential energy per kilogram;
\item if the water leaves the pipe at speed $v$ (having started essentially at rest in the reservoir), give it kinetic energy $\tfrac{1}{2}v^{2}$ joules per kilogram.
\end{itemize}
Hence the energy supplied \emph{per kilogram} is $gh + \tfrac{1}{2}v^{2}$. Since $\dot{m}$ kilograms pass each second, the energy supplied \emph{per second} — that is, the power — follows immediately.

\begin{propriete}[Power of a pump and of a water jet]
A pump handling a mass flow rate $\dot{m}$ (in $\mathrm{kg\,s^{-1}}$) delivers a useful power
\[
\boxed{P = \dot{m}gh + \tfrac{1}{2}\dot{m}v^{2}}
\]
where $h$ is the height through which the water is raised and $v$ the speed at which it leaves. In particular:
\begin{itemize}
\item water merely \emph{lifted} through $h$ (delivered at negligible speed): $P = \dot{m}gh$;
\item a \emph{water jet} of cross-section $A$, density $\rho$, speed $v$ (no lifting): with $\dot{m} = \rho A v$,
\[
P = \tfrac{1}{2}\dot{m}v^{2} = \tfrac{1}{2}\rho A v^{3}.
\]
\end{itemize}
(Proof examinable as a derivation: it is a direct application of the work--energy theorem per unit time.)
\end{propriete}

\begin{methode}[Solving a pumping or jet problem]
\begin{enumerate}
\item \textbf{Find the mass flow rate.} For a jet or pipe, $\dot{m} = \rho A v$; for a tank filled in time $t$, $\dot{m} = \rho V/t$. Remember $\rho_{\text{water}} = 1000\ \mathrm{kg\,m^{-3}}$ and convert areas to $\mathrm{m^{2}}$.
\item \textbf{Energy per kilogram.} Add $gh$ (lifting) and $\tfrac{1}{2}v^{2}$ (speed given), each only if present.
\item \textbf{Multiply.} Useful power $P = \dot{m}\left(gh + \tfrac{1}{2}v^{2}\right)$.
\item \textbf{Account for efficiency} if given: $P_{\text{input}} = P_{\text{useful}}/\eta$.
\end{enumerate}
\end{methode}

\begin{attention}[Do not forget the kinetic energy term]
When water leaves a pipe or hose \emph{at speed}, the pump (or the fire engine) supplies \emph{both} potential and kinetic energy. Using only $\dot{m}gh$ underestimates the power. Conversely, if the water is simply poured gently into the tank, the term $\tfrac{1}{2}\dot{m}v^{2}$ must not be invented. Read the statement carefully: ``the water issues from the pipe at $6\ \mathrm{m\,s^{-1}}$'' means the kinetic term is required.
\end{attention}

\begin{exoresolu}[Pump filling an overhead tank]
A pump raises water from a ground-level reservoir to an overhead tank whose inlet is $12\ \mathrm{m}$ above the reservoir surface. The water flows through a pipe of cross-sectional area $8.0\times10^{-3}\ \mathrm{m^{2}}$ and leaves the pipe at $5.0\ \mathrm{m\,s^{-1}}$. Take $\rho = 1000\ \mathrm{kg\,m^{-3}}$ and $g = 9.81\ \mathrm{m\,s^{-2}}$. Find the useful power of the pump.
\tcblower
\textbf{Step 1: mass flow rate.}
\[
\dot{m} = \rho A v = 1000 \times 8.0\times10^{-3} \times 5.0 = 40\ \mathrm{kg\,s^{-1}}.
\]
\textbf{Step 2: energy per kilogram.} Lifting: $gh = 9.81 \times 12 = 117.72\ \mathrm{J\,kg^{-1}}$. Kinetic: $\tfrac{1}{2}v^{2} = \tfrac{1}{2}(5.0)^{2} = 12.5\ \mathrm{J\,kg^{-1}}$.
\textbf{Step 3: power.}
\[
P = \dot{m}\left(gh + \tfrac{1}{2}v^{2}\right) = 40\,(117.72 + 12.5) = 40 \times 130.22 \approx 5.21\times10^{3}\ \mathrm{W}.
\]
The pump delivers a useful power of about $\boxed{5.2\ \mathrm{kW}}$, of which $40\times 12.5 = 500\ \mathrm{W}$ (about $10\%$) goes into the kinetic energy of the jet.
\end{exoresolu}

\begin{exoresolu}[Power of a fire-hose jet]
Water issues horizontally from a hose nozzle of diameter $7.0\ \mathrm{cm}$ at a speed of $25\ \mathrm{m\,s^{-1}}$. Find the kinetic power of the jet. ($\rho = 1000\ \mathrm{kg\,m^{-3}}$.)
\tcblower
The cross-sectional area is
\[
A = \pi r^{2} = \pi (0.035)^{2} = 3.848\times10^{-3}\ \mathrm{m^{2}}.
\]
Hence $\dot{m} = \rho A v = 1000 \times 3.848\times10^{-3} \times 25 = 96.2\ \mathrm{kg\,s^{-1}}$, and
\[
P = \tfrac{1}{2}\dot{m}v^{2} = \tfrac{1}{2}\rho A v^{3} = \tfrac{1}{2}\times 1000 \times 3.848\times10^{-3}\times 25^{3} \approx 3.0\times10^{4}\ \mathrm{W}.
\]
The jet carries about $\boxed{30\ \mathrm{kW}}$ of kinetic power. Note the $v^{3}$ dependence: doubling the speed multiplies the power by $8$.
\end{exoresolu}

\courssub{Efficiency of Pumps and Engines}

No real machine converts all the energy it consumes into useful work: friction in bearings, turbulence in pipes and heat losses all take their share. This is quantified by the \cle{efficiency}.

\begin{definition}[Efficiency]
The \cle{efficiency} $\eta$ of a machine is
\[
\eta = \frac{\text{useful power output}}{\text{power input}} = \frac{\text{useful energy output}}{\text{energy input}},
\]
usually expressed as a percentage. Consequently $P_{\text{useful}} = \eta\, P_{\text{input}}$, i.e.\ the input power required is $P_{\text{input}} = P_{\text{useful}}/\eta$.
\end{definition}

\begin{exemplebox}[Motor driving a pump]
The pump of the first worked example is driven by an electric motor of efficiency $80\%$. The electrical power consumed is
\[
P_{\text{in}} = \frac{5.21\ \mathrm{kW}}{0.80} \approx 6.5\ \mathrm{kW}.
\]
The missing $1.3\ \mathrm{kW}$ is dissipated as heat and sound in the motor and pump.
\end{exemplebox}

\courssub{Choosing Between Energy and Power Methods}

By now you have two toolboxes that overlap: the \emph{work--energy theorem} (and conservation of energy) on one side, \emph{power} ($P = Fv$, $P = \dot{W}$) on the other. Choosing the right one saves precious exam minutes.

\begin{methode}[Which method for which problem?]
\begin{itemize}
\item \textbf{Speeds and displacements, no time involved} (``find the speed at the bottom'', ``how far does it slide''): use the \emph{work--energy theorem} or conservation of energy.
\item \textbf{Rates, times, continuous flow} (``how long to fill the tank'', ``power developed'', ``mass per second''): use \emph{power}.
\item \textbf{Constant speed against resistance} (vehicle at steady speed): the driving force balances resistance, and $P = Fv$ links them directly.
\item \textbf{Acceleration at a given instant} (vehicle speeding up): write $P = Fv$ for the tractive force, then Newton's second law $F - R = ma$.
\end{itemize}
\end{methode}

\begin{attention}[GCE exam tip]
In structured questions, later parts almost always reuse earlier results. Keep \emph{exact} values (fractions, surds, unrounded decimals) throughout your working and round only the final answer to the accuracy requested. A premature rounding in part (a) can cost all the accuracy marks of parts (b) and (c).
\end{attention}

\courssub{GCE-Style Structured Problems}

\begin{exoresolu}[Vehicle on a hill with resistance]
A car of mass $1200\ \mathrm{kg}$ climbs a hill inclined at an angle $\alpha$ to the horizontal, where $\sin\alpha = \tfrac{1}{14}$. A constant resistance to motion of $600\ \mathrm{N}$ acts on the car. Take $g = 9.8\ \mathrm{m\,s^{-2}}$.
\begin{enumerate}
\item[(a)] The car climbs at a steady speed of $20\ \mathrm{m\,s^{-1}}$. Find the power developed by the engine.
\item[(b)] With the engine now developing $45\ \mathrm{kW}$, find the acceleration of the car at the instant its speed is $20\ \mathrm{m\,s^{-1}}$ up the hill.
\end{enumerate}
\tcblower
\textbf{(a)} At steady speed the tractive force $F$ balances the total opposing force:
\[
F = R + mg\sin\alpha = 600 + 1200 \times 9.8 \times \tfrac{1}{14} = 600 + 840 = 1440\ \mathrm{N}.
\]
Then
\[
P = Fv = 1440 \times 20 = 28\,800\ \mathrm{W} = 28.8\ \mathrm{kW}.
\]
\textbf{(b)} With $P = 45\,000\ \mathrm{W}$ at $v = 20\ \mathrm{m\,s^{-1}}$, the new tractive force is
\[
F' = \frac{P}{v} = \frac{45\,000}{20} = 2250\ \mathrm{N}.
\]
Newton's second law up the slope:
\[
F' - R - mg\sin\alpha = ma \;\Longrightarrow\; 2250 - 1440 = 1200\,a,
\]
so $a = \dfrac{810}{1200} = 0.675\ \mathrm{m\,s^{-2}}$.
\end{exoresolu}

\begin{exoresolu}[Pump, tank and time]
A small water pump of useful power $1.5\ \mathrm{kW}$ fills an empty tank of capacity $2.4\ \mathrm{m^{3}}$ situated $8.0\ \mathrm{m}$ above the water source. The water enters the tank at negligible speed. Taking $\rho = 1000\ \mathrm{kg\,m^{-3}}$ and $g = 9.81\ \mathrm{m\,s^{-2}}$, find the time needed to fill the tank.
\tcblower
The total potential energy to supply is
\[
E = mgh = \rho V g h = 1000 \times 2.4 \times 9.81 \times 8.0 = 188\,352\ \mathrm{J}.
\]
Since $P = \dfrac{E}{t}$,
\[
t = \frac{E}{P} = \frac{188\,352}{1500} \approx 125.6\ \mathrm{s} \approx 2\ \mathrm{min}\ 6\ \mathrm{s}.
\]
\end{exoresolu}

\begin{exoresolu}[Springs and energy combined]
A light spring of natural length $0.50\ \mathrm{m}$ and modulus of elasticity $60\ \mathrm{N}$ hangs vertically. A mass of $0.80\ \mathrm{kg}$ is attached to its lower end and released from rest with the spring just unstretched. Take $g = 9.8\ \mathrm{m\,s^{-2}}$ and $\lambda = 60\ \mathrm{N}$, natural length $l = 0.50\ \mathrm{m}$.
\begin{enumerate}
\item[(a)] Using energy, find the extension of the spring when the mass first comes momentarily to rest.
\item[(b)] Find the speed of the mass when it has fallen $0.20\ \mathrm{m}$.
\end{enumerate}
\tcblower
\textbf{(a)} Let $x$ be the maximum extension. The gravitational potential energy lost equals the elastic energy stored (kinetic energy is zero at both ends):
\[
mgx = \frac{\lambda x^{2}}{2l} \;\Longrightarrow\; x = \frac{2mgl}{\lambda} = \frac{2 \times 0.80 \times 9.8 \times 0.50}{60} = 0.1307\ \mathrm{m} \approx 0.13\ \mathrm{m}.
\]
\textbf{(b)} The maximum extension is $0.13\ \mathrm{m}$, so the mass never reaches a fall of $0.20\ \mathrm{m}$; the question likely contains a misprint. We illustrate the method for a fall of $y = 0.10\ \mathrm{m}$:
\[
mgy = \tfrac{1}{2}mv^{2} + \frac{\lambda y^{2}}{2l}
\;\Longrightarrow\;
0.80 \times 9.8 \times 0.10 = \tfrac{1}{2}(0.80)v^{2} + \frac{60(0.10)^{2}}{2 \times 0.50},
\]
\[
0.784 = 0.40\,v^{2} + 0.60 \;\Longrightarrow\; v^{2} = 0.46, \quad v \approx 0.68\ \mathrm{m\,s^{-1}}.
\]
This illustrates the exam tip: energy methods handle springs and gravity in one line, and checking consistency of data is part of the exercise.
\end{exoresolu}

\begin{exercice}[Fire pump in Yaounde]
A pump draws water from a well $6.0\ \mathrm{m}$ below ground and delivers it through a nozzle of cross-section $2.0\times10^{-3}\ \mathrm{m^{2}}$ at $12\ \mathrm{m\,s^{-1}}$. The pump-motor unit has efficiency $70\%$. Find (a) the mass flow rate, (b) the useful power, (c) the electrical power consumed. ($\rho = 1000\ \mathrm{kg\,m^{-3}}$, $g = 9.81\ \mathrm{m\,s^{-2}}$.)
\end{exercice}

\begin{corrige}[Exercise 1]
(a) $\dot{m} = \rho A v = 1000 \times 2.0\times10^{-3} \times 12 = 24\ \mathrm{kg\,s^{-1}}$.
(b) $P_{\text{useful}} = \dot{m}gh + \tfrac{1}{2}\dot{m}v^{2} = 24(9.81\times 6.0) + \tfrac{1}{2}(24)(144) = 1412.6 + 1728 = 3140.6\ \mathrm{W} \approx 3.1\ \mathrm{kW}$.
(c) $P_{\text{in}} = \dfrac{3140.6}{0.70} \approx 4.5\ \mathrm{kW}$.
\end{corrige}

\courssub{Revision Grid of the Chapter}

\begin{center}
\begin{tabularx}{\linewidth}{|l|X|X|}
\hline
\rowcolor{popblueL} \textbf{Quantity / Law} & \textbf{Formula} & \textbf{Conditions of validity} \\
\hline
Work of a constant force & $W = Fs\cos\theta$ & Force constant; $\theta$ angle between $\vec{F}$ and displacement \\
\hline
Work against gravity & $W = mgh$ & Near Earth's surface, $g$ constant; $h$ vertical rise \\
\hline
Work of a variable force & $W = \displaystyle\int_{x_1}^{x_2} F(x)\,dx$ & Force depends on position; area under $F$--$x$ graph \\
\hline
Elastic potential energy & $E_e = \dfrac{\lambda x^{2}}{2l}$ & Spring of modulus $\lambda$, natural length $l$, extension $x$ \\
\hline
Kinetic energy & $E_k = \tfrac{1}{2}mv^{2}$ & Translation of a particle of mass $m$ \\
\hline
Gravitational potential energy & $E_p = mgh$ & Measured from a chosen reference level \\
\hline
Work--energy theorem & $W_{\text{all forces}} = \Delta E_k$ & Always valid (resultant work) \\
\hline
Conservation of energy & $E_k + E_p = \text{const}$ & Only conservative forces do work \\
\hline
Power & $P = \dfrac{dW}{dt} = Fv$ & $P = Fv$ for force $F$ at speed $v$ in its direction \\
\hline
Pump / jet power & $P = \dot{m}gh + \tfrac{1}{2}\dot{m}v^{2}$, \; $\dot{m} = \rho A v$ & Steady flow; include each term only when present \\
\hline
Efficiency & $\eta = \dfrac{P_{\text{useful}}}{P_{\text{input}}}$ & Usually quoted as a percentage \\
\hline
\end{tabularx}
\end{center}

\begin{popfigure}
\begin{center}
\begin{tikzpicture}[node distance=0.55cm, every node/.style={font=\small}]
\node[draw=popblue, fill=popblueL, rounded corners, thick, align=center, text width=3.2cm] (work) {\textbf{Work}\\ $W = Fs\cos\theta$,\; $W = \int F\,dx$};
\node[draw=poppurple, fill=poppurple!12, rounded corners, thick, align=center, text width=3.4cm, right=of work] (wet) {\textbf{Work--Energy Theorem}\\ $W = \Delta E_k$};
\node[draw=popgreen, fill=popgreenL, rounded corners, thick, align=center, text width=3.4cm, right=of wet] (energy) {\textbf{Energy}\\ $E_k = \tfrac12 mv^2$,\; $E_p = mgh$};
\node[draw=poporange, fill=poporange!15, rounded corners, thick, align=center, text width=3.0cm, below=0.8cm of wet] (power) {\textbf{Power}\\ $P = \dot W = Fv$};
\node[draw=popdark, fill=gray!12, rounded corners, thick, align=center, text width=4.6cm, right=of power] (apps) {\textbf{Applications}\\ pumps, jets, vehicles, springs};
\draw[->, thick, popdark] (work) -- (wet);
\draw[->, thick, popdark] (wet) -- (energy);
\draw[->, thick, popdark] (wet) -- (power);
\draw[->, thick, popdark] (energy) -- (power);
\draw[->, thick, popdark] (power) -- (apps);
\end{tikzpicture}
\end{center}
\legende{The logical chain of the chapter: work leads, through the work--energy theorem, to energy, power and their applications.}
\end{popfigure}

\begin{aretenir}[Key points of the section]
\begin{itemize}
\item For continuous flow, work with the \cle{mass flow rate} $\dot{m} = \rho A v$ and the energy given to \emph{each kilogram}.
\item Pump power: $P = \dot{m}gh + \tfrac{1}{2}\dot{m}v^{2}$; a pure jet has $P = \tfrac{1}{2}\rho A v^{3}$ — note the $v^{3}$.
\item Efficiency links input and output: $P_{\text{input}} = P_{\text{useful}}/\eta$.
\item Energy methods for speeds and displacements; power methods for rates and times; $P = Fv$ with Newton's second law for vehicles.
\item Keep exact values until the final line of a structured question.
\end{itemize}
\end{aretenir}

\begin{savaistu}[Water power in Cameroon]
Cameroon's great hydroelectric dams — Song Loulou and Edea on the Sanaga river — are giant, perfected versions of the pump problems of this section, run in reverse: instead of spending power to raise water, they collect the power $\dot{m}gh$ released by water \emph{falling} through the height of the dam. The same formula you have just used for a domestic pump governs the electricity supply of entire cities — a beautiful illustration that the physics of the GCE classroom is the physics of the nation.
\end{savaistu}

\newpage

\coursec{Integration activity}

\courssub{Aims of the activity}

By the end of this group activity, each student should be able to:

\begin{itemize}
\item compute a \cle{mass flow rate} from a volume delivered per unit time;
\item calculate the \cle{work done against gravity} when a mass is raised through a given height;
\item calculate the \cle{kinetic energy} communicated to a moving fluid;
\item apply the \cle{work--energy theorem} and the \cle{conservation of energy} to obtain a total energy balance;
\item deduce a \cle{minimum power}, then correct it for the \cle{efficiency} of a real machine;
\item make and justify an engineering choice from a catalogue, and verify a design by an independent calculation;
\item present and defend a scientific argument in front of the class.
\end{itemize}

\courssub{Situation}

The Mfoundi neighbourhood, in the heart of Yaoundé, is supplied with water from a community well. The water level in the well is $12\ \mathrm{m}$ below the point where the delivery pipe discharges into an overhead tank. The neighbourhood committee wants to install an electric pump that satisfies the following requirements:

\begin{itemize}
\item the pump must deliver $600$ litres of water per hour;
\item the water must leave the delivery pipe with a speed of $2\ \mathrm{m\,s^{-1}}$;
\item the overall efficiency of the pump and its motor is $70\%$;
\item the overhead storage tank has a capacity of $2000$ litres.
\end{itemize}

Take $g = 9.81\ \mathrm{m\,s^{-2}}$ and the density of water $\rho = 1000\ \mathrm{kg\,m^{-3}}$. Recall that $1$ litre $= 10^{-3}\ \mathrm{m^{3}}$, so that $1$ litre of water has a mass of $1\ \mathrm{kg}$.

The local dealer proposes the following catalogue of motors:

\begin{center}
\begin{tabularx}{\linewidth}{|l|X|X|}
\hline
\rowcolor{popblueL} \textbf{Model} & \textbf{Rated power (W)} & \textbf{Price (FCFA)} \\
\hline
M1 & $120$ & $45\,000$ \\
\hline
M2 & $150$ & $58\,000$ \\
\hline
M3 & $200$ & $72\,000$ \\
\hline
M4 & $250$ & $85\,000$ \\
\hline
\end{tabularx}
\end{center}

Working in groups, you are the engineering team of the committee. You must determine the energy the pump must supply every second, choose the cheapest motor that can do the job, and verify your design. Each group will present its complete \cle{energy balance} on a poster, and the class will compare the designs.

\courssub{Tasks}

\begin{enumerate}
\item \textbf{Mass flow rate.} Convert the delivery of $600$ litres per hour into a mass of water pumped \emph{per second}. Denote this mass flow rate by $q$ (in $\mathrm{kg\,s^{-1}}$).
\item \textbf{Work against gravity.} Each kilogram of water must be raised through $12\ \mathrm{m}$. Write down the work done against gravity to raise a mass $m$ through a height $h$, and deduce the work done against gravity \emph{per second} (i.e. the potential energy gained by the water each second).
\item \textbf{Kinetic energy.} The water arrives at the top with speed $v = 2\ \mathrm{m\,s^{-1}}$. Write down the kinetic energy of a mass $m$ moving at speed $v$, and deduce the kinetic energy given to the water \emph{per second}.
\item \textbf{Minimum pump power.} Using the work--energy theorem, explain why the minimum useful power of the pump is the sum of the two quantities found in questions 2 and 3. Compute this minimum power $P_{\min}$.
\item \textbf{Choice of motor.} The pump--motor set has an efficiency of $70\%$. Explain why the electrical power required is $P_{\mathrm{elec}} = \dfrac{P_{\min}}{0.7}$. Compute $P_{\mathrm{elec}}$ and choose, from the catalogue, the \emph{cheapest} motor whose rated power is sufficient. Justify your choice.
\item \textbf{Verification.} Compute the time needed to fill the $2000$-litre tank at the delivery rate of $600$ litres per hour. Give your answer in hours and minutes.
\item \textbf{Poster.} Present your full energy balance (mass flow rate, potential energy per second, kinetic energy per second, minimum power, electrical power, chosen motor, filling time) on a poster, with one sentence justifying each step.
\end{enumerate}

\courssub{Model answer}

\textbf{1. Mass flow rate.} Since $1$ litre of water has a mass of $1\ \mathrm{kg}$, a delivery of $600$ litres per hour corresponds to $600\ \mathrm{kg}$ per hour. Hence
\[
q=\frac{600}{3600}=\frac{1}{6}\approx 0.1667\ \mathrm{kg\,s^{-1}}.
\]

\textbf{2. Work done against gravity per second.} The work needed to raise a mass $m$ through a height $h$ is $W = mgh$. In one second the mass raised is $q = \tfrac{1}{6}\ \mathrm{kg}$, so the potential energy gained per second is
\[
\dot{E}_{p}=q\,g\,h=\frac{1}{6}\times 9.81\times 12 = 19.62\ \mathrm{W}.
\]

\textbf{3. Kinetic energy per second.} The kinetic energy of a mass $m$ at speed $v$ is $E_{k}=\tfrac{1}{2}mv^{2}$. In one second a mass $q$ is given the speed $v = 2\ \mathrm{m\,s^{-1}}$, so the kinetic energy supplied per second is
\[
\dot{E}_{k}=\frac{1}{2}\,q\,v^{2}=\frac{1}{2}\times \frac{1}{6}\times 2^{2}=\frac{1}{3}\approx 0.333\ \mathrm{W}.
\]

\textbf{4. Minimum pump power.} By the work--energy theorem, the total work done by the pump on each kilogram of water equals the change in its mechanical energy, namely the gain in gravitational potential energy plus the gain in kinetic energy. Therefore the minimum useful power is
\[
P_{\min}=\dot{E}_{p}+\dot{E}_{k}=19.62+0.333\approx 19.95\ \mathrm{W}.
\]

\textbf{5. Choice of motor.} Only $70\%$ of the electrical power consumed is converted into useful mechanical work on the water, so
\[
P_{\mathrm{elec}}=\frac{P_{\min}}{0.7}=\frac{19.95}{0.7}\approx 28.5\ \mathrm{W}.
\]
Every motor in the catalogue has a rated power greater than $28.5\ \mathrm{W}$; the cheapest adequate model is therefore \textbf{M1} ($120\ \mathrm{W}$, $45\,000$ FCFA). Its rated power is comfortably above the requirement, which also gives a safety margin for friction in the pipes.

\textbf{6. Verification.} At $600$ litres per hour, the time to fill the $2000$-litre tank is
\[
t=\frac{2000}{600}=\frac{10}{3}\ \mathrm{h}\approx 3\ \mathrm{h}\ 20\ \mathrm{min}.
\]

\textbf{7. Conclusion.} The design is consistent: a $120\ \mathrm{W}$ motor (model M1) pumps $600$ litres per hour to a height of $12\ \mathrm{m}$ with an exit speed of $2\ \mathrm{m\,s^{-1}}$, and fills the $2000$-litre tank in about $3\ \mathrm{h}\ 20\ \mathrm{min}$.

\begin{attention}[Common errors to avoid]
\begin{itemize}
\item Forgetting to convert hours into seconds: the mass flow rate must be in $\mathrm{kg\,s^{-1}}$, not $\mathrm{kg\,h^{-1}}$.
\item Omitting the kinetic energy term: the water does not only rise, it also leaves the pipe \emph{moving}.
\item Dividing by the efficiency instead of multiplying: since the machine wastes energy, the electrical power must be \emph{larger} than the useful power, so we divide $P_{\min}$ by $0.7$.
\end{itemize}
\end{attention}

\newpage

\coursec{Methods and worked examples}

\coursec{Work, Energy and Power}

\courssub{Work Done by a Constant Force and Work Done Against Gravity}

\begin{definition}[Work of a constant force]
Let a constant force $\vec{F}$ act on a particle. If the point of application of the force undergoes a displacement $\vec{d}$ (from point $A$ to point $B$), the \emph{work} done by the force is the scalar quantity
\[
W = \vec{F} \cdot \vec{d} = F\,d\cos\theta,
\]
where $F = \|\vec{F}\|$, $d = \|\vec{d}\| = AB$ and $\theta$ is the angle between the directions of $\vec{F}$ and $\vec{d}$ ($0 \le \theta \le \pi$). The SI unit of work is the \emph{joule} (J): $1\ \mathrm{J} = 1\ \mathrm{N\,m}$.
\end{definition}

\begin{propriete}[Sign of the work]
\begin{itemize}
\item $W > 0$ (driving work) if $0 \le \theta < 90^{\circ}$: the force has a component in the direction of motion.
\item $W = 0$ if $\theta = 90^{\circ}$: the force is perpendicular to the displacement.
\item $W < 0$ (resisting work) if $90^{\circ} < \theta \le 180^{\circ}$: the force has a component opposite to the motion.
\end{itemize}
Work is a scalar; the total work of several forces is the algebraic sum of their individual works.
\end{propriete}

\begin{methode}[Computing the work of a constant force]
\begin{enumerate}
\item Draw a clear diagram showing the force $\vec{F}$, the displacement $\vec{d}$ (or the path), and mark the angle $\theta$ between them.
\item Identify the magnitude $F$ of the force and the distance $d$ travelled.
\item Apply $W = F d \cos\theta$.
\item Decide the sign from the value of $\cos\theta$ (or from the physical sense: driving or resisting).
\item For the weight $\vec{W} = m\vec{g}$: if the body rises vertically by $h$, the work done \emph{by} the weight is $W_{\text{weight}} = -mgh$; the work done \emph{against} gravity is $W_{\text{against}} = +mgh$. This depends only on the vertical height change, not on the path.
\item State the answer in joules with its sign and a brief physical interpretation.
\end{enumerate}
\end{methode}

\begin{popfigure}
\begin{tikzpicture}[scale=0.9, >=stealth]
% Force and displacement vectors
\draw[->, thick, popblue] (0,0) -- (3,0) node[midway, below] {$\vec{d}$};
\draw[->, thick, poporange] (0,0) -- (2.5,1.5) node[midway, above left] {$\vec{F}$};
\draw[dashed] (2.5,1.5) -- (2.5,0);
\draw (0.8,0) arc (0:31:0.8) node[midway, right] {$\theta$};
% Components
\draw[->, thick, popgreen] (0,0) -- (2.5,0) node[midway, below] {$F\cos\theta$};
\draw[->, thick, poppink] (2.5,0) -- (2.5,1.5) node[midway, right] {$F\sin\theta$};
\end{tikzpicture}
\legende{Work of a constant force: only the component $F\cos\theta$ along the displacement contributes.}
\end{popfigure}

\begin{exemplebox}[Work on an inclined plane]
A block of mass $5\ \mathrm{kg}$ is pulled a distance of $8\ \mathrm{m}$ up a rough plane inclined at $30^{\circ}$ to the horizontal. The pulling force of $50\ \mathrm{N}$ acts parallel to the plane, up the slope. The coefficient of friction is $0.2$. Take $g = 10\ \mathrm{m\,s^{-2}}$. Find the work done by (a) the pulling force, (b) the weight, (c) the friction, (d) the normal reaction, (e) the total work.
\tcblower
\textbf{Diagram and data:} $d = 8\ \mathrm{m}$, $\alpha = 30^{\circ}$, $F = 50\ \mathrm{N}$ (up the plane), $mg = 50\ \mathrm{N}$, $R = \mu N = \mu mg\cos\alpha = 0.2 \times 50 \times \cos30^{\circ} = 10 \times \frac{\sqrt{3}}{2} = 5\sqrt{3}\ \mathrm{N}$.

\textbf{(a) Pulling force:} $\theta = 0^{\circ}$, $W_F = 50 \times 8 \times \cos0^{\circ} = 400\ \mathrm{J}$.

\textbf{(b) Weight:} The weight $mg$ acts vertically downwards. The angle between $\vec{mg}$ and the displacement (up the plane) is $90^{\circ}+\alpha = 120^{\circ}$.
\[
W_{mg} = mg\,d\cos(90^{\circ}+\alpha) = 50 \times 8 \times \cos120^{\circ} = 400 \times \left(-\frac{1}{2}\right) = -200\ \mathrm{J}.
\]
\textit{Check:} vertical rise $h = d\sin\alpha = 8 \times \frac{1}{2} = 4\ \mathrm{m}$, so work against gravity $= mgh = 5 \times 10 \times 4 = 200\ \mathrm{J}$; work by gravity $= -200\ \mathrm{J}$. Consistent.

\textbf{(c) Friction:} Acts down the plane, opposite to displacement: $\theta = 180^{\circ}$.
\[
W_f = R d \cos180^{\circ} = 5\sqrt{3} \times 8 \times (-1) = -40\sqrt{3} \approx -69.3\ \mathrm{J}.
\]

\textbf{(d) Normal reaction:} Perpendicular to displacement: $\theta = 90^{\circ}$, $W_N = 0$.

\textbf{(e) Total work:}
\[
W_{\text{total}} = 400 - 200 - 40\sqrt{3} + 0 = 200 - 40\sqrt{3} \approx 130.7\ \mathrm{J}.
\]
The positive total work means the block gains kinetic energy.
\end{exemplebox}

\begin{attention}[Choosing the angle $\theta$]
The most frequent error is confusing the angle of the slope with the angle between the \emph{force} and the \emph{displacement}. Always draw the two vectors tail-to-tail and measure the angle between them. For the weight on an incline, $\theta = 90^{\circ} + \alpha$ (if moving up) or $\theta = 90^{\circ} - \alpha$ (if moving down). Remember: work against gravity depends \emph{only} on the vertical height gained, never on the length or shape of the path.
\end{attention}

\courssub{Work Done by Variable Forces}

\begin{definition}[Work of a variable force along a straight line]
If a force $F(x)$ varies with position $x$ and acts along the $x$-axis (or along a straight line taken as the $x$-axis), the work done in moving from $x = a$ to $x = b$ is
\[
W = \int_{a}^{b} F(x)\,dx.
\]
Geometrically, $W$ is the signed area between the graph of $F(x)$ and the $x$-axis from $a$ to $b$.
\end{definition}

\begin{methode}[Work of a variable force by integration]
\begin{enumerate}
\item Choose an origin and a positive direction along the line of motion. Express the force as a function $F(x)$ of the position $x$.
\item Write the elementary work $\delta W = F(x)\,\delta x$ for a small displacement $\delta x$.
\item Integrate between the initial and final positions: $W = \int_{x_1}^{x_2} F(x)\,dx$.
\item For an elastic string or spring obeying Hooke's law: tension $T = \dfrac{\lambda}{l}x$, where $\lambda$ is the modulus of elasticity, $l$ the natural length, and $x$ the extension. The work done in stretching from extension $x_1$ to $x_2$ is
\[
W = \int_{x_1}^{x_2} \frac{\lambda}{l}x\,dx = \frac{\lambda}{2l}(x_2^2 - x_1^2).
\]
This work is stored as elastic potential energy $E_{\text{el}} = \dfrac{\lambda x^2}{2l}$.
\item Check the sign: if the force acts in the direction of increasing $x$, the work is positive; if opposite, negative.
\end{enumerate}
\end{methode}

\begin{popfigure}
\begin{tikzpicture}[scale=0.8, >=stealth]
\begin{axis}[
    axis lines=middle,
    xlabel=$x$, ylabel=$F(x)$,
    xmin=0, xmax=5,
    ymin=0, ymax=4,
    xtick={1,3}, xticklabels={$x_1$,$x_2$},
    ytick=\empty,
    domain=0:5,
    samples=100,
    width=10cm, height=5cm,
]
\addplot[popblue, thick] {0.5*x + 1};
\addplot[poporange, fill=poporange, opacity=0.3] coordinates {(1,0) (1,1.5) (3,2.5) (3,0)} \closedcycle;
\draw[dashed] (1,0) node[below] {$x_1$} -- (1,1.5);
\draw[dashed] (3,0) node[below] {$x_2$} -- (3,2.5);
\end{axis}
\end{tikzpicture}
\legende{Work of a variable force $F(x)$ from $x_1$ to $x_2$ equals the area under the curve.}
\end{popfigure}

\begin{exoresolu}[Stretching an elastic string]
An elastic string of natural length $0.8\ \mathrm{m}$ and modulus of elasticity $60\ \mathrm{N}$ is fixed at one end. Find the work done in stretching it slowly from a length of $1.0\ \mathrm{m}$ to a length of $1.4\ \mathrm{m}$.
\tcblower
\textbf{Step 1 — Hooke's law.} Natural length $l = 0.8\ \mathrm{m}$, modulus $\lambda = 60\ \mathrm{N}$. Tension at extension $x$:
\[
F(x) = T(x) = \frac{\lambda}{l}x = \frac{60}{0.8}x = 75x\ \mathrm{N}.
\]

\textbf{Step 2 — Limits of integration.} Length $1.0\ \mathrm{m} \Rightarrow$ extension $x_1 = 1.0 - 0.8 = 0.2\ \mathrm{m}$. Length $1.4\ \mathrm{m} \Rightarrow$ extension $x_2 = 1.4 - 0.8 = 0.6\ \mathrm{m}$.

\textbf{Step 3 — Integrate.}
\[
W = \int_{0.2}^{0.6} 75x\,dx = 75\left[\frac{x^2}{2}\right]_{0.2}^{0.6} = \frac{75}{2}\left(0.36 - 0.04\right) = 37.5 \times 0.32 = 12\ \mathrm{J}.
\]

\textbf{Step 4 — Check via elastic potential energy.}
\[
E_{\text{el}} = \frac{\lambda x^2}{2l} \quad\Rightarrow\quad \Delta E_{\text{el}} = \frac{60}{2(0.8)}(0.6^2 - 0.2^2) = 37.5 \times 0.32 = 12\ \mathrm{J}.
\]
The work done equals the gain in elastic potential energy, as expected for a slow (quasi-static) stretch with no kinetic energy change.
\end{exoresolu}

\begin{exoresolu}[Work from a force–displacement graph]
A force $F(x)$ acts on a body moving along the $x$-axis. The graph of $F$ against $x$ consists of: a straight line from $(0,0)$ to $(2,10)$, a horizontal segment from $(2,10)$ to $(4,10)$, and a straight line from $(4,10)$ to $(6,0)$. Find the total work done from $x=0$ to $x=6$.
\tcblower
The work is the area under the graph. We compute it as the sum of three areas:
\begin{itemize}
\item Triangle $0 \le x \le 2$: $\frac{1}{2} \times 2 \times 10 = 10\ \mathrm{J}$.
\item Rectangle $2 \le x \le 4$: $2 \times 10 = 20\ \mathrm{J}$.
\item Triangle $4 \le x \le 6$: $\frac{1}{2} \times 2 \times 10 = 10\ \mathrm{J}$.
\end{itemize}
Total work $W = 10 + 20 + 10 = 40\ \mathrm{J}$.
\end{exoresolu}

\courssub{Energy: Kinetic and Potential}

\begin{definition}[Kinetic energy]
The \emph{kinetic energy} of a body of mass $m$ moving with speed $v$ is
\[
E_k = \frac{1}{2}mv^2.
\]
It is a scalar, always non-negative, measured in joules (J). It depends only on the speed, not on the direction of motion.
\end{definition}

\begin{definition}[Gravitational potential energy]
Near the Earth's surface, the \emph{gravitational potential energy} of a body of mass $m$ at a height $h$ above a chosen reference level is
\[
E_p = mgh.
\]
Only \emph{changes} in potential energy are physically meaningful; the reference level (where $E_p = 0$) can be chosen arbitrarily for convenience.
\end{definition}

\begin{definition}[Elastic potential energy]
For an elastic string or spring of natural length $l$ and modulus $\lambda$, stretched (or compressed) by an amount $x$, the elastic potential energy stored is
\[
E_{\text{el}} = \frac{\lambda x^2}{2l}.
\]
\end{definition}

\begin{propriete}[Conservative forces]
A force is \emph{conservative} if the work it does depends only on the initial and final positions, not on the path taken. Gravity and elastic forces are conservative. For a conservative force, one can define a potential energy $E_p$ such that the work done by the force equals the \emph{negative} change in potential energy: $W = -\Delta E_p$. Friction and air resistance are \emph{non-conservative} (dissipative) forces.
\end{propriete}

\begin{aretenir}[Key energy formulas]
\begin{center}
\begin{tabularx}{\linewidth}{|l|X|}\hline
\textbf{Quantity} & \textbf{Formula} \\ \hline
Kinetic energy & $E_k = \frac{1}{2}mv^2$ \\ \hline
Gravitational potential energy & $E_p = mgh$ (relative to a reference level) \\ \hline
Elastic potential energy & $E_{\text{el}} = \dfrac{\lambda x^2}{2l}$ \\ \hline
Work by gravity (vertical rise $h$) & $W_{\text{grav}} = -mgh$ \\ \hline
Work against gravity & $W_{\text{against}} = +mgh$ \\ \hline
Work by elastic force (extension $x_1 \to x_2$) & $W = -\dfrac{\lambda}{2l}(x_2^2 - x_1^2)$ \\ \hline
\end{tabularx}
\end{center}
\end{aretenir}

\courssub{The Work–Energy Theorem and Conservation of Energy}

\begin{propriete}[Work–Energy Theorem]
The total work done by \emph{all} forces acting on a particle (or a rigid body in translation) equals the change in its kinetic energy:
\[
W_{\text{total}} = \Delta E_k = \frac{1}{2}mv^2 - \frac{1}{2}mu^2.
\]
This theorem holds whether the forces are conservative or not. It is a scalar equation, often simpler than using Newton's second law when only speeds and distances are involved.
\end{propriete}

\begin{methode}[Applying the work–energy theorem]
\begin{enumerate}
\item Identify the body, its mass $m$, initial speed $u$ and final speed $v$.
\item Compute $\Delta E_k = \frac{1}{2}m(v^2 - u^2)$.
\item List \emph{all} forces that do work on the body (weight, tension, friction, driving force, normal reaction does no work if perpendicular to motion, etc.).
\item Compute the work done by each force over the displacement. For constant forces: $W = Fd\cos\theta$. For variable forces: integrate.
\item Sum the works algebraically to get $W_{\text{total}}$.
\item Set $W_{\text{total}} = \Delta E_k$ and solve for the unknown.
\end{enumerate}
\end{methode}

\begin{exoresolu}[Lorry braking downhill]
A lorry of mass $3000\ \mathrm{kg}$ travels at $20\ \mathrm{m\,s^{-1}}$ down a hill inclined at $\sin^{-1}(0.1)$ to the horizontal. The brakes apply a constant resisting force of $9000\ \mathrm{N}$. Using the work–energy theorem, find the distance travelled before the speed falls to $8\ \mathrm{m\,s^{-1}}$. Take $g = 10\ \mathrm{m\,s^{-2}}$.
\tcblower
\textbf{Step 1 — Change in kinetic energy.}
\[
\Delta E_k = \frac{1}{2}(3000)\left(8^2 - 20^2\right) = 1500(64 - 400) = -504\,000\ \mathrm{J}.
\]

\textbf{Step 2 — Works over distance $d$.} 
Component of weight down the slope: $mg\sin\alpha = 3000 \times 10 \times 0.1 = 3000\ \mathrm{N}$ (driving).
\[
W_{\text{weight}} = +3000\,d.
\]
Braking force opposes motion:
\[
W_{\text{brake}} = -9000\,d.
\]
Normal reaction does no work.

\textbf{Step 3 — Apply the theorem.}
\begin{align*}
W_{\text{total}} &= \Delta E_k\\
3000d - 9000d &= -504\,000\\
-6000d &= -504\,000\\
d &= 84\ \mathrm{m}.
\end{align*}
The lorry travels $84\ \mathrm{m}$ before its speed drops to $8\ \mathrm{m\,s^{-1}}$. Note that gravity \emph{supplies} energy (positive work), so the brakes must dissipate both the initial kinetic energy and the work done by gravity.
\end{exoresolu}

\begin{propriete}[Conservation of mechanical energy]
If \emph{only conservative forces} (gravity, elastic forces) do work on a system, the total mechanical energy $E = E_k + E_p$ (including elastic potential energy if present) remains constant:
\[
E_{\text{initial}} = E_{\text{final}} \quad\Longleftrightarrow\quad \frac{1}{2}mu^2 + mgh_1 + \frac{\lambda x_1^2}{2l} = \frac{1}{2}mv^2 + mgh_2 + \frac{\lambda x_2^2}{2l}.
\]
If non-conservative forces (friction, resistance, driving force) do work, mechanical energy is not conserved; use the work–energy theorem instead, or write
\[
E_{\text{final}} = E_{\text{initial}} + W_{\text{non-conservative}}.
\]
\end{propriete}

\begin{methode}[Applying conservation of mechanical energy]
\begin{enumerate}
\item Verify that no non-conservative forces do work (smooth surfaces, no engine, no resistance). If they do, use the work–energy theorem.
\item Choose a reference level for gravitational potential energy (usually the lowest point).
\item Write the total mechanical energy at the initial state: $E_i = \frac{1}{2}mu^2 + mgh_i (+ E_{\text{el},i})$.
\item Write it at the final state: $E_f = \frac{1}{2}mv^2 + mgh_f (+ E_{\text{el},f})$.
\item Set $E_i = E_f$ and solve. The mass $m$ often cancels.
\item If a constant resistance $R$ acts over a distance $s$, the energy lost is $Rs$: $E_i - E_f = Rs$.
\end{enumerate}
\end{methode}

\begin{exoresolu}[Bead on a smooth vertical wire]
A small bead of mass $0.5\ \mathrm{kg}$ slides from rest at the top $A$ of a smooth curved wire. Point $A$ is $5\ \mathrm{m}$ vertically above the lowest point $B$, and point $C$ is $2\ \mathrm{m}$ above $B$. Find the speed of the bead (a) at $B$, (b) at $C$. Take $g = 10\ \mathrm{m\,s^{-2}}$.
\tcblower
The wire is smooth $\Rightarrow$ reaction does no work $\Rightarrow$ mechanical energy conserved. Choose $B$ as reference level ($h=0$).

\textbf{(a) Speed at $B$.} At $A$: $v_A = 0$, $h_A = 5\ \mathrm{m}$. $E_A = 0 + mg(5) = 5mg$.
At $B$: $h_B = 0$, $E_B = \frac{1}{2}mv_B^2$.
\[
\frac{1}{2}mv_B^2 = 5mg \;\Longrightarrow\; v_B^2 = 10g = 100 \;\Longrightarrow\; v_B = 10\ \mathrm{m\,s^{-1}}.
\]

\textbf{(b) Speed at $C$.} $h_C = 2\ \mathrm{m}$, $E_C = \frac{1}{2}mv_C^2 + mg(2)$.
\[
\frac{1}{2}mv_C^2 + 2mg = 5mg \;\Longrightarrow\; \frac{1}{2}v_C^2 = 3g = 30 \;\Longrightarrow\; v_C = \sqrt{60} = 2\sqrt{15} \approx 7.75\ \mathrm{m\,s^{-1}}.
\]

\textbf{Remark.} The mass cancels: any bead would have the same speeds. The shape of the wire is irrelevant — only the vertical heights matter because gravity is conservative.
\end{exoresolu}

\begin{exoresolu}[Block on a rough inclined plane with a spring]
A block of mass $2\ \mathrm{kg}$ is placed on a rough plane inclined at $30^{\circ}$. It is attached to a spring of natural length $0.5\ \mathrm{m}$ and modulus $40\ \mathrm{N}$, fixed at the top of the plane. The block is released from rest with the spring unstretched. The coefficient of friction is $0.25$. Find the maximum distance the block moves down the plane before coming to rest. Take $g = 10\ \mathrm{m\,s^{-2}}$.
\tcblower
Let $x$ be the distance moved down the plane (extension of the spring). At the lowest point, $v=0$.
Initial energy (top): $E_i = 0$ (KE) $+ 0$ (GPE, ref at start) $+ 0$ (spring).
Final energy (lowest point): 
- KE $= 0$.
- GPE: vertical drop $= x\sin30^{\circ} = x/2$, so $E_p = -mg(x/2) = -2 \times 10 \times x/2 = -10x$.
- Spring energy: $E_{\text{el}} = \frac{\lambda x^2}{2l} = \frac{40 x^2}{2 \times 0.5} = 40x^2$.
Work done against friction: friction force $F_f = \mu N = \mu mg\cos30^{\circ} = 0.25 \times 20 \times \frac{\sqrt{3}}{2} = 2.5\sqrt{3}\ \mathrm{N}$. Work lost $= F_f x = 2.5\sqrt{3}\,x$.

Energy equation: $E_i = E_f + \text{work against friction}$.
\[
0 = -10x + 40x^2 + 2.5\sqrt{3}\,x.
\]
Divide by $x$ ($x>0$): $40x = 10 - 2.5\sqrt{3} \approx 10 - 4.33 = 5.67$.
\[
x = \frac{10 - 2.5\sqrt{3}}{40} \approx 0.142\ \mathrm{m}.
\]
The block moves about $14.2\ \mathrm{cm}$ down the plane.
\end{exoresolu}

\courssub{Power: Definition, Instantaneous Power and Power–Velocity}

\begin{definition}[Power]
\emph{Power} is the rate at which work is done or energy is transferred.
\begin{itemize}
\item \textbf{Average power} over a time interval $\Delta t$: $P_{\text{av}} = \dfrac{W}{\Delta t} = \dfrac{\Delta E}{\Delta t}$.
\item \textbf{Instantaneous power} of a force $\vec{F}$ acting on a particle moving with velocity $\vec{v}$:
\[
P = \vec{F} \cdot \vec{v} = Fv\cos\theta,
\]
where $\theta$ is the angle between $\vec{F}$ and $\vec{v}$.
\end{itemize}
The SI unit is the \emph{watt} (W): $1\ \mathrm{W} = 1\ \mathrm{J\,s^{-1}}$. $1\ \mathrm{kW} = 1000\ \mathrm{W}$; $1\ \mathrm{hp} \approx 746\ \mathrm{W}$.
\end{definition}

\begin{propriete}[Power and velocity for a vehicle]
If a vehicle's engine works at a constant power $P$ and the vehicle moves at speed $v$, the driving force (tractive force) is
\[
F = \frac{P}{v}.
\]
At \emph{maximum speed} on a given incline, acceleration $a = 0$, so the driving force exactly balances the total resistance (including the component of weight if on a slope):
\[
\frac{P}{v_{\max}} = R_{\text{total}} \quad\Longrightarrow\quad v_{\max} = \frac{P}{R_{\text{total}}}.
\]
If the vehicle is accelerating, Newton's second law gives
\[
\frac{P}{v} - R_{\text{total}} = ma.
\]
\end{propriete}

\begin{methode}[Solving power–velocity–acceleration problems]
\begin{enumerate}
\item Identify the given power $P$ (constant or instantaneous), mass $m$, resistances $R$ (may depend on $v$), and incline $\alpha$.
\item Write the driving force: $F = \dfrac{P}{v}$.
\item Apply Newton's second law along the direction of motion:
\[
\frac{P}{v} - R - mg\sin\alpha = ma.
\]
\item For maximum speed: set $a = 0$ and solve for $v_{\max}$.
\item For acceleration at a given speed $v$: substitute $v$ and solve for $a$.
\item If resistance varies with speed (e.g. $R = kv$ or $kv^2$), substitute the expression before solving.
\end{enumerate}
\end{methode}

\begin{exoresolu}[Car on a hill and on level ground]
A car of mass $1200\ \mathrm{kg}$ climbs a hill inclined at $\alpha$ where $\sin\alpha = \frac{1}{14}$, at a constant speed of $18\ \mathrm{m\,s^{-1}}$. Non-gravitational resistances total $400\ \mathrm{N}$. (a) Find the power developed by the engine. (b) Find the maximum speed of the car on level ground with the same power and the same resistance. Take $g = 10\ \mathrm{m\,s^{-2}}$.
\tcblower
\textbf{(a) Power on the hill.} Constant speed $\Rightarrow a=0 \Rightarrow$ driving force $F$ balances resistance + weight component.
\[
F = R + mg\sin\alpha = 400 + 1200 \times 10 \times \frac{1}{14} = 400 + \frac{12000}{14} = 400 + \frac{6000}{7} = \frac{8800}{7}\ \mathrm{N}.
\]
\[
P = Fv = \frac{8800}{7} \times 18 = \frac{158400}{7} \approx 22\,629\ \mathrm{W} = 22.6\ \mathrm{kW}.
\]

\textbf{(b) Maximum speed on level ground.} On level ground, $\sin\alpha = 0$, so at $v_{\max}$:
\[
\frac{P}{v_{\max}} = R = 400 \;\Longrightarrow\; v_{\max} = \frac{P}{400} = \frac{158400/7}{400} = \frac{396}{7} \approx 56.6\ \mathrm{m\,s^{-1}}.
\]

\textbf{Interpretation.} On the flat, the same engine power yields a much higher top speed because no power is spent raising the car's gravitational potential energy.
\end{exoresolu}

\begin{exoresolu}[Acceleration of a car with constant power]
A car of mass $1000\ \mathrm{kg}$ moves on a horizontal road against a constant resistance of $500\ \mathrm{N}$. The engine works at a constant rate of $30\ \mathrm{kW}$. Find (a) the maximum speed, (b) the acceleration when the speed is $20\ \mathrm{m\,s^{-1}}$.
\tcblower
\textbf{(a) Maximum speed.} $a=0 \Rightarrow F = R = 500\ \mathrm{N}$.
\[
v_{\max} = \frac{P}{F} = \frac{30\,000}{500} = 60\ \mathrm{m\,s^{-1}}.
\]

\textbf{(b) Acceleration at $v = 20\ \mathrm{m\,s^{-1}}$.}
Driving force: $F = \frac{P}{v} = \frac{30\,000}{20} = 1500\ \mathrm{N}$.
Newton's second law: $F - R = ma$.
\[
1500 - 500 = 1000\,a \;\Longrightarrow\; a = 1.0\ \mathrm{m\,s^{-2}}.
\]

\textbf{Remark.} As $v$ increases, $F = P/v$ decreases, so acceleration decreases until it reaches zero at $v_{\max} = 60\ \mathrm{m\,s^{-1}}$.
\end{exoresolu}

\begin{attention}[Constant speed vs. acceleration]
A classic GCE error: using $P = Rv$ (which assumes $a=0$) when the vehicle is \emph{accelerating}. If the speed is changing, you \emph{must} write $F = P/v$ and then $F - R - mg\sin\alpha = ma$. Only at maximum speed (or constant speed) does $F = R + mg\sin\alpha$ hold.
\end{attention}

\courssub{Pumping Water, Water Jets and Applications of Work, Energy and Power}

\begin{definition}[Power of a pump or jet]
A pump delivers a mass of water $\dfrac{dm}{dt}$ (in $\mathrm{kg\,s^{-1}}$) and gives it:
\begin{itemize}
\item Gravitational potential energy by raising it through a height $h$: power $P_p = \dfrac{dm}{dt}\,gh$.
\item Kinetic energy by ejecting it at speed $v$: power $P_k = \dfrac{1}{2}\dfrac{dm}{dt}\,v^2$.
\end{itemize}
The total useful hydraulic power is
\[
P_{\text{useful}} = \frac{dm}{dt}\left(gh + \frac{1}{2}v^2\right).
\]
If the pump has efficiency $\eta$ ($0 < \eta \le 1$), the input power required is
\[
P_{\text{in}} = \frac{P_{\text{useful}}}{\eta}.
\]
For a jet of cross-sectional area $A$ and speed $v$, the mass flow rate is $\dfrac{dm}{dt} = \rho A v$, where $\rho = 1000\ \mathrm{kg\,m^{-3}}$ for water.
\end{definition}

\begin{methode}[Power of a pump or water jet]
\begin{enumerate}
\item Determine the mass flow rate $\dfrac{dm}{dt}$ (given directly, or computed from volume flow rate $\times 1000$, or from $\rho A v$ for a jet).
\item Compute the potential energy rate: $P_p = \dfrac{dm}{dt}\,gh$ (if water is lifted).
\item Compute the kinetic energy rate: $P_k = \frac{1}{2}\dfrac{dm}{dt}\,v^2$ (if water is ejected with speed $v$).
\item Add them: $P_{\text{useful}} = P_p + P_k$.
\item If efficiency $\eta$ is given, compute input power: $P_{\text{in}} = P_{\text{useful}}/\eta$.
\item Check units: $\mathrm{kg\,s^{-1}} \times \mathrm{m\,s^{-2}} \times \mathrm{m} = \mathrm{W}$.
\end{enumerate}
\end{methode}

\begin{popfigure}
\begin{tikzpicture}[scale=0.7, >=stealth]
% Well
\draw[thick, popblue] (-1.5,-4) -- (1.5,-4) -- (1.5,0) -- (-1.5,0) -- cycle;
\draw[popblue, fill=popblue!20] (-1.5,-4) rectangle (1.5,-3.5);
\node at (0,-3.75) {Water level};
% Pipe
\draw[thick, popdark] (0,-4) -- (0,2);
% Pump
\draw[thick, poporange, fill=poporange!20] (-1,2) rectangle (1,3);
\node at (0,2.5) {Pump};
% Tank
\draw[thick, popgreen] (2,3) -- (5,3) -- (5,5) -- (2,5) -- cycle;
\draw[popgreen, fill=popgreen!20] (2,5) rectangle (5,4.5);
\node at (3.5,4.75) {Tank};
% Jet
\draw[->, thick, popblue] (0,3) -- (2,3) node[midway, above] {$v$};
% Dimensions
\draw[<->, thin] (-2,-4) -- (-2,3) node[midway, left] {$h$};
\draw[<->, thin] (0,3) -- (2,3) node[midway, above] {pipe};
\end{tikzpicture}
\legende{Pump raising water from a well to an overhead tank with ejection speed $v$.}
\end{popfigure}

\begin{exoresolu}[Village water pump]
A pump raises water from a well $12\ \mathrm{m}$ deep and delivers it at $4\ \mathrm{m\,s^{-1}}$ into an overhead tank at a rate of $20\ \mathrm{kg}$ per second. (a) Find the useful power of the pump. (b) If the pump is only $70\%$ efficient, find the power that must be supplied to it. Take $g = 10\ \mathrm{m\,s^{-2}}$.
\tcblower
\textbf{Given:} $\dfrac{dm}{dt} = 20\ \mathrm{kg\,s^{-1}}$, $h = 12\ \mathrm{m}$, $v = 4\ \mathrm{m\,s^{-1}}$, $\eta = 0.70$.

\textbf{(a) Useful power.}
Potential energy rate:
\[
P_p = \frac{dm}{dt}\,gh = 20 \times 10 \times 12 = 2400\ \mathrm{W}.
\]
Kinetic energy rate:
\[
P_k = \frac{1}{2}\frac{dm}{dt}\,v^2 = \frac{1}{2} \times 20 \times 16 = 160\ \mathrm{W}.
\]
Total useful power:
\[
P_{\text{useful}} = 2400 + 160 = 2560\ \mathrm{W} = 2.56\ \mathrm{kW}.
\]

\textbf{(b) Input power.}
\[
P_{\text{in}} = \frac{P_{\text{useful}}}{\eta} = \frac{2560}{0.70} \approx 3657\ \mathrm{W} = 3.66\ \mathrm{kW}.
\]

\textbf{Remark.} The kinetic energy contribution ($160\ \mathrm{W}$) is small compared to the potential energy contribution ($2400\ \mathrm{W}$): most power is spent lifting the water.
\end{exoresolu}

\begin{exoresolu}[Water jet from a fire hose]
A fire hose has a nozzle of cross-sectional area $2.5 \times 10^{-4}\ \mathrm{m^2}$. Water emerges at $25\ \mathrm{m\,s^{-1}}$. Find (a) the mass of water leaving the nozzle per second, (b) the kinetic energy given to the water per second, (c) the force exerted on a vertical wall if the jet hits it perpendicularly and falls vertically (assuming no splash back). Take $\rho = 1000\ \mathrm{kg\,m^{-3}}$.
\tcblower
\textbf{(a) Mass flow rate:}
\[
\frac{dm}{dt} = \rho A v = 1000 \times 2.5 \times 10^{-4} \times 25 = 6.25\ \mathrm{kg\,s^{-1}}.
\]

\textbf{(b) Kinetic energy per second (power of the jet):}
\[
P_k = \frac{1}{2}\frac{dm}{dt}v^2 = \frac{1}{2} \times 6.25 \times 625 = 1953.125\ \mathrm{W} \approx 1.95\ \mathrm{kW}.
\]

\textbf{(c) Force on the wall.} The water's horizontal momentum is destroyed. Rate of change of momentum = force.
Initial horizontal momentum per second $= \dfrac{dm}{dt} \times v = 6.25 \times 25 = 156.25\ \mathrm{kg\,m\,s^{-2}} = 156.25\ \mathrm{N}$.
Final horizontal momentum $= 0$. Force on water $= -156.25\ \mathrm{N}$; by Newton's third law, force on wall $= 156.25\ \mathrm{N}$.
\end{exoresolu}

\courssub{Mixed Applications of Work, Energy and Power}

\begin{methode}[Strategy for GCE-type mixed problems]
\begin{enumerate}
\item Read the problem twice; draw a clear diagram with all forces, heights, speeds, and distances.
\item Decide the appropriate tool:
\begin{itemize}
\item Conservation of energy: smooth surfaces, no engine, no resistance.
\item Work–energy theorem: resistance, driving force, or non-conservative forces present.
\item $P = Fv$ with Newton's second law: engine power given, possibly with acceleration.
\end{itemize}
\item Write one energy or force equation per stage of the motion.
\item Keep exact values (surds, fractions) as long as possible; round only at the final answer to 3 significant figures.
\item Check physical sense: positive power, speed $\le v_{\max}$, energy lost to resistance, etc.
\end{enumerate}
\end{methode}

\begin{exoresolu}[Cyclist on the Buea mountain road]
A cyclist and her bicycle have total mass $80\ \mathrm{kg}$. She descends a straight hill inclined at $\sin^{-1}\!\left(\frac{1}{20}\right)$ to the horizontal. Starting from rest, she freewheels (no pedalling) against a constant resistance of $30\ \mathrm{N}$. (a) Use the work–energy theorem to find her speed after travelling $200\ \mathrm{m}$ down the slope. (b) She then reaches level ground and pedals, working at a constant rate of $180\ \mathrm{W}$ against the same resistance of $30\ \mathrm{N}$. Find her maximum speed on the level. Take $g = 10\ \mathrm{m\,s^{-2}}$.
\tcblower
\textbf{(a) Speed after $200\ \mathrm{m}$ downhill.}
Vertical drop: $h = 200\sin\alpha = 200 \times \frac{1}{20} = 10\ \mathrm{m}$.
Forces doing work along the slope:
- Component of gravity down the slope: $mg\sin\alpha = 80 \times 10 \times \frac{1}{20} = 40\ \mathrm{N}$ (driving).
- Resistance: $30\ \mathrm{N}$ (opposing).
Work–energy theorem ($u=0$):
\[
W_{\text{gravity}} + W_{\text{resistance}} = \frac{1}{2}mv^2.
\]
\[
40 \times 200 - 30 \times 200 = \frac{1}{2}(80)v^2.
\]
\[
8000 - 6000 = 40v^2 \;\Longrightarrow\; v^2 = 50 \;\Longrightarrow\; v = \sqrt{50} = 5\sqrt{2} \approx 7.07\ \mathrm{m\,s^{-1}}.
\]

\textbf{(b) Maximum speed on level ground.}
At maximum speed, $a=0$, so driving force $F = R = 30\ \mathrm{N}$.
Power $P = 180\ \mathrm{W}$ constant.
\[
v_{\max} = \frac{P}{F} = \frac{180}{30} = 6\ \mathrm{m\,s^{-1}}.
\]

\textbf{Interpretation.} Freewheeling downhill she reaches $7.07\ \mathrm{m\,s^{-1}}$, faster than the $6\ \mathrm{m\,s^{-1}}$ she can sustain on the flat with $180\ \mathrm{W}$. During the descent, gravity supplied power $mgv\sin\alpha = 40 \times 7.07 \approx 283\ \mathrm{W}$, well above her own $180\ \mathrm{W}$.
\end{exoresolu}

\begin{exoresolu}[Particle projected up a rough inclined plane]
A particle of mass $0.5\ \mathrm{kg}$ is projected up a rough plane inclined at $30^{\circ}$ with an initial speed of $10\ \mathrm{m\,s^{-1}}$. The coefficient of friction is $0.2$. Find (a) the distance travelled up the plane before coming to rest, (b) the speed when it returns to the starting point. Take $g = 10\ \mathrm{m\,s^{-2}}$.
\tcblower
\textbf{Data:} $m=0.5\ \mathrm{kg}$, $u=10\ \mathrm{m\,s^{-1}}$, $\alpha=30^{\circ}$, $\mu=0.2$, $g=10$.
Weight component down plane: $mg\sin\alpha = 0.5 \times 10 \times 0.5 = 2.5\ \mathrm{N}$.
Normal reaction: $N = mg\cos\alpha = 0.5 \times 10 \times \frac{\sqrt{3}}{2} = 2.5\sqrt{3}\ \mathrm{N}$.
Friction: $F_f = \mu N = 0.2 \times 2.5\sqrt{3} = 0.5\sqrt{3}\ \mathrm{N}$.

\textbf{(a) Distance $x$ up the plane to rest ($v=0$).}
Work–energy theorem (upward displacement $x$):
Forces doing work: weight component ($-2.5x$), friction ($-0.5\sqrt{3}\,x$).
\[
\Delta E_k = 0 - \frac{1}{2}(0.5)(10^2) = -25\ \mathrm{J}.
\]
\[
-2.5x - 0.5\sqrt{3}\,x = -25 \;\Longrightarrow\; x(2.5 + 0.5\sqrt{3}) = 25.
\]
\[
x = \frac{25}{2.5 + 0.5\sqrt{3}} = \frac{50}{5 + \sqrt{3}} = \frac{50(5-\sqrt{3})}{22} = \frac{25(5-\sqrt{3})}{11} \approx 7.42\ \mathrm{m}.
\]

\textbf{(b) Speed on return to start.}
Total distance travelled $= 2x$. Work done by friction over round trip $= -2F_f x = -2(0.5\sqrt{3})x = -\sqrt{3}\,x$.
Gravity is conservative $\Rightarrow$ net work by gravity over round trip $= 0$.
Work–energy theorem for round trip:
\[
W_{\text{friction}} = \Delta E_k = \frac{1}{2}mv^2 - \frac{1}{2}mu^2.
\]
\[
-\sqrt{3}\,x = \frac{1}{2}(0.5)v^2 - 25.
\]
Substitute $x = \frac{25(5-\sqrt{3})}{11}$:
\[
-\sqrt{3} \times \frac{25(5-\sqrt{3})}{11} = 0.25v^2 - 25.
\]
\[
-\frac{125\sqrt{3} - 75}{11} = 0.25v^2 - 25.
\]
\[
0.25v^2 = 25 - \frac{125\sqrt{3} - 75}{11} = \frac{275 - 125\sqrt{3} + 75}{11} = \frac{350 - 125\sqrt{3}}{11}.
\]
\[
v^2 = \frac{4(350 - 125\sqrt{3})}{11} = \frac{1400 - 500\sqrt{3}}{11} \approx 48.3 \;\Longrightarrow\; v \approx 6.95\ \mathrm{m\,s^{-1}}.
\]
The return speed is less than the initial speed because friction dissipates energy.
\end{exoresolu}

\begin{attention}[Examination checklist]
\begin{itemize}
\item Work and energy are scalars: add them with signs, never as vectors.
\item $E_k = \frac{1}{2}mv^2 \ge 0$; a negative $\Delta E_k$ means kinetic energy is lost.
\item In $P = Fv$, $v$ is the \emph{instantaneous} speed and $F$ the force \emph{in the direction of motion}.
\item For pumps/jets: never forget the $mgh$ term (lifting) or the $\frac{1}{2}mv^2$ term (ejection) when both occur.
\item Always quote units: joules (J) for work and energy, watts (W) for power.
\item On an incline, the angle for work is between the \emph{force} and the \emph{displacement along the plane}, not the angle of the plane itself.
\item If a problem gives power and asks for acceleration, use $F = P/v$ then $F - R = ma$; do \emph{not} assume $F = R$.
\end{itemize}
\end{attention}

\begin{savaistu}[Historical note: James Watt and the horsepower]
The watt is named after the Scottish engineer James Watt (1736–1819), who improved the steam engine. To market his engines, he compared their output to the work a horse could do, defining the \emph{horsepower} (hp). He found a horse could turn a mill wheel of radius $12\ \mathrm{ft}$ at $2.4$ revolutions per minute with a force of $180\ \mathrm{lbf}$, giving $P \approx 33\,000\ \mathrm{ft\,lbf/min} \approx 746\ \mathrm{W}$. In Cameroon, generator and pump ratings are often given in kW or hp — remember $1\ \mathrm{hp} \approx 0.746\ \mathrm{kW}$.
\end{savaistu}

\begin{exercice}[Practice problems]
\begin{enumerate}
\item A force $\vec{F} = (3\vec{i} + 4\vec{j})\ \mathrm{N}$ moves a particle from $(1,2)$ to $(4,6)$ (coordinates in metres). Find the work done by $\vec{F}$.
\item A $2\ \mathrm{kg}$ block is pushed $5\ \mathrm{m}$ up a rough $30^{\circ}$ incline by a horizontal force of $30\ \mathrm{N}$. The coefficient of friction is $0.15$. Find the work done by each force and the final speed if it starts from rest. ($g=10$)
\item A spring of natural length $0.4\ \mathrm{m}$ and modulus $50\ \mathrm{N}$ is compressed to $0.25\ \mathrm{m}$. Find the work done in compressing it further to $0.15\ \mathrm{m}$.
\item A $1500\ \mathrm{kg}$ car climbs a $1$ in $10$ hill at a constant $20\ \mathrm{m\,s^{-1}}$. Resistance is $500\ \mathrm{N}$. Find the engine power. If the same power is used on level ground, find the maximum speed.
\item A pump delivers $0.05\ \mathrm{m^3/s}$ of water through a hose of cross-section $2 \times 10^{-4}\ \mathrm{m^2}$ to a height of $20\ \mathrm{m}$. Find the useful power. If efficiency is $75\%$, find the input power. ($\rho=1000$)
\item A particle slides from rest down a smooth quarter-circle wire of radius $2\ \mathrm{m}$. Find its speed at the bottom. It then moves on a rough horizontal table with $\mu=0.3$. How far does it travel before stopping? ($g=10$)
\end{enumerate}
\end{exercice}

\begin{corrige}[Exercise 1]
$\vec{d} = (3\vec{i} + 4\vec{j})\ \mathrm{m}$. $W = \vec{F}\cdot\vec{d} = 3\times3 + 4\times4 = 9+16 = 25\ \mathrm{J}$.
\end{corrige}

\begin{corrige}[Exercise 2]
Forces: horizontal $F=30\ \mathrm{N}$ ($\theta=30^{\circ}$ to plane), weight $mg=20\ \mathrm{N}$, friction $f=\mu N$, normal $N$.
$N = mg\cos30^{\circ} - F\sin30^{\circ} = 20\frac{\sqrt{3}}{2} - 15 = 10\sqrt{3}-15 \approx 2.32\ \mathrm{N}$.
$f = 0.15(10\sqrt{3}-15) \approx 0.348\ \mathrm{N}$.
$W_F = 30 \times 5 \times \cos30^{\circ} = 150\frac{\sqrt{3}}{2} = 75\sqrt{3} \approx 129.9\ \mathrm{J}$.
$W_{mg} = -mg(5\sin30^{\circ}) = -20 \times 2.5 = -50\ \mathrm{J}$.
$W_f = -f \times 5 \approx -1.74\ \mathrm{J}$.
$W_N = 0$.
$W_{\text{total}} \approx 78.2\ \mathrm{J} = \frac{1}{2}(2)v^2 \Rightarrow v \approx 8.84\ \mathrm{m\,s^{-1}}$.
\end{corrige}

\begin{corrige}[Exercise 3]
$l=0.4$, $\lambda=50$. $x_1 = 0.4-0.25 = 0.15$, $x_2 = 0.4-0.15 = 0.25$.
$W = \frac{\lambda}{2l}(x_2^2 - x_1^2) = \frac{50}{0.8}(0.0625 - 0.0225) = 62.5 \times 0.04 = 2.5\ \mathrm{J}$.
\end{corrige}

\begin{corrige}[Exercise 4]
$\sin\alpha = 0.1$. $F = R + mg\sin\alpha = 500 + 15000 \times 0.1 = 2000\ \mathrm{N}$.
$P = Fv = 2000 \times 20 = 40\,000\ \mathrm{W} = 40\ \mathrm{kW}$.
Level: $v_{\max} = P/R = 40\,000/500 = 80\ \mathrm{m\,s^{-1}}$.
\end{corrige}

\begin{corrige}[Exercise 5]
$\frac{dm}{dt} = \rho \times \text{volume rate} = 1000 \times 0.05 = 50\ \mathrm{kg/s}$.
$v = \frac{\text{volume rate}}{A} = \frac{0.05}{2\times10^{-4}} = 250\ \mathrm{m/s}$.
$P_p = 50 \times 10 \times 20 = 10\,000\ \mathrm{W}$.
$P_k = \frac{1}{2} \times 50 \times 250^2 = 1\,562\,500\ \mathrm{W}$.
$P_{\text{useful}} = 1\,572\,500\ \mathrm{W} \approx 1.57\ \mathrm{MW}$.
$P_{\text{in}} = 1.5725/0.75 \approx 2.10\ \mathrm{MW}$.
\end{corrige}

\begin{corrige}[Exercise 6]
Smooth quarter-circle: $h = R = 2\ \mathrm{m}$. $v_B = \sqrt{2gh} = \sqrt{40} = 2\sqrt{10} \approx 6.32\ \mathrm{m/s}$.
Rough horizontal: $W_f = -\mu mg d = \Delta E_k = 0 - \frac{1}{2}m(40)$.
$0.3 \times 10 \times d = 20 \Rightarrow d = 20/3 \approx 6.67\ \mathrm{m}$.
\end{corrige}

\newpage

\coursec{Exercises}

\courssub{I check my knowledge}

\begin{exercice}[Multiple choice]
A force of magnitude $F$ acts on a body which moves a distance $s$ in a straight line, the force making an angle $\theta$ with the direction of motion. The work done by the force is:
\begin{enumerate}
\item $Fs$
\item $Fs\sin\theta$
\item $Fs\cos\theta$
\item $\dfrac{Fs}{\cos\theta}$
\end{enumerate}
Choose the correct answer and justify your choice.
\end{exercice}

\begin{exercice}[True or false]
For each of the following statements, say whether it is \emph{true} or \emph{false} and justify your answer.
\begin{enumerate}
\item The work done against gravity in lifting a body of mass $m$ through a vertical height $h$ is $mgh$, whatever path is followed.
\item The kinetic energy of a body is doubled when its speed is doubled.
\item Power is the rate at which work is done, and its SI unit is the joule.
\item When a car moves at constant velocity on a level road, the net work done on the car is zero.
\end{enumerate}
\end{exercice}

\begin{exercice}[Definitions]
Define each of the following terms, giving in each case the SI unit and stating whether the quantity is a scalar or a vector:
\begin{enumerate}
\item the \cle{work} done by a constant force;
\item the \cle{kinetic energy} of a particle;
\item the \cle{power} developed by a force.
\end{enumerate}
\end{exercice}

\begin{exercice}[Direct calculations]
\begin{enumerate}
\item A crate of mass $25\ \mathrm{kg}$ is lifted vertically through $3.2\ \mathrm{m}$. Taking $g = 9.8\ \mathrm{m\,s^{-2}}$, calculate the work done against gravity.
\item A particle of mass $0.5\ \mathrm{kg}$ moves with speed $6\ \mathrm{m\,s^{-1}}$. Calculate its kinetic energy.
\item A machine does $4.5\times 10^{4}\ \mathrm{J}$ of work in $30\ \mathrm{s}$. Calculate its average power.
\end{enumerate}
\end{exercice}

\courssub{I practise}

\begin{exercice}[Work done by several forces]
A force of $50\ \mathrm{N}$ acts at $30^{\circ}$ to the horizontal and pulls a box $8\ \mathrm{m}$ along a rough horizontal floor against a friction force of $20\ \mathrm{N}$.
\begin{enumerate}
\item Calculate the work done by the pulling force.
\item Calculate the work done by the friction force.
\item Calculate the work done by the weight of the box and by the normal reaction of the floor.
\item Deduce the net work done on the box.
\end{enumerate}
\end{exercice}

\begin{exercice}[Stretching a spring]
A spring of stiffness $200\ \mathrm{N\,m^{-1}}$ is stretched from an extension of $4\ \mathrm{cm}$ to an extension of $10\ \mathrm{cm}$.
\begin{enumerate}
\item Calculate the work done by the stretching force.
\item Hence state the extra elastic potential energy stored in the spring.
\end{enumerate}
\end{exercice}

\begin{exercice}[Work--energy theorem for a car]
A car of mass $1200\ \mathrm{kg}$ accelerates from $5\ \mathrm{m\,s^{-1}}$ to $20\ \mathrm{m\,s^{-1}}$ over a distance of $100\ \mathrm{m}$ on a level road against a constant resistance of $400\ \mathrm{N}$.
\begin{enumerate}
\item Calculate the change in kinetic energy of the car.
\item Use the work--energy theorem to find the work done by the engine.
\item Deduce the average driving force exerted by the engine.
\end{enumerate}
\end{exercice}

\begin{exercice}[Conservation of energy]
A ball is thrown vertically upwards with speed $14\ \mathrm{m\,s^{-1}}$. Taking $g = 9.8\ \mathrm{m\,s^{-2}}$ and neglecting air resistance, use conservation of energy to find:
\begin{enumerate}
\item the maximum height reached above the point of projection;
\item the speed of the ball when it is $5\ \mathrm{m}$ above the point of projection.
\end{enumerate}
\end{exercice}

\courssub{I prepare for the GCE}

\begin{exercice}[Rough slope and energy method]
A block of mass $m$ kg slides from rest down a rough slope of length $10\ \mathrm{m}$ inclined at $25^{\circ}$ to the horizontal. The coefficient of friction between the block and the slope is $0.2$. Take $g = 9.8\ \mathrm{m\,s^{-2}}$.
\begin{enumerate}
\item Draw a diagram showing all the forces acting on the block. \hfill (2 marks)
\item Show that the magnitude of the friction force is $0.2\,mg\cos 25^{\circ}$. \hfill (2 marks)
\item Calculate the work done against friction as the block descends the $10\ \mathrm{m}$. \hfill (2 marks)
\item Use the work--energy theorem to find the speed of the block at the bottom of the slope. \hfill (4 marks)
\item State one assumption you have made. \hfill (1 mark)
\end{enumerate}
\end{exercice}

\begin{exercice}[Power, hills and pumping water]
A car of mass $1000\ \mathrm{kg}$ has an engine of maximum power $60\ \mathrm{kW}$ and experiences a constant resistance to motion of $500\ \mathrm{N}$. Take $g = 9.8\ \mathrm{m\,s^{-2}}$.
\begin{enumerate}
\item Find the maximum speed of the car on a level road. \hfill (3 marks)
\item The car now climbs a hill inclined at an angle $\sin^{-1}\!\left(\tfrac{1}{20}\right)$ to the horizontal, with the same resistance and the engine working at maximum power. Find its maximum speed up the hill. \hfill (4 marks)
\item A pump raises water from a river through a vertical height of $15\ \mathrm{m}$ at a rate of $40\ \mathrm{kg\,s^{-1}}$ and delivers it through a nozzle at $8\ \mathrm{m\,s^{-1}}$. Find the total power developed by the pump. \hfill (4 marks)
\item Given that the pumping system is only $75\ \%$ efficient, find the power of the motor required to drive it. \hfill (2 marks)
\end{enumerate}
\end{exercice}

\begin{exercice}[Cyclist, lorry and water jet]
\textbf{Part A.} A cyclist of total mass $80\ \mathrm{kg}$ freewheels down a hill inclined at an angle $\sin^{-1}\!\left(\tfrac{1}{25}\right)$ to the horizontal and reaches a terminal speed of $12\ \mathrm{m\,s^{-1}}$. Take $g = 9.8\ \mathrm{m\,s^{-2}}$.
\begin{enumerate}
\item By considering the forces acting at terminal speed, find the magnitude of the resistance to motion, assumed constant. \hfill (3 marks)
\item Find the power the cyclist must develop to cycle \emph{up} the same hill at a steady $6\ \mathrm{m\,s^{-1}}$ against the same resistance. \hfill (4 marks)
\end{enumerate}
\textbf{Part B.} A lorry of mass $8000\ \mathrm{kg}$ moves on level ground with its engine working at a constant power of $90\ \mathrm{kW}$ against a resistance of magnitude $250v$ newtons, where $v\ \mathrm{m\,s^{-1}}$ is its speed.
\begin{enumerate}
\item[(a)] Find the maximum speed of the lorry on level ground. \hfill (3 marks)
\item[(b)] Find the acceleration of the lorry at the instant when its speed is $10\ \mathrm{m\,s^{-1}}$. \hfill (3 marks)
\end{enumerate}
\textbf{Part C.} A fire hose of cross-sectional area $2\times 10^{-3}\ \mathrm{m^{2}}$ projects water of density $1000\ \mathrm{kg\,m^{-3}}$ horizontally at $25\ \mathrm{m\,s^{-1}}$.
\begin{enumerate}
\item[(a)] Find the mass of water discharged per second. \hfill (2 marks)
\item[(b)] Find the kinetic energy delivered to the water per second, and comment on the minimum power of the pump. \hfill (2 marks)
\end{enumerate}
\end{exercice}

\newpage

\coursec{Solutions to the exercises}

\begin{corrige}[Exercise 1 — Multiple choice]
The correct answer is \textbf{(c): $W = Fs\cos\theta$}.

\textbf{Justification.} Work is done only by the component of the force in the direction of the displacement. Resolving the force $\vec{F}$ into components parallel and perpendicular to the motion:
\begin{itemize}
\item component along the motion: $F\cos\theta$;
\item component perpendicular to the motion: $F\sin\theta$ (this does no work, since there is no displacement in that direction).
\end{itemize}
Hence
\[
W = (F\cos\theta)\times s = Fs\cos\theta.
\]
Equivalently, work is the scalar product $W = \vec{F}\cdot\vec{s} = Fs\cos\theta$, where $\theta$ is the angle between $\vec{F}$ and $\vec{s}$.

\textbf{Check of special cases:} if $\theta = 0$ (force along the motion), $W = Fs$; if $\theta = 90^{\circ}$ (force perpendicular to the motion, e.g.\ the normal reaction), $W = 0$; if $\theta = 180^{\circ}$ (force opposing the motion, e.g.\ friction), $W = -Fs$. These confirm formula (c) and reject (a), (b) and (d).

\begin{aretenir}[Work done by a constant force]
For a constant force $\vec{F}$ and a displacement $\vec{s}$:
\[
W = \vec{F}\cdot\vec{s} = Fs\cos\theta \quad (\text{joule, J})
\]
\begin{itemize}
\item $W>0$ if the force has a component in the direction of motion ($0\le\theta<90^{\circ}$);
\item $W=0$ if the force is perpendicular to the motion ($\theta=90^{\circ}$);
\item $W<0$ if the force opposes the motion ($90^{\circ}<\theta\le180^{\circ}$).
\end{itemize}
Work is a \textbf{scalar} quantity.
\end{aretenir}
\end{corrige}

\begin{corrige}[Exercise 2 — True or false]
\begin{enumerate}
\item \textbf{True.} The weight $mg$ is a \emph{conservative} force: its work depends only on the vertical displacement, not on the path. Whatever the path, the work done \emph{by} the weight is $-mgh$ when the body rises through $h$, so the work done \emph{against} gravity (by the lifting agent) is $mgh$.

\item \textbf{False.} Kinetic energy is $T = \tfrac{1}{2}mv^{2}$, i.e.\ it is proportional to the \emph{square} of the speed. Doubling the speed multiplies the kinetic energy by $2^{2} = 4$, not by $2$.

\item \textbf{False} (as stated). It is true that power is the rate at which work is done, $P = \dfrac{W}{t}$ (average) or $P = \dfrac{dW}{dt}$ (instantaneous). However, its SI unit is the \textbf{watt} ($\mathrm{W}$), where $1\ \mathrm{W} = 1\ \mathrm{J\,s^{-1}}$. The joule is the unit of work and energy, not of power.

\item \textbf{True.} At constant velocity the kinetic energy is constant, so by the work--energy theorem the \emph{net} work done on the car is zero: the positive work of the driving force exactly cancels the negative work of the resistances (friction, air resistance). Note that the engine still does work — it is the \emph{total} work that is zero.
\end{enumerate}

\begin{aretenir}[Key concepts check]
\begin{itemize}
\item Work against gravity: $W = mgh$ (path independent).
\item Kinetic energy: $T = \frac{1}{2}mv^2 \propto v^2$.
\item Power: $P = \frac{dW}{dt}$, unit = watt (W) = J/s.
\item Constant velocity $\Rightarrow$ net work = 0 (but individual forces may do work).
\end{itemize}
\end{aretenir}
\end{corrige}

\begin{corrige}[Exercise 3 — Definitions]
\begin{enumerate}
\item \textbf{Work done by a constant force.} The work done by a constant force $\vec{F}$ whose point of application undergoes a displacement $\vec{s}$ is the scalar product
\[
W = \vec{F}\cdot\vec{s} = Fs\cos\theta,
\]
where $\theta$ is the angle between the force and the displacement. SI unit: the \textbf{joule} ($\mathrm{J}$), with $1\ \mathrm{J} = 1\ \mathrm{N\,m}$. Work is a \textbf{scalar} (it can be positive, negative or zero).

\item \textbf{Kinetic energy.} The kinetic energy of a particle of mass $m$ moving with speed $v$ is the energy it possesses by virtue of its motion:
\[
T = \tfrac{1}{2}mv^{2}.
\]
SI unit: the \textbf{joule} ($\mathrm{J}$). Kinetic energy is a \textbf{scalar}, and it is always non-negative ($T \ge 0$).

\item \textbf{Power.} The power developed by a force is the rate at which the force does work:
\[
P = \frac{dW}{dt}.
\]
For a constant force moving its point of application at velocity $\vec{v}$, $P = \vec{F}\cdot\vec{v} = Fv\cos\theta$. SI unit: the \textbf{watt} ($\mathrm{W}$), with $1\ \mathrm{W} = 1\ \mathrm{J\,s^{-1}}$. Power is a \textbf{scalar}.
\end{enumerate}

\begin{propriete}[Instantaneous power for a constant force]
If a constant force $\vec{F}$ acts on a particle moving with instantaneous velocity $\vec{v}$, the instantaneous power is
\[
P = \vec{F} \cdot \vec{v} = Fv\cos\theta.
\]
\textit{Proof:} $W = \vec{F}\cdot\vec{s} \Rightarrow \frac{dW}{dt} = \vec{F}\cdot\frac{d\vec{s}}{dt} = \vec{F}\cdot\vec{v}$.
\end{propriete}
\end{corrige}

\begin{corrige}[Exercise 4 — Direct calculations]
\begin{enumerate}
\item \textbf{Work done against gravity:}
\[
W = mgh = 25 \times 9.8 \times 3.2 = 784\ \mathrm{J}.
\]

\item \textbf{Kinetic energy:}
\[
T = \tfrac{1}{2}mv^{2} = \tfrac{1}{2}\times 0.5 \times 6^{2} = \tfrac{1}{2}\times 0.5\times 36 = 9\ \mathrm{J}.
\]

\item \textbf{Average power:}
\[
P = \frac{W}{t} = \frac{4.5\times 10^{4}}{30} = 1.5\times 10^{3}\ \mathrm{W} = 1.5\ \mathrm{kW}.
\]
\end{enumerate}

\begin{methode}[Calculating work, energy and power]
\begin{enumerate}
\item Identify the known quantities: mass $m$, height $h$, speed $v$, work $W$, time $t$.
\item Choose the correct formula: $W=mgh$, $T=\frac{1}{2}mv^2$, $P=W/t$ (average) or $P=Fv$ (instantaneous, constant force).
\item Substitute values with consistent SI units (kg, m, s).
\item Compute and state the answer with the correct unit.
\end{enumerate}
\end{methode}
\end{corrige}

\begin{corrige}[Exercise 5 — Work done by several forces]
\begin{enumerate}
\item \textbf{Work done by the pulling force.} Only the horizontal component $50\cos 30^{\circ}$ does work:
\[
W_{P} = 50\cos 30^{\circ}\times 8 = 50\times\frac{\sqrt{3}}{2}\times 8 = 200\sqrt{3}\approx 346.4\ \mathrm{J}.
\]

\item \textbf{Work done by friction.} Friction opposes the motion ($\theta = 180^{\circ}$):
\[
W_{f} = -20\times 8 = -160\ \mathrm{J}.
\]

\item \textbf{Work done by the weight and the normal reaction.} Both forces are perpendicular to the (horizontal) displacement, so $\cos 90^{\circ} = 0$:
\[
W_{\text{weight}} = 0\ \mathrm{J}, \qquad W_{R} = 0\ \mathrm{J}.
\]

\item \textbf{Net work.} The net work is the sum of the works of all the forces:
\[
W_{\text{net}} = 200\sqrt{3} - 160 + 0 + 0 \approx 346.4 - 160 = 186.4\ \mathrm{J}.
\]
By the work--energy theorem, this positive net work equals the gain in kinetic energy of the box.
\end{enumerate}

\begin{aretenir}[Work done by multiple forces]
The total (net) work done on a particle is the algebraic sum of the works done by each individual force:
\[
W_{\text{net}} = \sum_i W_i = \sum_i \vec{F}_i \cdot \vec{s}.
\]
Forces perpendicular to the displacement do zero work. The work-energy theorem states $W_{\text{net}} = \Delta T$.
\end{aretenir}
\end{corrige}

\begin{corrige}[Exercise 6 — Stretching a spring]
The tension in a spring of stiffness $k$ at extension $x$ is $T = kx$ (Hooke's law). Since the tension varies with the extension, the work done by the stretching force from extension $x_{1}$ to $x_{2}$ is
\[
W = \int_{x_{1}}^{x_{2}} kx\,dx = \tfrac{1}{2}k\left(x_{2}^{2} - x_{1}^{2}\right).
\]

\begin{enumerate}
\item Here $k = 200\ \mathrm{N\,m^{-1}}$, $x_{1} = 0.04\ \mathrm{m}$, $x_{2} = 0.10\ \mathrm{m}$:
\[
W = \tfrac{1}{2}\times 200\times\left(0.10^{2} - 0.04^{2}\right) = 100\times(0.01 - 0.0016) = 100\times 0.0084 = 0.84\ \mathrm{J}.
\]

\item The work done by the stretching force is stored as elastic potential energy, so the \textbf{extra elastic potential energy} stored is
\[
\Delta E_{e} = 0.84\ \mathrm{J}.
\]
\end{enumerate}

\begin{attention}[Common mistake]
Do \emph{not} compute $\tfrac{1}{2}k(x_{2}-x_{1})^{2}$: the spring already stores energy $\tfrac{1}{2}kx_{1}^{2}$ at extension $x_{1}$. The \emph{extra} energy is the difference $\tfrac{1}{2}kx_{2}^{2}-\tfrac{1}{2}kx_{1}^{2}$.
\end{attention}

\begin{propriete}[Elastic potential energy]
For a spring of stiffness $k$ extended by $x$ from its natural length, the elastic potential energy stored is
\[
E_e = \frac{1}{2}kx^2.
\]
The work done in stretching from $x_1$ to $x_2$ equals the increase in elastic potential energy: $W = \Delta E_e = \frac{1}{2}k(x_2^2 - x_1^2)$.
\end{propriete}
\end{corrige}

\begin{corrige}[Exercise 7 — Work--energy theorem for a car]
\begin{enumerate}
\item \textbf{Change in kinetic energy:}
\[
\Delta T = \tfrac{1}{2}m\left(v^{2}-u^{2}\right) = \tfrac{1}{2}\times 1200\times\left(20^{2}-5^{2}\right) = 600\times(400-25) = 600\times 375 = 225\,000\ \mathrm{J} = 2.25\times 10^{5}\ \mathrm{J}.
\]

\item \textbf{Work done by the engine.} The work--energy theorem states:
\[
W_{\text{engine}} + W_{\text{resistance}} = \Delta T.
\]
The resistance does work $W_{\text{resistance}} = -400\times 100 = -40\,000\ \mathrm{J}$. Hence
\[
W_{\text{engine}} = \Delta T - W_{\text{resistance}} = 225\,000 + 40\,000 = 265\,000\ \mathrm{J} = 2.65\times 10^{5}\ \mathrm{J}.
\]

\item \textbf{Average driving force.} Since $W_{\text{engine}} = F_{\text{av}}\times s$:
\[
F_{\text{av}} = \frac{265\,000}{100} = 2650\ \mathrm{N}.
\]
\end{enumerate}

\begin{methode}[Applying the work-energy theorem]
\begin{enumerate}
\item Calculate the change in kinetic energy: $\Delta T = \frac{1}{2}m(v^2 - u^2)$.
\item Calculate the work done by each force (positive, negative, or zero).
\item Apply $W_{\text{net}} = \Delta T$ or $\sum W_i = \Delta T$.
\item Solve for the unknown (work, force, distance, or final speed).
\end{enumerate}
\end{methode}
\end{corrige}

\begin{corrige}[Exercise 8 — Conservation of energy]
Take the point of projection as the level of zero potential energy. With no air resistance, the total mechanical energy $E = \tfrac{1}{2}mv^{2} + mgh$ is conserved.

\begin{enumerate}
\item \textbf{Maximum height.} At the highest point the speed is zero. Conservation of energy between projection and the top:
\[
\tfrac{1}{2}m(14)^{2} + 0 = 0 + mgH
\quad\Longrightarrow\quad
H = \frac{14^{2}}{2g} = \frac{196}{2\times 9.8} = \frac{196}{19.6} = 10\ \mathrm{m}.
\]

\item \textbf{Speed at height $5\ \mathrm{m}$.} Conservation of energy between projection and the point at height $h = 5\ \mathrm{m}$:
\[
\tfrac{1}{2}m(14)^{2} = \tfrac{1}{2}mv^{2} + mg(5).
\]
Dividing by $m$ and multiplying by $2$:
\[
196 = v^{2} + 2\times 9.8\times 5 = v^{2} + 98
\quad\Longrightarrow\quad
v^{2} = 98
\quad\Longrightarrow\quad
v = \sqrt{98} = 7\sqrt{2}\approx 9.9\ \mathrm{m\,s^{-1}}.
\]
(The same speed is found on the way up and on the way down, since energy is conserved.)
\end{enumerate}

\begin{propriete}[Conservation of mechanical energy]
When only conservative forces (gravity, elastic forces) do work, the total mechanical energy is constant:
\[
E = T + V = \frac{1}{2}mv^2 + mgh + \frac{1}{2}kx^2 = \text{constant}.
\]
If non-conservative forces (friction, air resistance) act, then $W_{\text{nc}} = \Delta E = \Delta T + \Delta V$.
\end{propriete}
\end{corrige}

\begin{corrige}[Exercise 9 — Rough slope and energy method]
\begin{enumerate}
\item \textbf{Diagram (2 marks).} The forces are: the \textbf{weight} $mg$ acting vertically downwards; the \textbf{normal reaction} $R$ perpendicular to the slope; the \textbf{friction} $F$ acting up the slope, opposing the motion.
\begin{popfigure}
\begin{center}
\begin{tikzpicture}[scale=1.0]
\draw[thick] (0,0) -- (6,0);
\draw[thick] (0,0) -- (6,2.5);
\draw (1.2,0) arc (0:22.6:1.2) node[midway, right] {$25^{\circ}$};
\draw[fill=popblueL, draw=popdark] (2.6,1.083) rectangle (3.4,1.883);
\draw[->, very thick, poporange] (3.0,1.483) -- (3.0,-0.6) node[below] {$mg$};
\draw[->, very thick, popgreen] (3.0,1.883) -- (3.35,2.75) node[above right] {$R$};
\draw[->, very thick, poppurple] (2.6,1.35) -- (1.6,0.93) node[below left] {$F$};
\draw[->, thick, popdark, dashed] (3.4,1.35) -- (4.6,1.85) node[right] {motion};
\end{tikzpicture}
\end{center}
\legende{Forces on a block sliding down a rough slope.}
\end{popfigure}

\textit{Mark scheme: 1 mark for weight vertically down and $R$ perpendicular to the plane; 1 mark for friction directed up the slope.}

\item \textbf{Friction force (2 marks).} Resolving the weight perpendicular to the slope, there is no acceleration in that direction, so
\[
R = mg\cos 25^{\circ}.
\]
Since the block is sliding, friction is limiting: $F = \mu R$ with $\mu = 0.2$. Hence
\[
F = 0.2\,mg\cos 25^{\circ}. \qquad\blacksquare
\]
\textit{Mark scheme: 1 mark for $R = mg\cos 25^{\circ}$; 1 mark for $F = \mu R$.}

\item \textbf{Work done against friction (2 marks).} The friction acts opposite to the displacement over $10\ \mathrm{m}$:
\[
W_{f} = F\times 10 = 10\times 0.2\,mg\cos 25^{\circ} = 2\,mg\cos 25^{\circ}.
\]
Numerically, with $m$ in kg:
\[
W_{f} = 2\times 9.8\times\cos 25^{\circ}\times m \approx 17.76\,m\ \mathrm{J}.
\]
\textit{Mark scheme: 1 mark for method $W = Fs$; 1 mark for the correct expression/value.}

\item \textbf{Speed at the bottom (4 marks).} The vertical drop is
\[
h = 10\sin 25^{\circ}\ \mathrm{m}.
\]
The work--energy theorem between the top (rest) and the bottom gives
\[
\underbrace{mgh}_{\text{work of weight}} \;-\; \underbrace{0.2\,mg\cos 25^{\circ}\times 10}_{\text{work against friction}} = \tfrac{1}{2}mv^{2}.
\]
Dividing by $m$:
\[
\tfrac{1}{2}v^{2} = 10g\sin 25^{\circ} - 2g\cos 25^{\circ} = 9.8\left(10\sin 25^{\circ} - 2\cos 25^{\circ}\right).
\]
Numerically, $10\sin 25^{\circ} \approx 4.226$ and $2\cos 25^{\circ} \approx 1.813$:
\[
\tfrac{1}{2}v^{2} = 9.8\times(4.226 - 1.813) = 9.8\times 2.413 \approx 23.65
\quad\Longrightarrow\quad
v^{2} \approx 47.3,
\]
\[
\boxed{v \approx 6.9\ \mathrm{m\,s^{-1}}.}
\]
\textit{Mark scheme: 1 mark for the height $h = 10\sin 25^{\circ}$; 1 mark for a correct work--energy equation; 1 mark for correct algebra (mass cancels); 1 mark for the numerical answer.}

\item \textbf{Assumption (1 mark).} Any one of: friction is constant along the whole slope (coefficient uniform); the block is treated as a particle; air resistance is neglected; $g$ is constant.
\end{enumerate}

\begin{aretenir}[Energy method on an inclined plane]
For a particle sliding down a rough slope of length $s$ and angle $\theta$:
\begin{itemize}
\item Vertical drop: $h = s\sin\theta$.
\item Normal reaction: $R = mg\cos\theta$.
\item Friction force: $F = \mu mg\cos\theta$.
\item Work against friction: $W_f = -\mu mg s\cos\theta$.
\item Work-energy: $mgh - \mu mg s\cos\theta = \frac{1}{2}mv^2$.
\end{itemize}
Mass $m$ always cancels out.
\end{aretenir}
\end{corrige}

\begin{corrige}[Exercise 10 — Power, hills and pumping water]
\begin{enumerate}
\item \textbf{Maximum speed on the level (3 marks).} At maximum (constant) speed the acceleration is zero, so the driving force $F$ balances the resistance:
\[
F = 500\ \mathrm{N}.
\]
Using $P = Fv$ with $P = 60\,000\ \mathrm{W}$:
\[
v_{\max} = \frac{P}{F} = \frac{60\,000}{500} = 120\ \mathrm{m\,s^{-1}}.
\]
\textit{Mark scheme: 1 mark for $F = R$ at top speed; 1 mark for $P = Fv$; 1 mark for the answer.}

\item \textbf{Maximum speed up the hill (4 marks).} On the hill, the driving force must overcome the resistance \emph{and} the component of the weight down the slope. With $\sin\theta = \tfrac{1}{20}$:
\[
mg\sin\theta = 1000\times 9.8\times \frac{1}{20} = 490\ \mathrm{N}.
\]
At maximum speed (no acceleration):
\[
F = 500 + 490 = 990\ \mathrm{N}.
\]
Then
\[
v_{\max} = \frac{P}{F} = \frac{60\,000}{990} \approx 60.6\ \mathrm{m\,s^{-1}}.
\]
\textit{Mark scheme: 1 mark for resolving the weight; 1 mark for the total resisting force; 1 mark for $v = P/F$; 1 mark for the answer.}

\item \textbf{Power of the pump (4 marks).} Each second, a mass $\dot{m} = 40\ \mathrm{kg}$ is raised through $15\ \mathrm{m}$ \emph{and} given kinetic energy at $8\ \mathrm{m\,s^{-1}}$. The pump must supply both:
\[
P = \underbrace{\dot{m}gh}_{\text{gain of PE per second}} + \underbrace{\tfrac{1}{2}\dot{m}v^{2}}_{\text{gain of KE per second}}.
\]
\[
P = 40\times 9.8\times 15 + \tfrac{1}{2}\times 40\times 8^{2} = 5880 + 1280 = 7160\ \mathrm{W} \approx 7.16\ \mathrm{kW}.
\]
\textit{Mark scheme: 1 mark for the PE term; 1 mark for the KE term; 1 mark for adding both; 1 mark for the answer. A very common error is to forget the kinetic energy of the delivered water.}

\item \textbf{Power of the motor (2 marks).} With efficiency $75\%$:
\[
\text{efficiency} = \frac{\text{useful power output}}{\text{power input}}
\quad\Longrightarrow\quad
P_{\text{motor}} = \frac{7160}{0.75} \approx 9547\ \mathrm{W} \approx 9.55\ \mathrm{kW}.
\]
\textit{Mark scheme: 1 mark for the correct use of efficiency; 1 mark for the answer.}
\end{enumerate}

\begin{propriete}[Power for pumping fluids]
For a pump delivering a mass flow rate $\dot{m}$ (kg/s) of fluid:
\begin{itemize}
\item Power to raise through height $h$: $P_{\text{PE}} = \dot{m}gh$.
\item Power to give speed $v$: $P_{\text{KE}} = \frac{1}{2}\dot{m}v^2$.
\item Total useful power: $P = \dot{m}gh + \frac{1}{2}\dot{m}v^2$.
\item If pump efficiency is $\eta$, input power = $P/\eta$.
\end{itemize}
\end{propriete}
\end{corrige}

\begin{corrige}[Exercise 11 — Cyclist, lorry and water jet]
\textbf{Part A — the cyclist.}
\begin{enumerate}
\item \textbf{Resistance at terminal speed (3 marks).} At terminal (constant) speed, the forces along the slope balance: the component of the weight down the slope equals the resistance $R$. With $\sin\theta = \tfrac{1}{25}$:
\[
R = mg\sin\theta = 80\times 9.8\times \frac{1}{25} = \frac{784}{25} = 31.36\ \mathrm{N}.
\]
\textit{Mark scheme: 1 mark for resolving the weight; 1 mark for equilibrium at terminal speed; 1 mark for the value.}

\item \textbf{Power to climb at $6\ \mathrm{m\,s^{-1}}$ (4 marks).} Going \emph{up} the hill at steady speed, the cyclist must overcome both the resistance and the component of the weight:
\[
F = R + mg\sin\theta = 31.36 + 31.36 = 62.72\ \mathrm{N}.
\]
The power required is
\[
P = Fv = 62.72\times 6 = 376.32\ \mathrm{W} \approx 376\ \mathrm{W}.
\]
\textit{Mark scheme: 1 mark for recognising that both resistance and gravity oppose the motion uphill; 1 mark for $F = 62.72\ \mathrm{N}$; 1 mark for $P = Fv$; 1 mark for the answer.}
\end{enumerate}

\textbf{Part B — the lorry.}
\begin{enumerate}
\item[(a)] \textbf{Maximum speed (3 marks).} At maximum speed the driving force equals the resistance $250v$, and $P = Fv$:
\[
90\,000 = 250v \times v = 250v^{2}
\quad\Longrightarrow\quad
v^{2} = \frac{90\,000}{250} = 360
\quad\Longrightarrow\quad
v_{\max} = \sqrt{360} = 6\sqrt{10}\approx 19.0\ \mathrm{m\,s^{-1}}.
\]
\textit{Mark scheme: 1 mark for $F = 250v$ at top speed; 1 mark for $P = Fv$; 1 mark for the answer.}

\item[(b)] \textbf{Acceleration at $v = 10\ \mathrm{m\,s^{-1}}$ (3 marks).} The driving force at that instant is
\[
F = \frac{P}{v} = \frac{90\,000}{10} = 9000\ \mathrm{N}.
\]
The resistance is $250\times 10 = 2500\ \mathrm{N}$. Newton's second law:
\[
F - 250v = ma
\quad\Longrightarrow\quad
9000 - 2500 = 8000\,a
\quad\Longrightarrow\quad
a = \frac{6500}{8000} = 0.8125\ \mathrm{m\,s^{-2}} \approx 0.81\ \mathrm{m\,s^{-2}}.
\]
\textit{Mark scheme: 1 mark for $F = P/v$; 1 mark for applying Newton's second law with the resistance; 1 mark for the answer.}
\end{enumerate}

\textbf{Part C — the water jet.}
\begin{enumerate}
\item[(a)] \textbf{Mass discharged per second (2 marks).} In one second the water occupies a cylinder of length $v = 25\ \mathrm{m}$ and cross-section $A = 2\times 10^{-3}\ \mathrm{m^{2}}$, i.e.\ a volume $Av$:
\[
\dot{m} = \rho A v = 1000\times 2\times 10^{-3}\times 25 = 50\ \mathrm{kg\,s^{-1}}.
\]
\textit{Mark scheme: 1 mark for $\dot{m} = \rho Av$; 1 mark for the value.}

\item[(b)] \textbf{Kinetic energy delivered per second (2 marks).}
\[
\frac{\Delta T}{\Delta t} = \tfrac{1}{2}\dot{m}v^{2} = \tfrac{1}{2}\times 50\times 25^{2} = 25\times 625 = 15\,625\ \mathrm{W} \approx 15.6\ \mathrm{kW}.
\]
\textbf{Comment:} since the water leaves horizontally, no potential energy is involved; the kinetic energy given to the water each second is exactly the useful work done by the pump each second. Hence the pump must develop \emph{at least} $15.6\ \mathrm{kW}$ — in practice more, because of friction losses in the hose and the pump's own efficiency being less than $100\%$.
\textit{Mark scheme: 1 mark for $\tfrac{1}{2}\dot{m}v^{2}$; 1 mark for the value and the comment on minimum power.}
\end{enumerate}

\begin{aretenir}[Power and variable resistance]
When resistance depends on speed ($R = kv$ or $R = kv^2$):
\begin{itemize}
\item At maximum speed: $P = F_{\text{drive}} v = R v \Rightarrow P = kv^2$ or $kv^3$.
\item At any speed: $F_{\text{drive}} = P/v$, then $ma = P/v - R(v)$.
\end{itemize}
For water jets: mass flow rate $\dot{m} = \rho A v$, power = $\frac{1}{2}\dot{m}v^2 + \dot{m}gh$ (if height change).
\end{aretenir}
\end{corrige}

\newpage

\coursec{Summary sheet}

\begin{popfigure}
\begin{tikzpicture}[grow cyclic, mindmap, concept color=popblue,
  level 1 concept/.append style={level distance=3.6cm, sibling angle=60},
  every node/.append style={font=\small, text width=2.2cm, align=center}]
\node[concept, text=white] {Work, Energy \& Power}
  child[concept color=poporange] { node[concept, text=white] {Work of a force} 
    child[concept color=poporange!75] { node[concept, text=white] {$W = F s \cos\theta$} }
    child[concept color=poporange!75] { node[concept, text=white] {Against gravity: $W = mgh$} }
  }
  child[concept color=popgreen] { node[concept, text=white] {Variable force}
    child[concept color=popgreen!75] { node[concept, text=white] {$W = \int F\,dx$} }
    child[concept color=popgreen!75] { node[concept, text=white] {Spring: $W = \tfrac12 kx^2$} }
  }
  child[concept color=poppurple] { node[concept, text=white] {Energy}
    child[concept color=poppurple!75] { node[concept, text=white] {KE $= \tfrac12 mv^2$} }
    child[concept color=poppurple!75] { node[concept, text=white] {GPE $= mgh$} }
  }
  child[concept color=popgold] { node[concept, text=white] {Work--Energy}
    child[concept color=popgold!75] { node[concept, text=white] {$W_{\text{net}} = \Delta$KE} }
    child[concept color=popgold!75] { node[concept, text=white] {Conservation: KE+GPE const.} }
  }
  child[concept color=poppink] { node[concept, text=white] {Power}
    child[concept color=poppink!75] { node[concept, text=white] {$P = \frac{W}{t} = Fv$} }
    child[concept color=poppink!75] { node[concept, text=white] {1 W = 1 J/s} }
  }
  child[concept color=popblue!70!black] { node[concept, text=white] {Applications}
    child[concept color=popblue!60] { node[concept, text=white] {Pumps \& water jets} }
    child[concept color=popblue!60] { node[concept, text=white] {Vehicles on inclines} }
  };
\end{tikzpicture}
\legende{Mind map of the chapter Work, Energy and Power.}
\end{popfigure}

\begin{aretenir}[Chapter Summary — Work, Energy and Power]
\textbf{Key definitions}
\begin{itemize}
\item The \cle{work} done by a constant force $\vec{F}$ over a displacement $s$ is $W = F s \cos\theta$, where $\theta$ is the angle between the force and the displacement. Unit: the \cle{joule} (J), $1\ \mathrm{J} = 1\ \mathrm{N\,m}$.
\item Work done \cle{against gravity} in raising a mass $m$ through a vertical height $h$: $W = mgh$ (independent of the path taken).
\item Work of a \cle{variable force} $F(x)$: $W = \displaystyle\int_{x_1}^{x_2} F(x)\,dx$ (area under the force--displacement graph). For a spring obeying Hooke's law $F = kx$, the work done \emph{in stretching or compressing} the spring (work against the spring force) is $W = \tfrac12 kx^2$.
\item \cle{Kinetic energy}: $\mathrm{KE} = \tfrac12 mv^2$. \cle{Gravitational potential energy}: $\mathrm{GPE} = mgh$ (relative to a chosen reference level).
\item \cle{Power} is the rate of doing work: average power $P = \dfrac{W}{t}$, instantaneous power $P = \dfrac{dW}{dt} = Fv$ (when force and velocity are parallel). Unit: the \cle{watt} (W); $1\ \mathrm{kW} = 1000\ \mathrm{W}$; $1\ \mathrm{hp} \approx 746\ \mathrm{W}$.
\end{itemize}

\textbf{Laws and theorems to know}
\begin{itemize}
\item \cle{Work--Energy Theorem}: the net work done on a body equals its change in kinetic energy:
\[ W_{\text{net}} = \Delta \mathrm{KE} = \tfrac12 mv^2 - \tfrac12 mu^2. \]
\item \cle{Principle of Conservation of Mechanical Energy}: when only conservative forces (gravity, spring force) act,
\[ \mathrm{KE}_1 + \mathrm{GPE}_1 = \mathrm{KE}_2 + \mathrm{GPE}_2. \]
\item With a constant resistive force $R$ acting over a distance $s$: 
\[ \mathrm{KE}_1 + \mathrm{GPE}_1 = \mathrm{KE}_2 + \mathrm{GPE}_2 + Rs \]
(energy lost as work done against resistance).
\item \cle{Power--velocity relation}: a vehicle with driving force $F$ moving at speed $v$ develops power $P = Fv$; at maximum constant speed, driving force equals total resistance.
\item On an incline of angle $\alpha$: component of weight down the plane $= mg\sin\alpha$; power required to move uphill at constant speed $v$ against resistance $R$ is $P = (R + mg\sin\alpha)\,v$.
\item \cle{Pumping water}: to raise a mass $m$ of water through height $h$ in time $t$, $P = \dfrac{mgh}{t} = \dfrac{\rho V g h}{t}$. For a jet ejecting mass at rate $\dot m$ with speed $v$: kinetic power $P = \tfrac12 \dot m v^2$; using $\dot m = \rho A v$, $P = \tfrac12 \rho A v^3$.
\end{itemize}

\textbf{Methods}
\begin{enumerate}
\item Draw a clear force diagram; resolve the weight into components parallel and perpendicular to the direction of motion.
\item Choose the appropriate tool: constant acceleration $\Rightarrow$ energy methods are often quicker than $F=ma$; variable force $\Rightarrow$ integrate $F(x)$.
\item Fix a reference level for GPE and keep it consistent throughout the question.
\item For power questions, write $P = Fv$ and identify $F$ from equilibrium (constant speed) or from Newton's second law (accelerating).
\item For pumps and jets, first find the mass moved per second ($\dot m$), then apply energy per second.
\end{enumerate}

\textbf{Common mistakes}
\begin{itemize}
\item Using $W = Fs$ when the force is not parallel to the displacement — the factor $\cos\theta$ is essential.
\item Forgetting that work against gravity depends only on the \emph{vertical} height gained, not the path length.
\item Mixing units: convert km/h to m/s (divide by $3.6$), kW to W, minutes to seconds.
\item Sign errors in the work--energy theorem: work done \emph{by} a resistive force is negative.
\item Confusing mass $m$ with mass \emph{per second} $\dot m$ in water-jet problems.
\item Assuming $P = Fv$ implies constant $F$: at constant power, $F$ decreases as $v$ increases.
\end{itemize}

\textbf{GCE tips}
\begin{itemize}
\item Quote the theorem before substituting: examiners award method marks for ``by conservation of energy'' or ``by the work--energy theorem''.
\item Keep $g = 9.8\ \mathrm{m\,s^{-2}}$ (or $9.81$ as required) and give final answers to 3 significant figures unless stated otherwise.
\item In ``maximum speed'' vehicle problems, state explicitly: acceleration zero $\Rightarrow$ driving force equals total resistance.
\item Show the energy balance line by line (initial KE + initial GPE = final KE + final GPE + work against resistance); this structure earns full credit even if arithmetic slips.
\end{itemize}
\end{aretenir}

\findecours

\end{document}
